\chapter{物理探査}
\section{観測量から地下の物性を推定する}
\label{theory:exploration:observables}
地下の岩石をすべて掘削して見ることはできない。物理探査では地表・空中・孔内などで測定した物理量を、地下の物性と幾何構造へ結び付ける。まず「どの信号を測るか」を定め、物性から観測値を計算する順問題を作り、その逆を推定する。観測量が異なれば感度のある対象も異なる。

\begin{longtable}{p{0.23\textwidth}p{0.29\textwidth}p{0.34\textwidth}}
\hline
手法 & 主な観測量 & 主な物性・構造\\\hline
重力 & 重力加速度の差 & 密度差と形状・深度\\
磁気 & 磁場・磁気異常 & 誘導磁化・残留磁化\\
地震・音波 & 弾性波・音波の走時や波形 & 波速、インピーダンス、境界\\
電気・電磁 & 電位差、電場・磁場応答 & 比抵抗・導電率\\
地中レーダ & 高周波電磁波の反射 & 誘電率・導電率の境界\\
放射能 & 主に天然ガンマ線の強度 & 表層付近の核種分布や移動\\
地温 & 地中・地表面の温度 & 熱の伝導・移流、熱水等\\
\hline
\end{longtable}

地震・音波探査では発信源と受振器を混同しない。地上で用いるジオフォンは一般に地盤の速度に応答する。一方、水中曳航のストリーマでは、管内に並べたハイドロフォンが音圧変動を測る。海水の水中音波受信を、地上の速度計を防水にしただけの装置と説明するのは不適切である\cite{usgs_marine_seismic}。

AE（アコースティック・エミッション）は、岩石などの微小破壊・変形に伴って放出される過渡的な弾性波である。岩石破壊実験や発生位置の推定に用いられる\cite{sato_lei2009_ae}。対象が岩石であることは重力・磁気探査との共通点だが、直接の観測量は異なる。令和6年度問34ではこの区別が問われる。

\section{MT・CSAMTとレーダを区別する}
\label{theory:exploration:electromagnetic}
MT（地磁気地電流法）は自然に変動する電場と磁場を測定し、地下の比抵抗構造を推定する\cite{usgs_mt_metadata}。電場と磁場の比を複素インピーダンスと呼ぶ。均質・等方な導電半空間で、変位電流を無視した平面波近似を考えると、導電率\(\sigma=1/\rho\)について
\[
 |Z|^2=\frac{\omega\mu}{\sigma},\qquad
 \rho_a=\frac{|E/H|^2}{\omega\mu},\qquad
 \delta=\sqrt{\frac{2\rho}{\omega\mu}}
\]
である。\(\delta\)は振幅が減衰する代表距離で、周波数が低いほど、また比抵抗が高いほど大きい。実際の層構造では\(\rho_a\)は各層の抵抗率そのものではなく、周波数ごとの応答を表す見掛比抵抗である。観測された\(B\)から\(H\)へ変換する場合には\(B=\mu H\)の単位と係数をそろえる。

CSAMTは人工送信源を用い、十分遠方の条件でMTに類似する電磁応答を解析する。送信源の近傍や遷移域では単純な平面波式が成り立たず、その効果の扱いが必要になる。CSAMTによる実地の比抵抗断面の例として西条平野の調査を参照する\cite{ochi_csamt_saijo}。磁場も測るからといって磁化率を主要な推定対象と考えるのが、平成27年度問19の誤りである。

地中レーダはこれとは異なり、電磁波を放射して境界からの反射を測る。通常の極浅部調査で用いる比較的高い周波数帯と、誘電率・導電率の違いを押さえる。高周波で常に深く見えるわけではなく、導電性の高い土壌や水による減衰が到達深度を制限する。

\section{測定と解釈の限界を分ける}
同じ重力異常を、浅い小岩体と深い大岩体が近似的に生むことがある。磁気には磁化方向、電磁探査には含水・塩分・粘土などの複数要因がある。したがって「異常を検出した」ことから特定岩種や製品を直ちに断定しない。物性、幾何形状、測定感度、ノイズの前提を列挙し、別の探査や露頭・ボーリングで制約を加える。

重力補正の具体的な導出は重力・測地章、誘導磁化と残留磁化の区別は地磁気章に置く。以下ではそれらを参照し、個別問題では指定条件と誤った対応だけを特定する。地震探査の走時・信号処理は次節で詳しく扱う。孔内計測の詳細は残る原典バッチとともに追加する。
\section{探査データ：測る、整える、地下へ戻す}
\label{theory:exploration:data}
受振器に入る波形は地下構造の写しそのものではない。震源波形、伝播経路、反射・屈折、計器の応答、必要でない波や環境雑音が重なった記録である。計測配置を設計し、時間と振幅をそろえ、対象成分を取り出して地下モデルへ結び付ける。フィルタ、重合、走時補正、画像化、逆解析はそれぞれ異なる役割を持つ。

\subsection{線形性と時不変性：応答を重ねる条件}
\label{theory:exploration:lti}
線形性は入力の和や定数倍に対する応答が同じ操作を保つことである。時不変性は入力を時間移動したとき、応答も同じだけ移動することである。両方が成立するLTI系では単位インパルス応答 \(h\) を使い、
\[
 y(t)=\int_{-\infty}^{\infty}x(\tau)h(t-\tau)\,d\tau,
 \qquad y[n]=\sum_k x[k]h[n-k]
\]
と表せる\cite{mit_signals2010}。任意入力を小さなインパルスの和へ分解し、その応答を加えるという意味である。

時変の線形系でも重ね合わせはできるが、一般には入力時刻にも依存する \(h(t,\tau)\) が必要である。「線形」だけから単一応答の時間移動による畳み込みまで保証しない。計器の応答を補償するときも、応答の零点付近では雑音を強く増幅する場合がある。

\subsection{フィルタ：残す周波数と抑える周波数}
\label{theory:exploration:filters}
低域通過は低周波数を残し高周波数を抑える。\(T=1/f\) なので、長周期を残す処理である。帯域通過は範囲内を残す。ノッチは狭い帯域を阻止する処理であり、商用電源等の周期的混入に用いる\cite{scipy_notch}。ノッチを「特定周波数だけを通す」と覚えると保持と除去が逆になる。

理想フィルタの通過利得1・阻止利得0はモデルの定義である。有限長の実装には遷移帯域、リップル、位相変化、端の処理がある。抑えたい雑音と対象信号の帯域が重なれば、雑音だけを無条件に取り除くことはできない。二次元の画像には空間方向の平滑化や微分を行い、エッジや構造を調べる。フィルタは二次元だけに限定された概念ではない。

\subsection{移動平均：分散の低減と周期消去の反例}
\label{theory:exploration:moving}
\(M\) 点の等重み平均 \(y[n]=M^{-1}\sum_{k=0}^{M-1}x[n-k]\) は平滑化である。各点に独立・平均ゼロ・分散 \(\sigma^2\) の雑音が含まれると、平均後の雑音分散は \(\sigma^2/M\)、標準偏差は \(\sigma/\sqrt{M}\) となる。有限点でランダム雑音が全て消えるわけではない。中央窓は未来側の値も使うため、即時の処理と後処理を区別する。

周波数応答は等比級数を足して独自に導くと、
\[
 H(e^{i\omega})=\frac1M\sum_{k=0}^{M-1}e^{-i\omega k}
 =e^{-i\omega(M-1)/2}\frac{\sin(M\omega/2)}{M\sin(\omega/2)}.
\]
\(\omega=2\pi j/M\) の非DC成分では零点を持つ。このため、周期がちょうど \(M\) 標本の正弦波は消せる。\(M=4\) なら一周期の四つの値の和がゼロとなる。一方、周期がずれれば同じ窓で完全には消えない。MITの移動平均の応答図も周波数によって利得と零点が変わることを示す\cite{mit_moving2012}。平成28年度問20の公式2を、移動平均では周期雑音を決して除去できないという数学的命題へ変更しない。

\subsection{スタックの平方根則を信号と分散から導く}
\label{theory:exploration:stack}
同じ対象信号 \(s\) と雑音 \(n_j\) の記録を \(N\) 本重ねる。
\[
 \sum_{j=1}^{N}(s+n_j)=Ns+\sum_jn_j,
 \qquad \operatorname{Var}\!\left(\sum_jn_j\right)=N\sigma^2
\]
は無相関・同分散雑音の条件で成り立つ。信号振幅は \(N\) 倍、雑音標準偏差は \(\sqrt N\) 倍なので、振幅S/Nの改善倍率は \(\sqrt N\) である。平均として \(N\) で割っても比は同じである。電力S/Nなら改善倍率は \(N\) となる。\(N^2\) はどちらの改善倍率でもない。

一般の雑音では分散に \(2\sum_{i<j}\operatorname{Cov}(n_i,n_j)\) が加わる。同じ電源雑音が各記録に残れば、回数だけ増やして平方根則を期待できない。対象波形も時刻がずれれば相殺されるため、走時差や極性をそろえる。USGSの実データ処理ではNMOとCMP重合、重合後の処理が別々に記録される\cite{usgs_processing2000}。

受振器群は独立雑音の平均化に加え、入射方向と空間波長によるフィルタとして働く。表面波は空間的に相関した波であり、単純な独立雑音とは異なる。群間隔と配列長によって対象反射も弱まる場合がある。同じ起振点で反復する垂直スタックと、共通中点の異なる経路をそろえるCMP重合を区別する。

\subsection{標本化：二倍という数と境界条件}
\label{theory:exploration:sampling}
標本化は \(x[n]=x(nT_s)\)、標本化周波数は \(f_s=1/T_s\) である。帯域制限信号を理想的に再構成するには、周期的に複写されるスペクトルが重ならないことが必要である。最高周波数成分を含む信号には \(f_s>2f_{\max}\) と余裕を取る\cite{mit_signals2010}。

等号では \(x(t)=\sin(2\pi f_{\max}t)\) を \(t=n/(2f_{\max})\) で測ると全てゼロになる。最高周波数を含まない帯域や境界成分の扱いを定義しないまま、「二倍以上だから全て完全復元」と断定しない。現実には量子化、有限記録長、時刻誤差、帯域外雑音もある。aliasingは異なる連続信号が同じ標本を与える現象であり、取得後にFFTをかけるだけでは区別できない。

\subsection{FFT・相互相関・コヒーレンス}
\label{theory:exploration:fft}
有限標本のDFTは \(X[k]=\sum_{n=0}^{N-1}x[n]e^{-2\pi i kn/N}\) である。FFTは同じ和を構造を利用して速く計算する算法である。radix-2では二のべき乗の標本数が扱いやすいが、それ以外が計算不能ではない\cite{scipy_fft}。ゼロを追加して計算点を増やしても、実測時間が延びたり近い二周波数を区別する情報が増えたりするわけではない。

\label{theory:exploration:correlation}
相互相関 \(R_{xy}[\ell]=\sum_n x[n]y[n+\ell]\) は一方をずらして類似を測る。採用する符号 convention と規格化によって遅れの符号や振幅を確認する。周期信号なら複数のピークがあり、最大値だけで唯一の到達遅れを保証しない。

代表的な二乗コヒーレンスは \(C_{xy}(f)=|S_{xy}(f)|^2/(S_{xx}(f)S_{yy}(f))\) であり、周波数ごとの関係を測る\cite{scipy_coherence}。相互相関と関連するが、時間遅れを横軸にする同じ量ではない。非零の自己スペクトルと適切な統計推定の条件を考え、二信号の因果関係をこの値だけで断定しない。

\subsection{走時・静補正・反射・CMPを結ぶ}
\label{theory:exploration:seismicterms}
単純な水平二層の屈折波では走時直線の切片を原点走時と呼ぶ。観測距離ゼロの実測初動と同一ではなく、直線の外挿である。中間層が薄く、初動として一度も現れない場合がある。また、低速度層は通常の臨界屈折条件を満たさず検出しにくい。USGSはこれらを別々に説明する\cite{usgs_refraction2013}。ブラインド層の日本語の使用範囲が広い場合もあるため、層が見えない仕組みを明示する。

単純な垂直入射の反射では \(Z=\rho V\)、\(R=(Z_2-Z_1)/(Z_2+Z_1)\) が反射振幅を支配する。\(Z\) はインピーダンス、\(R\) は反射係数であり別の量である\cite{epa_seismic_reflection}。静補正は地形や表層の低速度層等が生む走時ずれをそろえる\cite{usgs_static2004}。ブーゲー補正は質量による重力を扱い、走時補正ではない。

CMP重合では共通中点の記録を集め、通常NMO等の時間差を補正する。水平・均質の近似から外れると、共通中点が同じ反射点を意味するとは限らない。速度モデル、走時と深度の対応、データの範囲を確かめる。

\subsection{下方接続とmigrationは地震波にも用いる}
\label{theory:exploration:continuation}
migrationは傾斜面の反射や回折記録を地下の位置へ再配置する画像化である。Stanfordの説明では地表の波動場を深度方向へ下方接続し、適切な時刻の波動場から画像を作る\cite{stanford_downward2001}。重力・磁気の下方接続と名前を共有するが、伝播させる場と方程式が異なる。下方接続をポテンシャル場専用の処理と覚えるのは誤りである。

令和元年度問34の初回5はこの限定から生じた。初回記録は保存し、波動下方接続の一次説明と、送受信コイルを用いるrigid-boom EMIの電磁誘導の原理を別に確認して4へ訂正した\cite{rigid_boom2019}。後者は一次論文の検索収録を確認した範囲であり、PDF全文の取得・読了には至っていない。装置の日本語名の最終照合は査読待ちとする。

\subsection{順問題・逆問題・正則化：何を更新するのか}
\label{theory:exploration:inverse}
順問題はモデル \(m\) を与えてデータ \(F(m)\) を計算する。逆問題はデータ \(d\) からモデルを推定する。計測はデータの取得であり、計測そのものを順問題の数値計算と同一視しない。代表的な目的関数は
\[
 \Phi(m)=\|W_d(F(m)-d)\|_2^2+\lambda\|L(m-m_0)\|_2^2
\]
である。最初の項は誤差で重み付けした残差、第二項は滑らかさや参照モデルからの差への拘束である\cite{pygimli_inversion}。データを誤差より細かく再現しても、雑音をモデルへ取り込んでいる可能性がある。

Marquardt型の減衰付き更新、共役勾配、遺伝的探索はモデル更新の算法である。正則化・平滑化拘束は逆解の安定性を制御する。有限要素法や差分法は支配方程式の離散化を担い、\(F(m)\) や感度を計算するため逆解析の内部でも利用できる。平成30年度問35の公式5は直接の役割の比較として読むが、FEMがインバージョンに無関係という絶対命題にはしない。

\subsection{波動の基本則と等価層・検層}
\label{theory:exploration:principles}
線形波動では重ね合わせが成立する。相反性を持つ媒体では震源と受振点を交換した応答に対応がある。ホイヘンスは波面構成、スネルは境界の位相整合による角度と波速、フェルマは波線の走時の停留条件を表す。最短時間という単語だけで複数経路の可能性を排除しない。レンツの法則は磁束変化を妨げる電磁誘導の向きに関する法則であり、弾性波の屈折則とは異なる。

USGSの等価層は重力の密度層や磁気の双極子層によってポテンシャル場を表す説明である\cite{usgs_equivalent2007}。場の表現を波動伝播の基本則と混同しない。ラテラル法は孔内の電極間隔と配置による電気検層の方式である\cite{usgs_lateral_logging}。計測配置の名称と、取得後のフィルタ・差分・最適化の算法を区別する。

\section{屈折・反射・表面波：地下を拘束する情報の違い}
\label{theory:exploration:seismic}
地震探査の記録には直接波、屈折波、反射波、表面波が重なる。屈折法では通常初動走時、反射法では反射イベントの走時と波形、表面波法では周波数ごとの位相速度を主に使う。名前の違いは処理だけではなく、地下構造を拘束する観測情報の違いである。

\subsection{速度と地質：波種・状態・境界を分ける}
\label{theory:exploration:waveproperties}
安定な等方弾性体では、体積弾性率 \(K\)、せん断弾性率 \(G\)、密度 \(\rho\) を使い、
\[
 V_P=\sqrt{\frac{K+4G/3}{\rho}},\qquad
 V_S=\sqrt{\frac{G}{\rho}}
\]
となる。\(K>0,G>0\) なら \(V_P>V_S\) である。流体ではせん断を保てず通常のS波が伝わらないので、二つの伝播速度を常に観測して比較できるとは限らない。

岩種が同じでも、割れ目の開口・接触、風化、変質、空隙、飽和が有効弾性係数を変える。空隙の空気を水が置き換えると圧縮性等へ影響し、P波速度は変わり得る。USGSの測定報告では飽和帯の深さとP波速度を読み、表面波のS波逆解析を拘束する\cite{usgs_refraction2013}。地質学の境界は岩種・年代・堆積単位等の区分であり、速度差の境界と同一ではない。速度から風化や破砕を推定するとき、露頭やボーリングと併せて解釈する。

通常の浅部屈折ではP波が多く使われるが、S波屈折も存在する\cite{epa_seismic_refraction}。土質のS波調査を比較的平坦な地形で行うという運用の目安を、急峻地形では物理的に不可能という命題へ変更しない。強震動にはS波構造が重要であり、通常の地山P波屈折の主用途と、複数手法の統合による寄与を分ける\cite{segj_integrated2017,japansea_spac2017}。

\subsection{屈折初動：はぎとりと往復の整合}
\label{theory:exploration:refraction}
上層速度 \(V_1\)、下層速度 \(V_2>V_1\)、上層厚 \(h\) の単純な水平二層では、\(\sin i_c=V_1/V_2\) を満たす臨界屈折波の走時は
\[
 t_{\rm head}(x)=t_i+\frac{x}{V_2},\qquad
 t_i=\frac{2h\cos i_c}{V_1}
\]
となる。直接波の \(x/V_1\) と比べて最初に到達する時刻を読む。原点走時 \(t_i\) は走時直線の外挿切片であり、起振位置の実測初動ではない。臨界距離と直接波を追い越す交差距離も同じ概念ではない\cite{epa_seismic_refraction}。

萩原のはぎとり法は表層の走時寄与を分離して下層の速度と境界を求める屈折初動の方法である。JICA調査団の比較表は直接波＋屈折波を扱い、トモグラフィとの構造・速度の仮定を分けている\cite{jica_refraction2015}。反射波の走時をそのまま同じはぎとりへ入力する説明は誤りである。

端点AからBへの全走時とBからAへの全走時は、相反的な同じ経路・媒体で整合する。全走時一致はこの往復整合などを確かめる意味であり、全ての受振点の到達時刻が同じという意味ではない。原点走時の整合、曲線の平行性も観測モデルに沿って確認する。薄い中間層は屈折波があっても初動として現れない距離範囲しか持たない場合がある。低速度逆転、薄いブラインド層、連続速度勾配を分ける\cite{usgs_refraction2013}。

連続的に深度増速するミラージュ構造では波線が曲がり、走時は折れ線ではなく曲線となる。トモグラフィは波線が通る格子速度又は遅さを反復更新し、連続勾配も扱える。ただし通らない領域の速度を一意に決めることはできない。測線は地形と地質の変化を見ながら設計し、等高線に垂直で一定勾配という望ましい典型配置だけを絶対条件にしない。

\subsection{反射走時：双曲線と放物線の適用範囲}
\label{theory:exploration:reflectiongeometry}
単一水平反射面、上部の一定速度 \(V\)、深さ \(h\)、震源と受振器の距離 \(x\) なら、反射点は両者の中点にある。左右対称の経路長を足して独自に導くと、
\[
 t(x)=\frac{2}{V}\sqrt{h^2+(x/2)^2},\qquad
 t^2=t_0^2+\frac{x^2}{V^2},\quad t_0=\frac{2h}{V}.
\]
\(t\)-\(x\) 平面では双曲線である。多層では小〜中オフセットの範囲で有効な重合速度を使って双曲線近似する。大オフセット、強い不均質、異方性等には非双曲線成分が残る。

さらに \(|x|\ll Vt_0\) なら平方根を展開して、
\[
 t(x)=t_0+\frac{x^2}{2V^2t_0}
 -\frac{x^4}{8V^4t_0^3}+\cdots
\]
となる。二次項までの近似は放物線である。このため平成28年度問21と令和5年度問32の公式選択を、放物線ではどんな近似もできないという命題へ書き換えない。通常の走時モデルが双曲線であることと、原文に近似範囲の指定がない問題点を併記する。

NMOは各オフセットの走時をゼロオフセット走時へそろえる。重合速度はその目的に最適な速度であり、真の区間速度と無条件に同じではない\cite{usgs_processing2000}。反射とは異なるmoveoutの多重反射等は重合で弱まる場合があり、相関雑音も全て残るとは限らない。同じmoveoutの雑音は抑えにくい。

浅層反射の最大オフセットを目標深度程度とする説明は通常の展開目安である\cite{coastal_reflection2016}。目標反射角、近接雑音、速度解析、深度、開口と重合数で設計を変える。二倍を一律に必要とする説明を誤りと選ぶ場合も、二倍の観測が常に不適切という意味ではない。

\subsection{分解能と検出：薄層が存在することを知る}
\label{theory:exploration:resolution}
厚さ \(\Delta h\) の上下面反射の時間差は単純な垂直経路で \(\Delta t=2\Delta h/V\) である。二つのwaveletが強く重なると別々の反射として分離しにくくなる。代表的な垂直分解能の目安 \(\Delta h\sim\lambda/4\) は、上面と下面の分離に関するものである\cite{seg_vertical_resolution}。

分離できない薄層も、合成振幅や周波数応答の変化として存在を検出できる場合がある。検出にはS/N、インピーダンス差、位相、帯域が関わる。「四分の一波長より薄ければ必ず見えない」と断定しない。Rayleigh基準やRicker基準、waveletは反射イベントの分離に関係し、屈折直線のintercept timeとは役割が異なる。

横方向にはFresnel zoneの有限範囲から反射が寄与する\cite{seg_lateral_resolution}。単純なゼロオフセットの幾何では往復経路差を半波長とする一次zoneから \(r_F\approx\sqrt{h\lambda/2}\) の目安を独自に得る。migration、開口、帯域と速度モデルによって横方向の分解能は変わる。zone径を全処理の最終解像度と固定しない。

\subsection{表面波分散とS波モデル}
\label{theory:exploration:surfacewaves}
Rayleigh波は表面に沿って進み、深度へ振幅が減少する。不均質な成層では周波数により感度深度が変わるため位相速度 \(c(f)\) が分散する。長波長は深部、短波長は浅部を拘束する傾向があるが、一つの波長が一つの深度だけを測るわけではない\cite{epa_surfacewaves}。均質等方半無限体は非分散という退化例である。

観測位相速度から地下モデルを作り、理論分散曲線を計算して残差を小さくする。位相速度は \(V_P,V_S,\rho\)、層厚に依存するが、通常の地盤分散は \(V_S\) に強い感度を持つ。Xiaらの一次論文要約は高周波側でS波速度と層厚の影響が大きいことを述べる\cite{xia_rayleigh1999}。P波・密度の影響をゼロにせず、別の測定や拘束を使う。

水平多層の一次元逆解析は通常のモデルであり、地盤自体が常に横方向に均質と分かったわけではない。高次モード、方向による変化、モデル自由度とデータ数を確認する。能動表面波の通常の浅部対象を20m程度とする資料がある\cite{moe_surface2004}。しかし20mや100mは普遍的上限ではなく、震源の長波長、受振器の低周波応答、配置、雑音と解析範囲に依存する。微動で長波長側を補えれば深部拘束を増やせるが、併用だけで深度増加を保証しない。

\subsection{微動アレイの短波長と長波長の限界}
\label{theory:exploration:arraylimits}
空間標本化で間隔 \(d\) に対し波長が短すぎると、異なる空間波数が同じ標本を与える。等間隔・単方向の単純モデルでは \(\lambda>2d\) の条件が基本となる。実アレイの限界は配置と波の方向にも依存する。空間aliasingは主に最短波長側の制約であり、最長側と取り違えない\cite{japansea_spac2017}。

長波長ではアレイ内の位相差が小さく、有限開口、ノイズ、推定算法と観測時間の制約が効く。F-Kは周波数ごとの波数分布を求め、SPACは空間自己相関と分散の関係を使う。通常の条件でSPACがF-Kより長波長まで推定する例が一次論文にあるが、入射方向と多数の波、構造、配置、実装に依存する\cite{array_limits2002}。一定の半径倍率を全観測の保証値へ転記しない。複数サイズで重なる帯域を確認し、疑わしい分散曲線を機械的に一つのモデルへ合わせない。

\section{電気探査：地下へ流した電流から何が分かるか}
\label{theory:exploration:dcarrays}
地層を掘る前に、水や粘土、空隙、金属などによる電流の流れやすさの差を調べたい。比抵抗法は電流を注入して地表や孔内の電位差を測り、地下の抵抗率分布を推定する。抵抗\(R\)は試料の形状も含むが、均質な棒では\(R=\rho L/A\)である。長さ\(L\)、断面積\(A\)から\(\rho=RA/L\)となるため、抵抗率の単位は\(\Omega\,\mathrm m\)であり、\(\Omega\,\mathrm m^3\)ではない。局所形では\(\boldsymbol E=\rho\boldsymbol J\)で、導電率は\(\sigma=1/\rho\)である。

定常電流の保存と\(\boldsymbol E=-\nabla V\)を合わせると、源のない部分では
\[
 \nabla\cdot(\sigma\nabla V)=0
\]
となる。注入電極を源・回収電極を吸込みとして、地表の境界条件と地形を含めて電位を計算するのが順問題である。多数の電位差へ一致する\(\sigma\)を探すのが逆問題で、電位だけで地質名が一意に決まるわけではない\cite{epa_resistivity}。

\subsection*{半空間の電位とWennerの幾何係数}
教材の説明用に均質・等方な導電半空間、絶縁した平らな地表、小さい点電極を仮定する。地表の注入点から半球面を通る電流を\(I\)とすると、\(J_r=I/(2\pi r^2)\)である。したがって無限遠基準の電位は\(V(r)=\rho I/(2\pi r)\)となる。注入Aと回収Bの重ね合わせより、電位電極M,Nの差は
\[
 \Delta V=\frac{\rho I}{2\pi}
 \left(\frac1{AM}-\frac1{BM}-\frac1{AN}+\frac1{BN}\right).
\]
Wenner配置を地表直線上のA--M--N--B、全隣接間隔\(a\)とすると、括弧は\(1/a\)であり
\[
 \rho_a=2\pi a\frac{\Delta V}{I}
\]
を得る。この導出は教材独自である。符号は電極と電位差の定義をそろえる。不均質な地盤の\(\rho_a\)は「この配列で同じV/Iを出す均質半空間」の値であり、MとNの真下にある一層の真の抵抗率ではない。

Schlumberger、dipole--dipole等も幾何係数と地下感度が異なる。垂直探査では中心をおおむね固定し間隔を広げ、水平多層等のモデルを推定する。水平探査では配置を移動して横の変化を調べる。固定間隔でも感度は有限体積に分布するので「特定の深度だけの水平抵抗率」を直接測ると読まない。二次元・三次元の比抵抗トモグラフィは、多数配置のデータを順モデルへ合わせて分布を求める。偽断面の表示位置を実深度と同一視しない。

\subsection*{線形フィルタは配置でなく順応答の計算法}
層状大地モデルの比抵抗変換から見掛比抵抗曲線を計算する際、線形フィルタを用いる方法がある。USGSのTI-59プログラムはその具体例である\cite{usgs_linear_filter1981}。平成23年度問30では、配列名に混じるリニアフィルタを一般の雑音除去だけと解釈しない。IPは追加の分極応答を測るが、比抵抗と同じ電極を使うため、無関係という全称も避ける。

\subsection*{四電極で分ける接触電圧降下}
\label{theory:exploration:fourpoint}
二端子で電流を流しその端子電圧を測ると、地盤だけでなく両電極の接触抵抗による電圧降下も含む。別のM,Nを高入力抵抗\(R_{\rm in}\)の電圧計へつなぎ、測定電流\(I_v\)を十分小さくすれば、接触誤差の代表値\(I_v(R_{cM}+R_{cN})\)を小さくできる。これが四電極の要点である。送信側の接地抵抗は電流を確保する問題として残り、四電極なら接地・地形・電磁カップリングの影響が全てゼロになるわけではない。平成25年度問34はこの分離を問う。

\section{含水岩石の電流経路とArchie式}
\label{theory:exploration:archie}
多くの岩石骨格は絶縁性が高く、孔隙水に溶けたイオンが電流を運ぶ。孔隙率\(\phi\)、水飽和度\(S_w\)、水の抵抗率\(\rho_w\)について、代表的な実験モデルは
\[
 \rho_t=a\rho_w\phi^{-m}S_w^{-n}
\]
である。\(a,m,n\)は骨格・連結性・濡れ条件等に依存する係数で、全岩石に同じ値を置かない。飽和では\(\rho_t/\rho_w=a\phi^{-m}\)が地層係数となる。孔隙ネットワークの少ない導電断面と迂回路をまとめたモデルであり、普遍的な保存則から全係数が決まる式ではない。

粘土表面、導電鉱物、電子伝導に無視できない電流があれば、孔隙水のみという前提が破れる。単純な並列伝導モデル\(\sigma_{\rm bulk}=\sigma_{\rm pore}+\sigma_{\rm surface}\)を考えると、追加の表面伝導は抵抗率を下げる。低抵抗の水は相対的に孔隙水伝導を優勢にしやすい。逆に高抵抗水だけでは式の精度を保証できない。ただし清浄な絶縁骨格なら淡水でも利用し得る\cite{epa_borehole_resistivity}。

海水の塩分は通常導電率を上げ抵抗率を下げる。不飽和化は連続水経路を減らし高抵抗化しやすい。粘土化・含水割れ目は低抵抗化しやすいが、乾燥した割れ目や珪化まで同じ傾向にはしない。「礫層は高抵抗」という典型も周辺の粘土層との対比・塩分条件による。原典の一般傾向を、岩石名から値を一意に決める規則にしない。

\section{IPとSP：人工分極と自然電位の機構}
\label{theory:exploration:ipsp}
IP（誘導分極）は人工電流を流した際の地下界面の電荷蓄積を利用する。電流遮断後の電位減衰又は周波数応答の位相を測る。硫化鉱物の電子伝導と電解質のイオン伝導の界面、鉱物表面等が応答へ寄与する。比抵抗だけでは区別しにくい界面状態を調べるための追加情報であり、IPでも比抵抗を測る。ケーブル間の電磁結合が分極応答へ混ざり得るため、周波数・配置・測定時間を吟味する\cite{epa_ip}。

SP（自然電位）は人工的に電流を注入せず、基準点との電位差を非分極電極で測る。地下水流による電気二重層の電荷輸送、酸化還元差を結ぶ導電性鉱体、温度や化学勾配等が源となる。説明用に\(\boldsymbol J=-\sigma\nabla V+\boldsymbol J_s\)と置けば、電流保存は
\[
 \nabla\cdot(\sigma\nabla V)=\nabla\cdot\boldsymbol J_s
\]
となる。観測から源\(\boldsymbol J_s\)を推定するには抵抗率と源機構のモデルが必要である。電位の絶対零点は測らず、基準点の選び方も記録する。

典型的な硫化鉱体の鉱化電池では地表の負異常を利用する。一方、流動・熱・化学が重なる場の符号は一つの法則へ固定できない。正異常だから必ず熱水、負異常だから必ず鉱体と断定しない。SP異常の幅や幾何モデル、他観測との同時逆解析による深度推定は可能だが、SP分布だけで深度が一意に読めるわけではない\cite{epa_sp}。令和3年度問32の公式4は保持するが、「SPでは深度解析ができない」という誤った禁止を教材にしない。原典の否定範囲に残る疑義を問題ごとに示す。

\section{有限送信源・周波数領域・時間領域}
\label{theory:exploration:nearfield}
自然MTと人工CSAMTの抵抗率・表皮深度は前節\ref{theory:exploration:electromagnetic}の式へつながる。均質導電半空間で変位電流を無視すると、周期場の拡散による振幅減衰尺度は\(\delta=\sqrt{2\rho/(\omega\mu)}\)である。送受信距離\(r\)に対して\(r/\delta\)が十分大きい遠方では平面波近似を使いやすい。固定rで低周波化又は高抵抗化すればδが増し、近接場へ移る。USGS実地報告は三表皮深度超を目安として用いるが、その数を全装置・全構造の厳密閾値とはしない\cite{usgs_csamt2012}。

\(\rho_a=|E/H|^2/(\omega\mu)\)のµは透磁率である。真空ではµ0を用い、誘電率ε0ではない。HをBと取り違えると換算係数と単位も変わる。平成28年度問23は式の形だけでなく、この変数定義の誤りを問う。

\subsection*{周波数領域では振幅と位相を測る}
\label{theory:exploration:fdtd}
FDEMでは正弦電流が送信コイルの一次磁場を作り、地盤の渦電流が二次磁場を生む。定常応答の複素振幅比を同相・直交成分又は振幅・位相として測る。Slingramやloop--loopはこの仕組みを使う。電極を挿さなくてもよいが、コイル間隔、周波数、導電率、姿勢に感度が依存する\cite{epa_fdem}。

TDEM/TEMでは送信を遮断した後の過渡応答を測る。均質導電体の準静的近似では\(\partial\boldsymbol B/\partial t=(1/(\mu\sigma))\nabla^2\boldsymbol B\)となるため、代表拡散距離は\(\sqrt{t/(\mu\sigma)}\)の尺度で増す。遅い時刻には一般に深部への感度があるが、信号は弱くなり、装置寸法・電流・遮断時間・測定時間・ノイズが限界を決める。central loopで深部を測れることを「送受信間隔によらず無制限に深く測れる」としない\cite{epa_tdem_archive}。一次場の遮断後観測と、電場の局所歪みによるMTのstatic shiftの少なさは別の性質である。

VLFは既存通信局の3--30kHz帯の信号を使う代表的方式で、局の方位と対象構造の向きが重要である。調査者は既存局の送信条件を任意には変更できない。Schumann共振も電気電磁現象であり、MT観測記録で調べられている\cite{schumann_mt2009}。反射地震のNMO補正とは分野が異なる。

\section{GPR：走時断面から電気境界を推定する}
\label{theory:exploration:gpr}
地中レーダは高周波電磁パルスを放射し、電気特性の異なる境界からの反射を受ける。低損失・非磁性の均質媒質なら
\[
 v\simeq\frac{1}{\sqrt{\mu_0\epsilon}}
 =\frac{c}{\sqrt{\epsilon_r}},\qquad
 \epsilon_r=\frac{\epsilon}{\epsilon_0},\qquad
 z\simeq\frac{vt}{2}
\]
である。tは往復走時であり、直下水平面などの単純な幾何と速度が必要である。εはF/m、εrは無次元で、空気を分母にするのは通常の近似に過ぎない。高導電性・分散媒質ではこの単純速度式の条件が破れる。

波形をアンテナ位置へ並べるprofileはまず時間断面である。wide-angleや共通中点配置から速度を推定し、深度変換に用いる。埋設管等を横切る測線は位置を捉えやすい。点対象の回折双曲線をそのまま地層形状と思わず、速度・偏位・必要なmigrationを考える\cite{epa_gpr}。

\subsection*{分解能力と最小識別距離を区別する}
\label{theory:exploration:gprresolution}
二つの反射を時間差\(\Delta t\)で識別するにはパルス幅・帯域が制約になる。説明用には帯域Bで\(\Delta t\sim1/B\)、距離分解尺度\(\Delta z\sim v/(2B)\)である。帯域が中心周波数に比例する同系統装置なら、高周波ほど最小識別距離が小さく、能力としての「分解能」は良い。この語の二つの使い方で比例の向きが変わるので、式と変数を明示する。

実際の探査深度は導電減衰、散乱、幾何拡がり、アンテナ帯域、検出閾値等で決まる。高周波で分解を良くすると到達深度が制限されるという交換条件はあるが、深度が厳密に1/fへ比例する普遍式ではない。低損失非分散媒質の導電損失近似は\(\alpha\simeq(\sigma/2)\sqrt{\mu/\epsilon}\)で、fだけの反比例を導くものではない。

通常のGPRはMHz以上のパルスで、原典の10MHz程度は分類目安である。10kHzとの三桁差は誤りだが、10MHzに絶対的な物理境界があるわけではない。力学強度を直接観測しない一方、含水・空隙・劣化等の共通要因による間接相関はあり得る。「ほとんど相関なし」という原典はその限定を保持する。数mの空洞や埋設物の用途を、全媒質の到達上限へ変えない。

\section{孔内検層：深度目盛と観測物性を結び付ける}
\label{theory:exploration:borehole}
地表探査だけでは深度の曖昧さが残る。ボーリング孔に沿う連続計測は、その位置の深度と物性応答を結び付ける。ただし孔周辺の有限体積を見ており、孔壁・泥水侵入・ケーシング・装置結合の補正がある。検層記録の深度が分かることと、地盤全体の構造が完全に分かることは別である。

速度検層のdownholeは地表起振・孔内受振、upholeは孔内起振・地表受振である。走時から速度を求めるには距離と伝播経路、屈折・傾斜の条件を扱う。孔壁へ密着させるdownhole受振器なら水位のない部分でも測定できる\cite{geebec_ps}。水位以深だけという限定は誤りだが、孔内状態によらず精度が保証される意味ではない。

suspensionは孔内振源と複数受振器を持つゾンデで、受振点間の時間差から区間速度を得る\cite{oyo_suspension}。地盤のP/S応答を測り、水中をS波が伝わるという説明にはしない。VSPでは通常の地表振源・孔内受振と逆配置等を区別し、反射記録と孔内深度・地層を対比する。初動速度だけの検層と反射の解析目的は同じではない。

normal電気検層は電極の注入・電位差による比抵抗計測である。導電水・泥水を通じて地層へ電流を流す通常の裸孔条件を考える。孔内水の有無・抵抗率を無視しない。一方、コイルの誘導検層には別の接触条件があり、電極型の水位条件を全検層へ拡張しない。P波速度は弾性率と密度に関わり、含水・飽和で変わり得るが、含水比だけで全地盤の増減が決まるわけではない。

\subsection*{ガンマ密度：粒子密度でなくかさ密度}
\label{theory:exploration:gammadensity}
通常の密度検層はガンマ線を放射し、主にCompton散乱を介する応答から電子密度を測る。一般的岩石の組成では校正によりかさ密度へ結び付く。ベータ線やRayleigh散乱へ置き換えない。飽和二成分モデルなら
\[
 \rho_b=(1-\phi)\rho_m+\phi\rho_f,\qquad
 \phi=\frac{\rho_m-\rho_b}{\rho_m-\rho_f}
\]
である。\(\rho_m\)は粒子・骨格、\(\rho_f\)は孔隙流体の密度であり、これらが既知という条件から孔隙率を推定する。粒子密度そのものの直接測定ではない。孔壁接触や孔内流体、線源・検出器校正が必要である\cite{epa_gamma_density}。

\section{時間差・幾何・熱を使う補足調査}
\label{theory:exploration:multimethod}
time lapseは同じ地点・方式の反復測定を比較し、地下水移動や物性変化を推定する。差分には機器ドリフトや測線ずれ・温度差も入り得る。tomographyの覆いが不足すれば、特定方向の構造や薄い領域を分解しにくい。重力では同じ形状深度なら密度差の大きさが異常を増すが、密度差だけで異なる対象の検出難易度を順位付けしない。

1m深地温探査は、地表下1mの温度を測り流動水・熱源による平常温度との差を調べる。1mは測温深度であり、水脈の最大深度ではない。水温との温度差が小さい季節には検出しにくく、地況・境界条件を扱う熱モデルと観測を組み合わせる\cite{naro_thermal2008,landslide_thermal2007}。特定地の適期を全国共通の月にしない。

環境振動の振動レベルは人の感覚に応じた周波数補正を施す加速度レベルである。原典のJIS準拠測定を、単なる最大変位測定へ取り違えない\cite{rion_vibration2016}。出題時と現行の規格本文の版差は未照合であり、現在版の番号や性能要件を推測して補わない。

\subsection*{試験での確認順序}
励起源、受信器、直接観測量、推定物性、単位、有限距離・帯域・媒質の条件を順に照合する。今回の29問での論点を全16試験回の頻度の確定値とはしない。SPの深度推定、Archieの表面伝導、GPRの比例表現等の疑義は各問題へ残し、公式正答と科学的な説明の保証範囲を区別する。

\section{地震・弾性波探査}

\subsection{平成23年度 問29}
\label{q:h23-reg-applied-29}

\noindent\textbf{原典：} sources/H23/17\_applied\_science.pdf、PDF物理10頁。

\textbf{公式正答：} 5（single）。

\paragraph{問題の要旨} 物理探査の観測対象と装置を比較し、海上音波探査の受振器を誤って説明した記述を選ぶ。

\Needspace{4\baselineskip}\paragraph{確認した条件・記述（独自要約）}\mbox{}\par

\noindent\textbf{選択方向：} 誤り又は最も不適切・関係の薄いものを1つ\par

\paragraph{問われる核心} 観測量と受振器の対応

\paragraph{必要知識} ハイドロフォンの圧力応答、地上速度計、相対重力計、MT、電磁波の反射。

\paragraph{考え方・正答を選ぶ理由（査読待ち初稿）} 測定量を対照する。重力は加速度、MTは自然電磁応答、レーダは電磁反射である。水中の流体はせん断波を伝えず、海上の受振系では音圧を検出するハイドロフォンを用いる。速度計を防水加工するという第5肢の説明が誤りである。

\Needspace{4\baselineskip}\paragraph{選択肢の検討（査読待ち初稿）}\mbox{}\par

\noindent\textbf{1：適切} 磁気から鉄類の有無と概略位置を探るが、細かな形状や製品種類は判別できない。通常の磁気異常だけでは磁化・形状・深度が非一意であり、細かな製品識別までは確定できないという探査上の限界を述べる。現代の補助情報を用いる識別可能性全般を否定する命題とは解さない。\par

\noindent\textbf{2：適切} MTは自然の地磁気・地電流から深部比抵抗を探る。自然変動の水平電場・磁場の応答から電気抵抗率を推定する方法である。\par

\noindent\textbf{3：適切} 相対重力計は吊った錘のばね伸びを利用する。ばね定数と伸びの変化から重力の差を測る。絶対値を直接定める落下式とは区別する。\par

\noindent\textbf{4：適切} 地中レーダは概ね10 MHz以上の電磁波で極浅部を探る。原典が述べる通常の浅部レーダの周波数帯と原理に対応する。到達深度は周波数だけでなく導電率や減衰にも依存する。\par

\noindent\textbf{5：不適切} ストリーマケーブルには防水加工した速度計を多数入れる。水中音波を測る受振器は通常ハイドロフォンであり、圧力変動を検出する。地上用速度計の防水化と同一ではない。\par

\paragraph{再利用する考え方} 観測量と受振器の対応を、測定量・推定対象・基準面・近似条件に分けて理解する。名称だけで探査法を対応させない。

\paragraph{共通理論} \hyperref[theory:exploration:observables]{第\ref*{theory:exploration:observables}節}, \hyperref[theory:exploration:electromagnetic]{第\ref*{theory:exploration:electromagnetic}節}, \hyperref[theory:gravity:measurement]{第\ref*{theory:gravity:measurement}節}

\par 参考：\cite{usgs_marine_seismic,usgs_mt_metadata}。

\par 難易度 3（校準初稿）、優先度未評価。

\section{重力・磁気探査}

\subsection{平成27年度 問22}
\label{q:h27-reg-applied-22}

\noindent\textbf{原典：} sources/H27/17\_applied\_science.pdf、PDF物理8頁。

\textbf{公式正答：} 2（single）。

\paragraph{問題の要旨} 地形、無限水平層、潮汐、ばねドリフト、緯度の重力効果から誤りを選ぶ。

\Needspace{4\baselineskip}\paragraph{確認した条件・記述（独自要約）}\mbox{}\par

\noindent\textbf{選択方向：} 誤り又は最も不適切・関係の薄いものを1つ\par

\paragraph{問われる核心} 無限平板の重力と重力補正の基準

\paragraph{必要知識} ニュートン引力積分、無限平板と有限岩体の違い、地形補正の符号、潮汐、ドリフト、正規重力。

\paragraph{考え方・正答を選ぶ理由（査読待ち初稿）} 無限水平層の上に観測点がある場合、各薄層の寄与を積分すると\(g_z=2\pi G\rho t\)である。層と測点の距離は消えるので、浅いほど大きいという第2肢が誤りである。第1肢は無限平板に対する地形差、第4肢は短時間の線形ドリフト近似として条件を明示して読む。

\Needspace{4\baselineskip}\paragraph{選択肢の検討（査読待ち初稿）}\mbox{}\par

\noindent\textbf{1：適切} 山・谷の地形効果はともに重力を小さくする。測点高まで無限平板があるというブーゲー基準に対する地形効果として読む。測点より上の山の引力は上向き、谷は基準平板の下向き引力を欠くため、いずれも下向き重力を減らす。任意の山と谷の絶対重力を比較する命題ではない。\par

\noindent\textbf{2：不適切} 密度と厚さ一定の無限水平層は浅いほど重力が大きい。無限平板の外部重力は2πGρtで深度に依存しない。有限の岩体では距離依存があるため、無限という条件が重要である。\par

\noindent\textbf{3：適切} 月と太陽の引力による潮汐補正が必要である。月・太陽による時間変動を地下構造による異常と区別する必要がある。\par

\noindent\textbf{4：適切} ばねの時間変化を一定と近似し閉合誤差を各点へ配分する。同一点への戻り測定の閉合差から時間に比例するドリフトを近似して補正する操作を述べる。一定の伸びは測定区間のドリフトを一定率とする通常の近似として解釈し、あらゆる時間変化が必ず線形とはしない。\par

\noindent\textbf{5：適切} 絶対重力と緯度から求めた正規重力との差をとるのが緯度補正である。緯度に依存する理論基準を差し引き、基準地球の変化を除く操作である。\par

\paragraph{再利用する考え方} 無限平板の重力と重力補正の基準を、測定量・推定対象・基準面・近似条件に分けて理解する。名称だけで探査法を対応させない。

\paragraph{共通理論} \hyperref[theory:gravity:slab]{第\ref*{theory:gravity:slab}節}, \hyperref[theory:gravity:corrections]{第\ref*{theory:gravity:corrections}節}, \hyperref[theory:gravity:measurement]{第\ref*{theory:gravity:measurement}節}

\par 参考：\cite{gsi_gravity_definitions}。

\paragraph{関連問題} \hyperref[q:r07-reg-applied-33]{R07-REG-APPLIED-33}

\par 難易度 3（校準初稿）、優先度未評価。

\subsection{平成29年度 問31}
\label{q:h29-reg-applied-31}

\noindent\textbf{原典：} sources/H29/17\_applied\_science.pdf、PDF物理10頁。

\textbf{公式正答：} 5（single）。

\paragraph{問題の要旨} 重力探査で用いる補正名のうち、主に反射法地震探査に属するものを選ぶ。

\Needspace{4\baselineskip}\paragraph{確認した条件・記述（独自要約）}\mbox{}\par

\noindent\textbf{選択方向：} 誤り又は最も不適切・関係の薄いものを1つ\par

\paragraph{問われる核心} 重力補正と地震走時補正の区別

\paragraph{必要知識} 運動・質量・潮汐・標高の重力効果、浅部層による走時差。

\paragraph{考え方・正答を選ぶ理由（査読待ち初稿）} 補正の対象量を並べると、第1〜4肢は重力加速度の変化を扱う。静補正は主に地震探査の走時をそろえるもので、重力探査の標準的補正名としては該当しない。

\Needspace{4\baselineskip}\paragraph{選択肢の検討（査読待ち初稿）}\mbox{}\par

\noindent\textbf{1：適切} エトベス補正。船や航空機の運動に伴う回転系の見掛加速度の差を補正する。\par

\noindent\textbf{2：適切} ブーゲー補正。基準面と測点の間の岩石質量の引力を除く。\par

\noindent\textbf{3：適切} 潮汐補正。月・太陽による時間変動を除く。\par

\noindent\textbf{4：適切} フリーエア補正。基準面からの高さに伴う重力変化を補正する。\par

\noindent\textbf{5：不適切} 静補正。浅部低速度層や標高による地震波の走時差を、共通の基準面へ補正する地震探査の用語である。\par

\paragraph{再利用する考え方} 重力補正と地震走時補正の区別を、測定量・推定対象・基準面・近似条件に分けて理解する。名称だけで探査法を対応させない。

\paragraph{共通理論} \hyperref[theory:gravity:corrections]{第\ref*{theory:gravity:corrections}節}, \hyperref[theory:exploration:observables]{第\ref*{theory:exploration:observables}節}

\par 参考：\cite{gsi_gravity_definitions}。

\par 難易度 3（校準初稿）、優先度未評価。

\subsection{平成30年度 問32}
\label{q:h30-reg-applied-32}

\noindent\textbf{原典：} sources/H30/17\_applied\_science.pdf、PDF物理11頁。

\textbf{公式正答：} 3（single）。

\paragraph{問題の要旨} 重力・磁気探査の基準面、磁化特性、移動観測補正と、別の観測分野の語を区別する。

\Needspace{4\baselineskip}\paragraph{確認した条件・記述（独自要約）}\mbox{}\par

\noindent\textbf{選択方向：} 誤り又は最も不適切・関係の薄いものを1つ\par

\paragraph{問われる核心} 重力・磁気探査の物理量と用語

\paragraph{必要知識} 等ポテンシャル面、残留磁化と誘導磁化、磁気秩序の温度依存、回転座標系。

\paragraph{考え方・正答を選ぶ理由（査読待ち初稿）} ジオイドとエトベス補正は重力、ケーニヒスベルガー比とキュリー温度は磁気に直接関係する。プランク関数は熱平衡放射のスペクトルを記述する関数であり、この候補群では重力・磁気探査の基本量に当たらない。したがってこの肢を選ぶ。

\Needspace{4\baselineskip}\paragraph{選択肢の検討（査読待ち初稿）}\mbox{}\par

\noindent\textbf{1：適切} ジオイド。基準となる重力等ポテンシャル面であり、重力場と標高の基準に関係する。\par

\noindent\textbf{2：適切} ケーニヒスベルガー比。自然残留磁化の大きさと現在の場による誘導磁化の大きさの比であり、磁気異常の解釈に関係する。\par

\noindent\textbf{3：不適切} プランク関数。熱平衡放射のスペクトルを記述する関数であり、この候補群では重力・磁気探査の基本量に当たらない。\par

\noindent\textbf{4：適切} キュリー温度。強磁性・フェリ磁性の自発磁化が消失する温度に対応し、磁性鉱物や磁気探査の解釈に関係する。\par

\noindent\textbf{5：適切} エトベス補正。移動観測系の運動に伴う見掛重力の差を補正する。\par

\paragraph{再利用する考え方} 重力・磁気探査の物理量と用語を、測定量・推定対象・基準面・近似条件に分けて理解する。名称だけで探査法を対応させない。

\paragraph{共通理論} \hyperref[theory:gravity:corrections]{第\ref*{theory:gravity:corrections}節}, \hyperref[theory:gravity:geoid]{第\ref*{theory:gravity:geoid}節}, \hyperref[theory:geomagnetism:magnetization]{第\ref*{theory:geomagnetism:magnetization}節}, \hyperref[theory:exploration:observables]{第\ref*{theory:exploration:observables}節}

\par 参考：\cite{gsi_gravity_definitions,kanaya_okuma2011}。

\paragraph{関連問題} \hyperref[q:r04-reg-applied-35]{R04-REG-APPLIED-35}

\par 難易度 3（校準初稿）、優先度未評価。

\subsection{令和4年度 問35}
\label{q:r04-reg-applied-35}

\noindent\textbf{原典：} sources/R04/17\_applied\_science.pdf、PDF物理12頁。

\textbf{公式正答：} 4（single）。

\paragraph{問題の要旨} 重力・磁気探査の基準面、磁化特性、移動観測補正と、別の観測分野の語を区別する。

\Needspace{4\baselineskip}\paragraph{確認した条件・記述（独自要約）}\mbox{}\par

\noindent\textbf{選択方向：} 誤り又は最も不適切・関係の薄いものを1つ\par

\paragraph{問われる核心} 重力・磁気探査の物理量と用語

\paragraph{必要知識} 等ポテンシャル面、残留磁化と誘導磁化、磁気秩序の温度依存、回転座標系。

\paragraph{考え方・正答を選ぶ理由（査読待ち初稿）} ジオイドとエトベス補正は重力、ケーニヒスベルガー比とキュリー温度は磁気に直接関係する。プランク関数は熱平衡放射のスペクトルを記述する関数であり、この候補群では重力・磁気探査の基本量に当たらない。したがってこの肢を選ぶ。

\Needspace{4\baselineskip}\paragraph{選択肢の検討（査読待ち初稿）}\mbox{}\par

\noindent\textbf{1：適切} エトベス補正。移動観測系の運動に伴う見掛重力の差を補正する。\par

\noindent\textbf{2：適切} ジオイド。基準となる重力等ポテンシャル面であり、重力場と標高の基準に関係する。\par

\noindent\textbf{3：適切} キュリー温度。強磁性・フェリ磁性の自発磁化が消失する温度に対応し、磁性鉱物や磁気探査の解釈に関係する。\par

\noindent\textbf{4：不適切} プランク関数。熱平衡放射のスペクトルを記述する関数であり、この候補群では重力・磁気探査の基本量に当たらない。\par

\noindent\textbf{5：適切} ケーニヒスベルガー比。自然残留磁化の大きさと現在の場による誘導磁化の大きさの比であり、磁気異常の解釈に関係する。\par

\paragraph{再利用する考え方} 重力・磁気探査の物理量と用語を、測定量・推定対象・基準面・近似条件に分けて理解する。名称だけで探査法を対応させない。

\paragraph{共通理論} \hyperref[theory:gravity:corrections]{第\ref*{theory:gravity:corrections}節}, \hyperref[theory:gravity:geoid]{第\ref*{theory:gravity:geoid}節}, \hyperref[theory:geomagnetism:magnetization]{第\ref*{theory:geomagnetism:magnetization}節}, \hyperref[theory:exploration:observables]{第\ref*{theory:exploration:observables}節}

\par 参考：\cite{gsi_gravity_definitions,kanaya_okuma2011}。

\paragraph{関連問題} \hyperref[q:h30-reg-applied-32]{H30-REG-APPLIED-32}, \hyperref[q:r06-reg-applied-34]{R06-REG-APPLIED-34}

\par 難易度 3（校準初稿）、優先度未評価。

\subsection{令和6年度 問34}
\label{q:r06-reg-applied-34}

\noindent\textbf{原典：} sources/R06/17\_applied\_science.pdf、PDF物理12頁。

\textbf{公式正答：} 3（single）。

\paragraph{問題の要旨} 重力・磁気探査の基準面、磁化特性、移動観測補正と、別の観測分野の語を区別する。

\Needspace{4\baselineskip}\paragraph{確認した条件・記述（独自要約）}\mbox{}\par

\noindent\textbf{選択方向：} 誤り又は最も不適切・関係の薄いものを1つ\par

\paragraph{問われる核心} 重力・磁気探査の物理量と用語

\paragraph{必要知識} 等ポテンシャル面、残留磁化と誘導磁化、磁気秩序の温度依存、回転座標系。

\paragraph{考え方・正答を選ぶ理由（査読待ち初稿）} ジオイドとエトベス補正は重力、ケーニヒスベルガー比とキュリー温度は磁気に直接関係する。AE（アコースティック・エミッション）は微小破壊等に伴う過渡的な弾性波の放出・計測である。岩石研究にも用いるが、重力や磁化の測定とは異なる物理量を扱う。したがってこの肢を選ぶ。

\Needspace{4\baselineskip}\paragraph{選択肢の検討（査読待ち初稿）}\mbox{}\par

\noindent\textbf{1：適切} ケーニヒスベルガー比。自然残留磁化の大きさと現在の場による誘導磁化の大きさの比であり、磁気異常の解釈に関係する。\par

\noindent\textbf{2：適切} ジオイド。基準となる重力等ポテンシャル面であり、重力場と標高の基準に関係する。\par

\noindent\textbf{3：不適切} AE（アコースティック・エミッション）。微小破壊等に伴う過渡的な弾性波の放出・計測である。岩石研究にも用いるが、重力や磁化の測定とは異なる物理量を扱う。\par

\noindent\textbf{4：適切} キュリー温度。強磁性・フェリ磁性の自発磁化が消失する温度に対応し、磁性鉱物や磁気探査の解釈に関係する。\par

\noindent\textbf{5：適切} エトベス補正。移動観測系の運動に伴う見掛重力の差を補正する。\par

\paragraph{再利用する考え方} 重力・磁気探査の物理量と用語を、測定量・推定対象・基準面・近似条件に分けて理解する。名称だけで探査法を対応させない。

\paragraph{共通理論} \hyperref[theory:gravity:corrections]{第\ref*{theory:gravity:corrections}節}, \hyperref[theory:gravity:geoid]{第\ref*{theory:gravity:geoid}節}, \hyperref[theory:geomagnetism:magnetization]{第\ref*{theory:geomagnetism:magnetization}節}, \hyperref[theory:exploration:observables]{第\ref*{theory:exploration:observables}節}

\par 参考：\cite{gsi_gravity_definitions,kanaya_okuma2011,sato_lei2009_ae}。

\paragraph{関連問題} \hyperref[q:r04-reg-applied-35]{R04-REG-APPLIED-35}

\par 難易度 3（校準初稿）、優先度未評価。

\subsection{令和7年度 問33}
\label{q:r07-reg-applied-33}

\noindent\textbf{原典：} sources/R07/17\_applied\_science.pdf、PDF物理12頁。

\textbf{公式正答：} 5（single）。

\paragraph{問題の要旨} 正規重力と各補正を区別し、密度差によるブーゲー異常の符号を誤る記述を選ぶ。

\Needspace{4\baselineskip}\paragraph{確認した条件・記述（独自要約）}\mbox{}\par

\noindent\textbf{選択方向：} 誤り又は最も不適切・関係の薄いものを1つ\par

\paragraph{問われる核心} 基準密度と重力異常の符号

\paragraph{必要知識} 正規重力、フリーエア・ブーゲー・エトベス補正、密度差の引力積分。

\paragraph{考え方・正答を選ぶ理由（査読待ち初稿）} 測定重力から正規場・高さ・地表質量等の効果を取り除き、地下の密度差を評価する。同じ幾何条件なら重力の差は密度差\(\Delta\rho\)に線形であり、高密度は正の、低密度は負の異常を生む。第5肢はこの符号を逆に記述する。

\Needspace{4\baselineskip}\paragraph{選択肢の検討（査読待ち初稿）}\mbox{}\par

\noindent\textbf{1：適切} 正規重力は地球の質量・自転・形状から導く理論値。基準回転楕円体の外部重力ポテンシャルと自転による遠心ポテンシャルから定める。\par

\noindent\textbf{2：適切} ブーゲー補正は測点と海水準面の間の地殻物質の効果を除く。標高差だけでなく、その間の岩石質量の引力を平板等で近似して除く。\par

\noindent\textbf{3：適切} 船・航空機の移動は遠心効果を変える。地球自転と観測機の相対運動により見掛加速度が変わるためエトベス補正が必要である。\par

\noindent\textbf{4：適切} フリーエア補正は高さの違いの影響を補正する。高さによる重力の鉛直勾配の影響をそろえる操作である。\par

\noindent\textbf{5：不適切} 高密度地域のブーゲー異常は小さく、低密度地域では大きくなる。形状等を同条件にした基準密度との差では、密度の過剰は正、欠損は負の重力異常を生む。第5肢は基本的対応を逆にしている。\par

\paragraph{再利用する考え方} 基準密度と重力異常の符号を、測定量・推定対象・基準面・近似条件に分けて理解する。名称だけで探査法を対応させない。

\paragraph{共通理論} \hyperref[theory:gravity:slab]{第\ref*{theory:gravity:slab}節}, \hyperref[theory:gravity:corrections]{第\ref*{theory:gravity:corrections}節}

\par 参考：\cite{gsi_gravity_definitions}。

\paragraph{関連問題} \hyperref[q:h27-reg-applied-22]{H27-REG-APPLIED-22}

\par 難易度 3（校準初稿）、優先度未評価。

\section{電気・電磁探査}

\subsection{平成27年度 問19}
\label{q:h27-reg-applied-19}

\noindent\textbf{原典：} sources/H27/17\_applied\_science.pdf、PDF物理7頁。

\textbf{公式正答：} 3（single）。

\paragraph{問題の要旨} 放射能・地温・電磁・レーダ・重力探査の推定対象を比べ、CSAMTの物性を誤る記述を選ぶ。

\Needspace{4\baselineskip}\paragraph{確認した条件・記述（独自要約）}\mbox{}\par

\noindent\textbf{選択方向：} 誤り又は最も不適切・関係の薄いものを1つ\par

\paragraph{問われる核心} 電磁応答と比抵抗の対応

\paragraph{必要知識} 比抵抗と導電率、平面波インピーダンス、CSAMT遠方場、ガンマ線・温度・重力の観測。

\paragraph{考え方・正答を選ぶ理由（査読待ち初稿）} CSAMTは人工源の遠方場を用いる電磁探査であり、電磁応答から比抵抗構造を推定する。均質・平面波近似なら見掛比抵抗は\(\rho_a=|E/H|^2/(\omega\mu)\)である。磁化率を主対象とする第3肢は物性の対応を誤っている。

\Needspace{4\baselineskip}\paragraph{選択肢の検討（査読待ち初稿）}\mbox{}\par

\noindent\textbf{1：適切} 放射能探査は主にガンマ線強度から断層等の地質情報を得る。天然放射性核種の分布や移動の差を地質情報として利用する。ガンマ線の減衰により観測は主に表層からの寄与であり、深部核種をそのまま透視する意味ではない。\par

\noindent\textbf{2：適切} 地温探査は温泉・地熱・地下水等の調査に用いる。温度分布と熱の移流・伝導を手掛かりに熱水・地下水等を探る。\par

\noindent\textbf{3：不適切} CSAMTは遠方人工源の直交水平電場・磁場から磁化率を得る。主要な推定対象は比抵抗である。水平電場と磁場の比を磁化率そのものと取り違えている。\par

\noindent\textbf{4：適切} レーダは誘電率や導電率の違う境界の電磁反射を捉える。電気特性の差が波の伝播と反射を変えることを利用する。\par

\noindent\textbf{5：適切} 重力は密度構造を求め、標高精度が重要である。標高1 mの誤差は標準鉛直勾配で約0.3086 mGalに相当し、微小重力異常の解釈を妨げる。\par

\paragraph{再利用する考え方} 電磁応答と比抵抗の対応を、測定量・推定対象・基準面・近似条件に分けて理解する。名称だけで探査法を対応させない。

\paragraph{共通理論} \hyperref[theory:exploration:electromagnetic]{第\ref*{theory:exploration:electromagnetic}節}, \hyperref[theory:exploration:observables]{第\ref*{theory:exploration:observables}節}

\par 参考：\cite{ochi_csamt_saijo,gsi_gravity_definitions}。

\par 難易度 3（校準初稿）、優先度未評価。

\section{孔隙水導電とArchie式の適用条件}

\subsection{平成27年度 問18}
\label{q:h27-reg-applied-18}

\noindent\textbf{原典：} sources/H27/17\_applied\_science.pdf、PDF物理6頁。

\textbf{公式正答：} 4（single）。

\par\textbf{検討未確定：} 粘土と砂、割れ目と健岩の比較は含水・塩分・鉱物条件に依存する。風化の全称則とはしない。

\paragraph{問題の要旨} 粘土・不飽和・割れ目・海水が比抵抗へ与える傾向を比較する。

\Needspace{4\baselineskip}\paragraph{確認した条件・記述（独自要約）}\mbox{}\par

\noindent\textbf{原典照合：} 29問21原典画像で否定、全五候補、数式・上付き単位・観測対象を照合。\par

\noindent\textbf{前提：} 孔隙水のイオン伝導と粘土表面伝導、含水した割れ目を考える。\par

\paragraph{問われる核心} 粘土・不飽和・割れ目・海水が比抵抗へ与える傾向を比較する。

\paragraph{必要知識} 孔隙水のイオン伝導と粘土表面伝導、含水した割れ目を考える。

\paragraph{考え方・正答を選ぶ理由（査読待ち初稿）} 4が不適切である。海水のイオン濃度は高く導電率を増すため、孔内・地層水へ海水が入れば通常は抵抗率が低下する。孔内流体の導電率は電極検層の測定にも影響する。

\Needspace{4\baselineskip}\paragraph{選択肢の検討（査読待ち初稿）}\mbox{}\par

\noindent\textbf{1：条件付きで適切} 粘性土は砂質土より一般に低比抵抗。粘土表面伝導と含水が低抵抗へ寄与する。同じ流体・飽和等の比較が必要である。\par

\noindent\textbf{2：条件付きで適切} 不飽和部分は一般に高比抵抗。水を通る連続導電経路が減り通常は抵抗率が上がる。鉱物・流体条件も影響する。\par

\noindent\textbf{3：条件付きで適切} 割れ目と風化が進んだ岩盤は低比抵抗。含水割れ目や粘土化がある場合の傾向である。乾燥空隙や珪化を同じ低下としない。\par

\noindent\textbf{4：不適切} 海水混入で比抵抗が著しく高くなる。イオンを含む海水は通常導電性が高く、淡水状態より抵抗率を下げる方向である。\par

\noindent\textbf{5：適切} 電気検層は裸孔が基本。通常の電極法では地層への電流経路を確保しケーシングによる影響を避ける。\par

\paragraph{再利用する考え方} 低比抵抗を岩石名だけに対応させず、水・塩分・粘土を分ける。

\paragraph{共通理論} \hyperref[theory:exploration:archie]{第\ref*{theory:exploration:archie}節}, \hyperref[theory:exploration:borehole]{第\ref*{theory:exploration:borehole}節}

\par 参考：\cite{epa_resistivity,epa_borehole_resistivity}。

\par 難易度 3（校準初稿）、優先度未評価。

\subsection{平成29年度 問32}
\label{q:h29-reg-applied-32}

\noindent\textbf{原典：} sources/H29/17\_applied\_science.pdf、PDF物理11頁。

\textbf{公式正答：} 4（single）。

\par\textbf{検討未確定：} 高比抵抗水でも絶縁清浄骨格でArchieは成立し得る。礫層・割れ目・地熱SPの符号の全称は留保する。

\paragraph{問題の要旨} Archieの孔隙水伝導条件と地下水・SP解釈を比較する。

\Needspace{4\baselineskip}\paragraph{確認した条件・記述（独自要約）}\mbox{}\par

\noindent\textbf{原典照合：} 29問21原典画像で否定、全五候補、数式・上付き単位・観測対象を照合。\par

\noindent\textbf{前提：} 孔隙水だけで電流が流れ、骨格と表面伝導を無視できる場合にArchieを使う。\par

\paragraph{問われる核心} Archieの孔隙水伝導条件と地下水・SP解釈を比較する。

\paragraph{必要知識} 孔隙水だけで電流が流れ、骨格と表面伝導を無視できる場合にArchieを使う。

\paragraph{考え方・正答を選ぶ理由（査読待ち初稿）} 4の高比抵抗孔隙水なら精度よく適用という保証が不適切である。孔隙水導電が小さいと相対的に表面伝導等が重要となり、仮定が破れやすい。清浄な絶縁骨格なら淡水でも成立し得るため、高比抵抗では絶対に不可とはしない。

\Needspace{4\baselineskip}\paragraph{選択肢の検討（査読待ち初稿）}\mbox{}\par

\noindent\textbf{1：適切} 囲む電極の比抵抗tomographyで二次元・三次元分布。多数配置電位差を逆解析し、覆いとモデル制約の範囲で分布を求める。\par

\noindent\textbf{2：適切} CSAMT近接場は小間隔・高比抵抗で生じやすい。同一周波数でδ∝√ρとなるので、距離/δが小さくなる方向である。\par

\noindent\textbf{3：条件付きで適切} 平野礫帯水層は高抵抗、山地裂帯変質帯は低抵抗。粘土層との比較や含水割れ目・変質粘土の対比として典型的である。塩分、空隙率、乾燥、母岩で変わる。\par

\noindent\textbf{4：不適切} 孔隙水比抵抗が高い場合にArchieを精度よく適用。孔隙水以外の導電を無視できることが必要で、高抵抗水だけで精度は保証できない。表面伝導の相対的重要性が増す。\par

\noindent\textbf{5：条件付きで適切} 鉱体頂部は負異常、地熱対流は正異常。典型鉱化機構と上昇流の特定条件の傾向である。熱電・化学・流動電位と基準点に依存し、全地熱地点の正符号を保証しない。\par

\paragraph{再利用する考え方} 実験式の対象と地質解釈の典型を、全材料に通じる十分条件へ拡張しない。

\paragraph{共通理論} \hyperref[theory:exploration:archie]{第\ref*{theory:exploration:archie}節}, \hyperref[theory:exploration:nearfield]{第\ref*{theory:exploration:nearfield}節}, \hyperref[theory:exploration:ipsp]{第\ref*{theory:exploration:ipsp}節}

\par 参考：\cite{epa_borehole_resistivity,epa_resistivity,epa_sp,usgs_csamt2012}。

\paragraph{関連問題} \hyperref[q:h26-reg-applied-22]{H26-REG-APPLIED-22}

\par 難易度 3（校準初稿）、優先度未評価。

\section{微動アレイの長短波長限界}

\subsection{令和4年度 問34}
\label{q:r04-reg-applied-34}

\noindent\textbf{原典：} sources/R04/17\_applied\_science.pdf、PDF物理12頁。

\textbf{公式正答：} 1（single）。

\par\textbf{検討未確定：} F-K対SPACの大小は解析・入射条件と実装に依存する。水平多層モデルは通常の一次元逆解析の前提であり、手法全般の絶対定義とはしない。

\paragraph{問題の要旨} 微動アレイの短波長限界、水平層モデル、F-K/SPAC、S波逆解析を比較する。

\Needspace{4\baselineskip}\paragraph{確認した条件・記述（独自要約）}\mbox{}\par

\noindent\textbf{原典照合：} 五候補の順序、最適切と不適切の方向、波種、深度目安、曲線名、最長と最短を原典画像で照合。\par

\noindent\textbf{前提：} 短波長側の空間aliasingと長波長側の開口・位相差・統計精度を分ける。\par

\paragraph{問われる核心} 微動アレイの短波長限界、水平層モデル、F-K/SPAC、S波逆解析を比較する。

\paragraph{必要知識} 短波長側の空間aliasingと長波長側の開口・位相差・統計精度を分ける。

\paragraph{考え方・正答を選ぶ理由（査読待ち初稿）} 1が不適切である。空間aliasingは主に短波長・高波数側の限界であり最長波長ではない。長波長側はアレイ内の位相差が小さくなることや有限開口・解析手法の精度等に制約される。

\Needspace{4\baselineskip}\paragraph{選択肢の検討（査読待ち初稿）}\mbox{}\par

\noindent\textbf{1：不適切} 最長波長がアレイ固有の空間aliasingで決まる。最長と最短の役割が逆である。短波長を粗い測点で標本化すると空間波数が折り返して区別できなくなる。\par

\noindent\textbf{2：条件付きで適切} 水平多層構造を仮定して表面波から地下構造を推定する。通常の一次元分散逆解析のモデルである。横方向変化が強い場所や他の多次元算法まで常に水平層であるとはいえない。\par

\noindent\textbf{3：適切} 人工地表加振の表面波で一次元地下モデルを推定する。通常の浅層能動表面波探査の説明であり、周波数・波長ごとの速度を地盤の深度構造に結ぶ。\par

\noindent\textbf{4：条件付きで適切} 同じアレイサイズでF-Kの最大波長はSPACより短い。通常の解析性能の比較であり、実験論文の条件ではSPACがより長波長まで推定した。配置・波の入射・ノイズ・推定算法に依存するため全ての実装に無条件の大小関係を課さない。\par

\noindent\textbf{5：適切} 観測位相速度に近い理論値を持つS波速度構造を作る。順モデルによる分散曲線と観測を照合し主にS波速度を更新する。P波速度・密度・層厚の前提も必要であり唯一解の保証ではない。\par

\paragraph{再利用する考え方} 周波数ではなく波数の標本化として空間aliasingを理解する。

\paragraph{共通理論} \hyperref[theory:exploration:surfacewaves]{第\ref*{theory:exploration:surfacewaves}節}, \hyperref[theory:exploration:arraylimits]{第\ref*{theory:exploration:arraylimits}節}

\par 参考：\cite{epa_surfacewaves,xia_rayleigh1999,japansea_spac2017,array_limits2002}。

\paragraph{関連問題} \hyperref[q:r06-reg-applied-33]{R06-REG-APPLIED-33}

\par 難易度 3（校準初稿）、優先度未評価。

\subsection{令和6年度 問33}
\label{q:r06-reg-applied-33}

\noindent\textbf{原典：} sources/R06/17\_applied\_science.pdf、PDF物理11頁。

\textbf{公式正答：} 2（single）。

\par\textbf{検討未確定：} P波速度・密度・層厚の影響もゼロでない。F-K/SPACの最大波長比較と水平層モデルは通常の解析条件の範囲で読む。

\paragraph{問題の要旨} F-K/SPAC、S波逆解析、水平層、能動起振、最短波長を比較する。

\Needspace{4\baselineskip}\paragraph{確認した条件・記述（独自要約）}\mbox{}\par

\noindent\textbf{原典照合：} 五候補の順序、最適切と不適切の方向、波種、深度目安、曲線名、最長と最短を原典画像で照合。\par

\noindent\textbf{前提：} 微動のRayleigh位相速度は複数の物性へ依存するが、主にS波速度構造に強い感度を持つ。\par

\paragraph{問われる核心} F-K/SPAC、S波逆解析、水平層、能動起振、最短波長を比較する。

\paragraph{必要知識} 微動のRayleigh位相速度は複数の物性へ依存するが、主にS波速度構造に強い感度を持つ。

\paragraph{考え方・正答を選ぶ理由（査読待ち初稿）} 2が不適切である。通常の微動分散逆解析で作るのは主にS波速度モデルであり、P波モデルへ置換する説明が誤りである。P波速度が位相速度へ影響しないという意味ではない。

\Needspace{4\baselineskip}\paragraph{選択肢の検討（査読待ち初稿）}\mbox{}\par

\noindent\textbf{1：条件付きで適切} 同じアレイサイズでF-Kの最大波長はSPACより短い。通常の解析性能の比較であり、実験論文の条件ではSPACがより長波長まで推定した。配置・波の入射・ノイズ・推定算法に依存するため全ての実装に無条件の大小関係を課さない。\par

\noindent\textbf{2：不適切} 観測位相速度に近いP波速度構造モデルを作る。対象とする主な未知量をP波へ取り違える。分散はP波にも依存するが、通常はS波構造を推定し他の物性を別条件で拘束する。\par

\noindent\textbf{3：条件付きで適切} 水平多層構造を仮定して表面波から地下構造を推定する。通常の一次元分散逆解析のモデルである。横方向変化が強い場所や他の多次元算法まで常に水平層であるとはいえない。\par

\noindent\textbf{4：適切} 人工地表加振の表面波で一次元地下モデルを推定する。通常の浅層能動表面波探査の説明であり、周波数・波長ごとの速度を地盤の深度構造に結ぶ。\par

\noindent\textbf{5：適切} 推定位相速度の最短波長は空間aliasingに制約される。波長に対し測点が粗いと異なる空間波数が区別できなくなる。長波長側の位相差の小ささ・有限開口の制約とは区別する。\par

\paragraph{再利用する考え方} 位相速度の感度が強い物性と、影響がゼロではない他物性を区別する。

\paragraph{共通理論} \hyperref[theory:exploration:surfacewaves]{第\ref*{theory:exploration:surfacewaves}節}, \hyperref[theory:exploration:arraylimits]{第\ref*{theory:exploration:arraylimits}節}

\par 参考：\cite{epa_surfacewaves,xia_rayleigh1999,japansea_spac2017,array_limits2002}。

\paragraph{関連問題} \hyperref[q:r04-reg-applied-34]{R04-REG-APPLIED-34}, \hyperref[q:r07-reg-applied-35]{R07-REG-APPLIED-35}

\par 難易度 3（校準初稿）、優先度未評価。

\section{孔内検層の観測量・水位と接触条件}

\subsection{平成23年度 問31}
\label{q:h23-reg-applied-31}

\noindent\textbf{原典：} sources/H23/17\_applied\_science.pdf、PDF物理10頁。

\textbf{公式正答：} 2（single）。

\par\textbf{検討未確定：} VSPの振源又は受振器は通常型と逆配置を含む。裸孔条件は電極検層の基本条件であり全検層へ拡張しない。

\paragraph{問題の要旨} 孔内速度・電気・密度検層の水位条件を比較する。

\Needspace{4\baselineskip}\paragraph{確認した条件・記述（独自要約）}\mbox{}\par

\noindent\textbf{原典照合：} 29問21原典画像で否定、全五候補、数式・上付き単位・観測対象を照合。\par

\noindent\textbf{前提：} 速度検層は弾性波、電極検層は比抵抗、ガンマ密度検層は校正かさ密度である。\par

\paragraph{問われる核心} 孔内速度・電気・密度検層の水位条件を比較する。

\paragraph{必要知識} 速度検層は弾性波、電極検層は比抵抗、ガンマ密度検層は校正かさ密度である。

\paragraph{考え方・正答を選ぶ理由（査読待ち初稿）} 2が誤りである。孔壁に密着させるダウンホール受振器は孔内水位より上でも測定できる。水による振動結合を使う他方式や、導電流体を要する電極検層と区別する。

\Needspace{4\baselineskip}\paragraph{選択肢の検討（査読待ち初稿）}\mbox{}\par

\noindent\textbf{1：条件付きで適切} VSPで孔内に振源又は受振器を設置し反射波を測る。通常の地表振源・孔内受振と逆VSPなどの構成を区別する。孔内深度と反射記録を結び地表反射探査の解釈を助ける。\par

\noindent\textbf{2：不適切} ダウンホールは孔内水位以深だけに適用可能。孔壁へ固定・密着した受振器なら水位のない部分でも測定できる。孔内水位以深だけという限定が誤りである。\par

\noindent\textbf{3：適切} サスペンションで孔内振源と複数受振器を用いる。孔内ゾンデの起振と受振点間の走時差により、孔軸に沿う地盤の区間速度を求める。水中のS波伝播を直接測るという意味ではない。\par

\noindent\textbf{4：条件付きで適切} 電極検層は孔内と地表の電極で比抵抗を測り裸孔を基本とする。通常のnormal電極検層では導電孔内流体を介して地層へ電流を流す。金属ケーシングの影響を避ける。誘導検層等の別方式へこの条件をそのまま一般化しない。\par

\noindent\textbf{5：適切} ガンマ線散乱で孔壁周辺の密度を測り孔隙率を推定する。主にCompton散乱に関わる電子密度を校正してかさ密度を求める。粒子と流体の密度が既知なら孔隙率を推定できる。\par

\paragraph{再利用する考え方} 流体、孔壁結合、線種、センサ位置を観測量ごとに整理する。

\paragraph{共通理論} \hyperref[theory:exploration:borehole]{第\ref*{theory:exploration:borehole}節}, \hyperref[theory:exploration:gammadensity]{第\ref*{theory:exploration:gammadensity}節}

\par 参考：\cite{epa_borehole_resistivity,epa_gamma_density,geebec_ps,oyo_suspension}。

\paragraph{関連問題} \hyperref[q:r01-reg-applied-33]{R01-REG-APPLIED-33}

\par 難易度 3（校準初稿）、優先度未評価。

\subsection{平成24年度 問22}
\label{q:h24-reg-applied-22}

\noindent\textbf{原典：} sources/H24/17\_applied\_science.pdf、PDF物理7頁。

\textbf{公式正答：} 2（single）。

\paragraph{問題の要旨} 孔内弾性波速度測定とnormal電気検層を区別する。

\Needspace{4\baselineskip}\paragraph{確認した条件・記述（独自要約）}\mbox{}\par

\noindent\textbf{原典照合：} 29問21原典画像で否定、全五候補、数式・上付き単位・観測対象を照合。\par

\noindent\textbf{前提：} 振源・受振器の位置と直接観測する物性を対応させる。\par

\paragraph{問われる核心} 孔内弾性波速度測定とnormal電気検層を区別する。

\paragraph{必要知識} 振源・受振器の位置と直接観測する物性を対応させる。

\paragraph{考え方・正答を選ぶ理由（査読待ち初稿）} 2のノルマル法は電極配置による比抵抗検層である。downhole、uphole、VSP、suspensionは弾性波走時又は波形を測り速度構造に関連する。

\Needspace{4\baselineskip}\paragraph{選択肢の検討（査読待ち初稿）}\mbox{}\par

\noindent\textbf{1：適切} ダウンホール法。地表起振、孔内受振で到達走時から速度を求める。\par

\noindent\textbf{2：不適切} ノルマル法。電極の間隔を定めて孔周辺の見掛比抵抗を測り、弾性波速度を直接測る方法ではない。\par

\noindent\textbf{3：適切} アップホール法。孔内起振、地表受振で弾性波走時を測る。\par

\noindent\textbf{4：適切} VSP法。孔内受振等の波形と深度を用い、速度と反射構造を対応させる。\par

\noindent\textbf{5：適切} サスペンション法。孔内振源と複数受振器の到達時間差から区間速度を求める。\par

\paragraph{再利用する考え方} 名称だけでなくセンサと観測量を思い出す。

\paragraph{共通理論} \hyperref[theory:exploration:borehole]{第\ref*{theory:exploration:borehole}節}

\par 参考：\cite{epa_borehole_resistivity,geebec_ps,oyo_suspension}。

\par 難易度 2（校準初稿）、優先度未評価。

\subsection{令和元年度 問33}
\label{q:r01-reg-applied-33}

\noindent\textbf{原典：} sources/R01/17\_applied\_science.pdf、PDF物理12頁。

\textbf{公式正答：} 3（single）。

\par\textbf{検討未確定：} 水の有無だけでdownhole可否は決まらないが、全孔内条件の非依存性を保証しない。VSPと裸孔は方式を特定する。

\paragraph{問題の要旨} 密度検層の放射線と粒子密度・かさ密度を区別する。

\Needspace{4\baselineskip}\paragraph{確認した条件・記述（独自要約）}\mbox{}\par

\noindent\textbf{原典照合：} 29問21原典画像で否定、全五候補、数式・上付き単位・観測対象を照合。\par

\noindent\textbf{前提：} 速度検層は弾性波、電極検層は比抵抗、ガンマ密度検層は校正かさ密度である。\par

\paragraph{問われる核心} 密度検層の放射線と粒子密度・かさ密度を区別する。

\paragraph{必要知識} 速度検層は弾性波、電極検層は比抵抗、ガンマ密度検層は校正かさ密度である。

\paragraph{考え方・正答を選ぶ理由（査読待ち初稿）} 3が不適切である。密度検層はガンマ線の主にCompton散乱による電子密度から孔周辺のかさ密度を求める。ベータ線でも粒子密度の直接測定でもない。

\Needspace{4\baselineskip}\paragraph{選択肢の検討（査読待ち初稿）}\mbox{}\par

\noindent\textbf{1：条件付きで適切} VSPで孔内に振源又は受振器を設置し反射波を測る。通常の地表振源・孔内受振と逆VSPなどの構成を区別する。孔内深度と反射記録を結び地表反射探査の解釈を助ける。\par

\noindent\textbf{2：条件付きで適切} ダウンホールは孔内水の有無が問題とならない。孔壁結合が確保できれば水位より上も測定できる。ただし孔壁崩壊や受振器の結合条件は測定品質へ影響する。\par

\noindent\textbf{3：不適切} ベータ線の散乱で地層の粒子密度を測る。線種はガンマ線、主観測は電子密度から校正するかさ密度である。二つの観測対象を取り違える。\par

\noindent\textbf{4：条件付きで適切} 電極検層は孔内と地表の電極で比抵抗を測り裸孔を基本とする。通常のnormal電極検層では導電孔内流体を介して地層へ電流を流す。金属ケーシングの影響を避ける。誘導検層等の別方式へこの条件をそのまま一般化しない。\par

\noindent\textbf{5：適切} サスペンションで孔内振源と複数受振器を用いる。孔内ゾンデの起振と受振点間の走時差により、孔軸に沿う地盤の区間速度を求める。水中のS波伝播を直接測るという意味ではない。\par

\paragraph{再利用する考え方} 流体、孔壁結合、線種、センサ位置を観測量ごとに整理する。

\paragraph{共通理論} \hyperref[theory:exploration:borehole]{第\ref*{theory:exploration:borehole}節}, \hyperref[theory:exploration:gammadensity]{第\ref*{theory:exploration:gammadensity}節}

\par 参考：\cite{epa_borehole_resistivity,epa_gamma_density,geebec_ps,oyo_suspension}。

\paragraph{関連問題} \hyperref[q:h23-reg-applied-31]{H23-REG-APPLIED-31}, \hyperref[q:r05-reg-applied-34]{R05-REG-APPLIED-34}

\par 難易度 3（校準初稿）、優先度未評価。

\subsection{令和5年度 問34}
\label{q:r05-reg-applied-34}

\noindent\textbf{原典：} sources/R05/17\_applied\_science.pdf、PDF物理12頁。

\textbf{公式正答：} 4（single）。

\par\textbf{検討未確定：} VSPに通常型と逆配置があり、電極検層と誘導検層の条件を区別する。ガンマ密度の孔壁・孔内流体補正も保持。

\paragraph{問題の要旨} 電極検層の孔内流体と各検層の信号を対応させる。

\Needspace{4\baselineskip}\paragraph{確認した条件・記述（独自要約）}\mbox{}\par

\noindent\textbf{原典照合：} 29問21原典画像で否定、全五候補、数式・上付き単位・観測対象を照合。\par

\noindent\textbf{前提：} 速度検層は弾性波、電極検層は比抵抗、ガンマ密度検層は校正かさ密度である。\par

\paragraph{問われる核心} 電極検層の孔内流体と各検層の信号を対応させる。

\paragraph{必要知識} 速度検層は弾性波、電極検層は比抵抗、ガンマ密度検層は校正かさ密度である。

\paragraph{考え方・正答を選ぶ理由（査読待ち初稿）} 4が不適切である。通常の電極型電気検層は孔内の水又は導電泥水を通じて地層へ電流を流すため、その有無と電気特性が問題になる。ダウンホールの孔壁結合とは別である。

\Needspace{4\baselineskip}\paragraph{選択肢の検討（査読待ち初稿）}\mbox{}\par

\noindent\textbf{1：適切} ダウンホールで地表起振・孔内受振の初動走時を測る。孔壁に結合した受振器の位置と走時から区間速度を推定する。経路と屈折の補正が必要である。\par

\noindent\textbf{2：適切} サスペンションで孔内振源と複数受振器を用いる。孔内ゾンデの起振と受振点間の走時差により、孔軸に沿う地盤の区間速度を求める。水中のS波伝播を直接測るという意味ではない。\par

\noindent\textbf{3：適切} ガンマ線散乱で孔壁周辺の密度を測り孔隙率を推定する。主にCompton散乱に関わる電子密度を校正してかさ密度を求める。粒子と流体の密度が既知なら孔隙率を推定できる。\par

\noindent\textbf{4：不適切} 電極検層で孔内水の有無は問題とならない。電極から地層へ電流を流す導電接触が必要で、孔内流体は電流経路と測定値に影響する。電磁誘導検層との混同である。\par

\noindent\textbf{5：条件付きで適切} VSPで孔内に振源又は受振器を設置し反射波を測る。通常の地表振源・孔内受振と逆VSPなどの構成を区別する。孔内深度と反射記録を結び地表反射探査の解釈を助ける。\par

\paragraph{再利用する考え方} 流体、孔壁結合、線種、センサ位置を観測量ごとに整理する。

\paragraph{共通理論} \hyperref[theory:exploration:borehole]{第\ref*{theory:exploration:borehole}節}, \hyperref[theory:exploration:gammadensity]{第\ref*{theory:exploration:gammadensity}節}

\par 参考：\cite{epa_borehole_resistivity,epa_gamma_density,geebec_ps,oyo_suspension}。

\paragraph{関連問題} \hyperref[q:r01-reg-applied-33]{R01-REG-APPLIED-33}, \hyperref[q:r07-reg-applied-34]{R07-REG-APPLIED-34}

\par 難易度 3（校準初稿）、優先度未評価。

\subsection{令和7年度 問34}
\label{q:r07-reg-applied-34}

\noindent\textbf{原典：} sources/R07/17\_applied\_science.pdf、PDF物理13頁。

\textbf{公式正答：} 4（single）。

\par\textbf{検討未確定：} 水位条件は通常の電極検層についてである。VSPの逆配置や孔内流体・孔壁接触による密度補正を別に扱う。

\paragraph{問題の要旨} ガンマ密度検層のCompton散乱と誤った線種を比較する。

\Needspace{4\baselineskip}\paragraph{確認した条件・記述（独自要約）}\mbox{}\par

\noindent\textbf{原典照合：} 29問21原典画像で否定、全五候補、数式・上付き単位・観測対象を照合。\par

\noindent\textbf{前提：} 速度検層は弾性波、電極検層は比抵抗、ガンマ密度検層は校正かさ密度である。\par

\paragraph{問われる核心} ガンマ密度検層のCompton散乱と誤った線種を比較する。

\paragraph{必要知識} 速度検層は弾性波、電極検層は比抵抗、ガンマ密度検層は校正かさ密度である。

\paragraph{考え方・正答を選ぶ理由（査読待ち初稿）} 4が不適切である。通常の密度検層はガンマ線のCompton散乱を利用して校正したかさ密度を求める。ベータ線のRayleigh散乱とする原典記述を正しい文章へ改変しない。

\Needspace{4\baselineskip}\paragraph{選択肢の検討（査読待ち初稿）}\mbox{}\par

\noindent\textbf{1：条件付きで適切} 電極検層の測定区間は孔内水位以深。通常の裸孔電極検層では導電水・泥水を介する区間を測る。電極接触を持たない誘導検層などには一般化しない。\par

\noindent\textbf{2：適切} ダウンホールで地表起振・孔内受振の初動走時を測る。孔壁に結合した受振器の位置と走時から区間速度を推定する。経路と屈折の補正が必要である。\par

\noindent\textbf{3：条件付きで適切} VSPで孔内に振源又は受振器を設置し反射波を測る。通常の地表振源・孔内受振と逆VSPなどの構成を区別する。孔内深度と反射記録を結び地表反射探査の解釈を助ける。\par

\noindent\textbf{4：不適切} ベータ線のレイリー散乱で密度分布を測る。ガンマ線のCompton散乱が主要な原理である。電子密度と岩石かさ密度の校正を使う。\par

\noindent\textbf{5：適切} サスペンションで孔内振源と複数受振器を用いる。孔内ゾンデの起振と受振点間の走時差により、孔軸に沿う地盤の区間速度を求める。水中のS波伝播を直接測るという意味ではない。\par

\paragraph{再利用する考え方} 流体、孔壁結合、線種、センサ位置を観測量ごとに整理する。

\paragraph{共通理論} \hyperref[theory:exploration:borehole]{第\ref*{theory:exploration:borehole}節}, \hyperref[theory:exploration:gammadensity]{第\ref*{theory:exploration:gammadensity}節}

\par 参考：\cite{epa_borehole_resistivity,epa_gamma_density,geebec_ps,oyo_suspension}。

\paragraph{関連問題} \hyperref[q:r05-reg-applied-34]{R05-REG-APPLIED-34}

\par 難易度 3（校準初稿）、優先度未評価。

\section{直流比抵抗の電極配置・見掛比抵抗}

\subsection{平成23年度 問30}
\label{q:h23-reg-applied-30}

\noindent\textbf{原典：} sources/H23/17\_applied\_science.pdf、PDF物理10頁。

\textbf{公式正答：} 2（single）。

\par\textbf{検討未確定：} 最も関係ないは相対比較である。IPは比抵抗測定と密接に関連するため完全な非関連とは説明しない。

\paragraph{問題の要旨} 比抵抗法の電極配置・計算法とIP分極測定を比較する。

\Needspace{4\baselineskip}\paragraph{確認した条件・記述（独自要約）}\mbox{}\par

\noindent\textbf{原典照合：} 29問21原典画像で否定、全五候補、数式・上付き単位・観測対象を照合。\par

\noindent\textbf{前提：} 直流比抵抗の電極配列・順計算と、時間又は周波数に依存する分極応答を分ける。\par

\paragraph{問われる核心} 比抵抗法の電極配置・計算法とIP分極測定を比較する。

\paragraph{必要知識} 直流比抵抗の電極配列・順計算と、時間又は周波数に依存する分極応答を分ける。

\paragraph{考え方・正答を選ぶ理由（査読待ち初稿）} 2のIP法が、静的な比抵抗を求める三配列とlinear filter計算に対して相対的に別の観測を加える。IPも比抵抗と電極を共有するため、無関係とまでは言えない。4は層状構造の比抵抗変換から見掛比抵抗を計算する線形フィルタである。

\Needspace{4\baselineskip}\paragraph{選択肢の検討（査読待ち初稿）}\mbox{}\par

\noindent\textbf{1：適切} ウェンナー法。等間隔の四電極配置で電流と電位差から見掛比抵抗を求める。\par

\noindent\textbf{2：条件付きで適切} IP法。電流印加で生じる分極の減衰又は位相を測る。静的比抵抗法とは目的量が異なるが、同じ配置で比抵抗も測るので関係のなさは相対的である。\par

\noindent\textbf{3：適切} シュランベルジャー法。電流電極と短い電位電極対を使う比抵抗垂直探査の代表的配置である。\par

\noindent\textbf{4：適切} リニアフィルタ法。層状比抵抗モデルの順応答計算に用いられる。一般の信号ノイズ処理だけと解釈しない。\par

\noindent\textbf{5：適切} ダイポール・ダイポール法。電流双極子と電位双極子を離した比抵抗配置で、二次元調査等に用いる。\par

\paragraph{再利用する考え方} 測定配置、順計算法、追加の分極観測を同じ種類の名称と思わない。

\paragraph{共通理論} \hyperref[theory:exploration:dcarrays]{第\ref*{theory:exploration:dcarrays}節}, \hyperref[theory:exploration:ipsp]{第\ref*{theory:exploration:ipsp}節}

\par 参考：\cite{epa_resistivity,epa_ip,usgs_linear_filter1981}。

\paragraph{関連問題} \hyperref[q:h24-reg-applied-20]{H24-REG-APPLIED-20}

\par 難易度 3（校準初稿）、優先度未評価。

\subsection{平成24年度 問20}
\label{q:h24-reg-applied-20}

\noindent\textbf{原典：} sources/H24/17\_applied\_science.pdf、PDF物理7頁。

\textbf{公式正答：} 5（single）。

\paragraph{問題の要旨} 電極による比抵抗測定とコイルによる電磁誘導を区別する。

\Needspace{4\baselineskip}\paragraph{確認した条件・記述（独自要約）}\mbox{}\par

\noindent\textbf{原典照合：} 29問21原典画像で否定、全五候補、数式・上付き単位・観測対象を照合。\par

\noindent\textbf{前提：} 電気的接触を持つ電極測定と送受信ループの誘導測定を区別する。\par

\paragraph{問われる核心} 電極による比抵抗測定とコイルによる電磁誘導を区別する。

\paragraph{必要知識} 電気的接触を持つ電極測定と送受信ループの誘導測定を区別する。

\paragraph{考え方・正答を選ぶ理由（査読待ち初稿）} 5のループ・ループ法が電磁誘導による手法である。残りは電極の接触、配列、見掛比抵抗という電流注入型比抵抗測定に直接関係する。

\Needspace{4\baselineskip}\paragraph{選択肢の検討（査読待ち初稿）}\mbox{}\par

\noindent\textbf{1：適切} 接地抵抗。地盤へ電流を流す電極の接触抵抗が送信と測定条件へ影響する。\par

\noindent\textbf{2：適切} ウェンナー法。等間隔四電極の比抵抗配置である。\par

\noindent\textbf{3：適切} 見掛け比抵抗。測ったV/Iへ配置の幾何係数を掛けた等価半空間の抵抗率である。\par

\noindent\textbf{4：適切} 電極。電流印加と電位差測定に用いる。\par

\noindent\textbf{5：不適切} ループ・ループ法。送受信コイルから地盤の誘導電流応答を測る手法で、電極を使う直流比抵抗測定の用語ではない。比抵抗を推定する目的自体は共通する。\par

\paragraph{再利用する考え方} 推定する物性が同じでも励起・受信の仕組みが異なる。

\paragraph{共通理論} \hyperref[theory:exploration:dcarrays]{第\ref*{theory:exploration:dcarrays}節}, \hyperref[theory:exploration:fdtd]{第\ref*{theory:exploration:fdtd}節}

\par 参考：\cite{epa_resistivity,epa_fdem}。

\paragraph{関連問題} \hyperref[q:h23-reg-applied-30]{H23-REG-APPLIED-30}

\par 難易度 2（校準初稿）、優先度未評価。

\subsection{平成26年度 問23}
\label{q:h26-reg-applied-23}

\noindent\textbf{原典：} sources/H26/17\_applied\_science.pdf、PDF物理10頁。

\textbf{公式正答：} 4（single）。

\par\textbf{検討未確定：} 低比抵抗層の誘導検層精度と迅速広域性は装置条件による。Ω・m\textsuperscript{3}の原典上付き3を画像照合。

\paragraph{問題の要旨} 電気・電磁法の比抵抗単位と測定配置を比較する。

\Needspace{4\baselineskip}\paragraph{確認した条件・記述（独自要約）}\mbox{}\par

\noindent\textbf{原典照合：} 29問21原典画像で否定、全五候補、数式・上付き単位・観測対象を照合。\par

\noindent\textbf{前提：} 抵抗率ρ=RA/LなのでΩmである。電極電流と誘導電流を区別する。\par

\paragraph{問われる核心} 電気・電磁法の比抵抗単位と測定配置を比較する。

\paragraph{必要知識} 抵抗率ρ=RA/LなのでΩmである。電極電流と誘導電流を区別する。

\paragraph{考え方・正答を選ぶ理由（査読待ち初稿）} 4が不適切である。R=ρL/Aからρ=RA/Lとなり、Ω×m\textsuperscript{2}/m=Ωmである。原典のΩ・m\textsuperscript{3}は誤った単位である。

\Needspace{4\baselineskip}\paragraph{選択肢の検討（査読待ち初稿）}\mbox{}\par

\noindent\textbf{1：適切} 直流又は低周波交替直流の電場から比抵抗分布を推定。極性切替を含む低周波電流と電位差の測定を地盤モデルへ逆解析する。\par

\noindent\textbf{2：条件付きで適切} 電磁検層はコイル誘導、低比抵抗層に適する。導電率が高い地層は誘導二次場を生みやすい。ただし誘導数・装置校正・孔状態にも依存し全範囲で精度増大とはしない。\par

\noindent\textbf{3：条件付きで適切} 人工又は自然磁場変動から比抵抗分布を効率的に測る。電磁誘導応答は電気物性へ関係する。効率と面積は方式・機器に依存する。\par

\noindent\textbf{4：不適切} 比抵抗の単位はΩ・m\textsuperscript{3}。抵抗率は電流密度と電場の比例係数でΩmである。体積の単位を掛けると次元が違う。\par

\noindent\textbf{5：適切} 領域を囲む電極から二次元・三次元比抵抗tomography。多方向・多数電流配置の電位データを逆解析して比抵抗分布を推定する。覆い方により分解できる構造が変わる。\par

\paragraph{再利用する考え方} 物性単位を試料幾何から導出し、測定名と物性を一対一に混同しない。

\paragraph{共通理論} \hyperref[theory:exploration:dcarrays]{第\ref*{theory:exploration:dcarrays}節}, \hyperref[theory:exploration:fdtd]{第\ref*{theory:exploration:fdtd}節}

\par 参考：\cite{epa_resistivity,epa_fdem,epa_borehole_resistivity}。

\par 難易度 3（校準初稿）、優先度未評価。

\section{探査深度の目安と観測条件}

\subsection{平成27年度 問21}
\label{q:h27-reg-applied-21}

\noindent\textbf{原典：} sources/H27/17\_applied\_science.pdf、PDF物理8頁。

\textbf{公式正答：} 1（single）。

\par\textbf{検討未確定：} 公式1の100mと通常20mの区別は当時の一般的装置条件の読みである。長波長表面波で100mが不可能とはいえない。比抵抗300mと音響基盤も一律保証ではない。

\paragraph{問題の要旨} 表面波の深度目安、屈折ブラインド、反射展開、比抵抗、音響基盤を比較する。

\Needspace{4\baselineskip}\paragraph{確認した条件・記述（独自要約）}\mbox{}\par

\noindent\textbf{原典照合：} 五候補の順序、最適切と不適切の方向、波種、深度目安、曲線名、最長と最短を原典画像で照合。\par

\noindent\textbf{前提：} 能動の浅部表面波探査の通常目安と、長波長の利用により深度が変わる事実を分ける。\par

\paragraph{問われる核心} 表面波の深度目安、屈折ブラインド、反射展開、比抵抗、音響基盤を比較する。

\paragraph{必要知識} 能動の浅部表面波探査の通常目安と、長波長の利用により深度が変わる事実を分ける。

\paragraph{考え方・正答を選ぶ理由（査読待ち初稿）} 公式1は通常の浅層表面波の目安を100m程度とした説明を選ぶ。環境省資料では20m程度、数m〜数十mの通常対象を述べる。一方、長周期起振や微動で100m級以上を調べる方法もあるため、100mが原理的に不可能とは解説しない。

\Needspace{4\baselineskip}\paragraph{選択肢の検討（査読待ち初稿）}\mbox{}\par

\noindent\textbf{1：条件付きで適切} 表面波の分散でS波構造を調べ探査深度は100m程度までである。分散によるS波推定は正しい。通常の浅部能動測定では20m級の目安であり公式は100mを不適切とするが、装置と最大波長を指定しない普遍的上限は定められない。\par

\noindent\textbf{2：適切} 深部へ増速する層構造でも薄い中間層が初動に現れない。臨界屈折波が生じても、直接波や深い層を通る波より早くなる距離範囲がなければ初動に現れない。低速度逆転だけが見えない原因ではない。\par

\noindent\textbf{3：条件付きで適切} 浅層反射で最大オフセットを探査深度程度に設定する。浅層反射の一般的な展開目安として読む。速度解析・近接雑音・反射角・地形・重合数で調整し、全ての計測に必要な厳密等式ではない。\par

\noindent\textbf{4：条件付きで適切} 二次元比抵抗法は低比抵抗部を検出し実用深度が300m程度である。導電性領域へ電流が集まりやすい等の特性があるが、電極展開、コントラスト、ノイズで感度と深度が変わる。300mは原典の典型目安であり装置非依存の保証ではない。\par

\noindent\textbf{5：条件付きで適切} 海上音波で海底下反射を測り固結層等の上面が音響基盤となる。大きなインピーダンス差と減衰等によって深部への透過が弱まり、内部を捉えにくくなる。基盤より下は常に観測不能という絶対命題ではない。\par

\paragraph{再利用する考え方} 探査深度を手法名に固定せず、波長・配列・震源・物性・雑音条件を確認する。

\paragraph{共通理論} \hyperref[theory:exploration:surfacewaves]{第\ref*{theory:exploration:surfacewaves}節}, \hyperref[theory:exploration:reflectiongeometry]{第\ref*{theory:exploration:reflectiongeometry}節}

\par 参考：\cite{moe_surface2004,epa_surfacewaves,segj_integrated2017,epa_seismic_reflection}。

\par 難易度 3（校準初稿）、優先度未評価。

\section{周波数領域と時間領域電磁応答}

\subsection{平成28年度 問19}
\label{q:h28-reg-applied-19}

\noindent\textbf{原典：} sources/H28/17\_applied\_science.pdf、PDF物理7頁。

\textbf{公式正答：} 1（single）。

\paragraph{問題の要旨} 電気・電磁現象と地震反射のNMO補正を区別する。

\Needspace{4\baselineskip}\paragraph{確認した条件・記述（独自要約）}\mbox{}\par

\noindent\textbf{原典照合：} 29問21原典画像で否定、全五候補、数式・上付き単位・観測対象を照合。\par

\noindent\textbf{前提：} NMOは反射走時のoffset依存を補正する地震探査処理である。\par

\paragraph{問われる核心} 電気・電磁現象と地震反射のNMO補正を区別する。

\paragraph{必要知識} NMOは反射走時のoffset依存を補正する地震探査処理である。

\paragraph{考え方・正答を選ぶ理由（査読待ち初稿）} 1のNMO補正である。表皮・近接場・誘導分極・Schumann共振は電気・電磁応答に直接関係する。SchumannはMT電磁記録でも観測される。

\Needspace{4\baselineskip}\paragraph{選択肢の検討（査読待ち初稿）}\mbox{}\par

\noindent\textbf{1：不適切} NMO補正。normal moveoutは震源受振器の距離による反射走時差を補正する。電磁現象の名称ではない。\par

\noindent\textbf{2：適切} 表皮効果。導電媒質内の周期電磁場の減衰と深度感度へ関係する。\par

\noindent\textbf{3：適切} ニアフィールド効果。有限送信源からの近距離場が平面波近似と異なることをいう。\par

\noindent\textbf{4：適切} IP効果。電流による地下界面の分極で減衰・位相が変わる。\par

\noindent\textbf{5：適切} シューマン共振。地球と電離層の電磁共振であり、自然MTの電場磁場記録にも現れる。\par

\paragraph{再利用する考え方} 同じ探査で扱う補正でも信号の物理原理を識別する。

\paragraph{共通理論} \hyperref[theory:exploration:nearfield]{第\ref*{theory:exploration:nearfield}節}, \hyperref[theory:exploration:ipsp]{第\ref*{theory:exploration:ipsp}節}, \hyperref[theory:exploration:reflectiongeometry]{第\ref*{theory:exploration:reflectiongeometry}節}

\par 参考：\cite{schumann_mt2009,epa_ip,usgs_csamt2012}。

\par 難易度 2（校準初稿）、優先度未評価。

\subsection{令和2年度 問34}
\label{q:r02-reg-applied-34}

\noindent\textbf{原典：} sources/R02/17\_applied\_science.pdf、PDF物理12頁。

\textbf{公式正答：} 4（single）。

\par\textbf{検討未確定：} 精度0.01mGalは機器/現場の条件値、検層孔隙率は間接推定、AEの先行は全称ではない。

\paragraph{問題の要旨} 周波数領域の位相と時間領域の過渡応答を区別する。

\Needspace{4\baselineskip}\paragraph{確認した条件・記述（独自要約）}\mbox{}\par

\noindent\textbf{原典照合：} 29問21原典画像で否定、全五候補、数式・上付き単位・観測対象を照合。\par

\noindent\textbf{前提：} 正弦波定常応答は振幅と位相、遮断後応答は時間に対する減衰を測る。\par

\paragraph{問われる核心} 周波数領域の位相と時間領域の過渡応答を区別する。

\paragraph{必要知識} 正弦波定常応答は振幅と位相、遮断後応答は時間に対する減衰を測る。

\paragraph{考え方・正答を選ぶ理由（査読待ち初稿）} 4が不適切である。Slingram等の周波数領域電磁法は連続正弦送信に対する受信振幅と位相差を測る。過渡応答は時間領域の測定である。

\Needspace{4\baselineskip}\paragraph{選択肢の検討（査読待ち初稿）}\mbox{}\par

\noindent\textbf{1：条件付きで適切} 微小重力差の相対測定や鉛直重力勾配を利用。0.01mGal級以下の精度を持つ計器の利用と勾配観測がある。全装置と測定現場に共通の精度保証ではない。\par

\noindent\textbf{2：適切} 領域を囲む信号源センサで物性を逆解析。tomographyの観測覆いと感度に沿って物性分布を推定する。\par

\noindent\textbf{3：条件付きで適切} 検層で比抵抗密度速度孔隙率等を深度に記録。物性を連続記録し、孔隙率のように複数物性と校正から推定する量もある。すべてが直接測定ではない。\par

\noindent\textbf{4：不適切} 周波数領域で連続波の振幅と過渡応答を測る。定常の振幅と位相を測る方式と、遮断後の時間領域過渡を混同する。\par

\noindent\textbf{5：条件付きで適切} AEは微小破壊で主破壊に先行することが多い。微小破壊に伴う弾性波放出が前兆の研究対象になる。主破壊前だけの定義又は常に発生する保証ではない。\par

\paragraph{再利用する考え方} f領域の位相とt領域の減衰、直接物性と校正推定量を分ける。

\paragraph{共通理論} \hyperref[theory:exploration:fdtd]{第\ref*{theory:exploration:fdtd}節}, \hyperref[theory:exploration:multimethod]{第\ref*{theory:exploration:multimethod}節}

\par 参考：\cite{epa_fdem,epa_tdem_archive,sato_lei2009_ae,gsi_gravity_definitions}。

\par 難易度 3（校準初稿）、優先度未評価。

\section{探査信号のフィルタ選択}

\subsection{平成23年度 問28}
\label{q:h23-reg-applied-28}

\noindent\textbf{原典：} sources/H23/17\_applied\_science.pdf、PDF物理9頁。

\textbf{公式正答：} 4（single）。

\par\textbf{検討未確定：} 移動平均の除去は低減であり完全除去ではない。線形性だけで単一応答の畳み込みは保証しない。標本化の原文は一周期と間隔を混ぜた表現であるため、帯域と最高周波数を明示して条件付きで読む。

\paragraph{問題の要旨} 空間フィルタ・平均・インパルス応答・低域通過・標本化を区別する。

\Needspace{4\baselineskip}\paragraph{確認した条件・記述（独自要約）}\mbox{}\par

\noindent\textbf{原典照合：} 否定の対象、指数、周波数と周期、五つの候補の列順を原典画像で照合。\par

\noindent\textbf{前提：} 通常の二次元画像と時系列を対象とする。低域は低周波数すなわち長周期である。\par

\paragraph{問われる核心} 空間フィルタ・平均・インパルス応答・低域通過・標本化を区別する。

\paragraph{必要知識} 通常の二次元画像と時系列を対象とする。低域は低周波数すなわち長周期である。

\paragraph{考え方・正答を選ぶ理由（査読待ち初稿）} 4は低域通過の役割を逆にしている。長周期成分を残し、短周期成分を減衰させるのが低域通過である。他の説明は典型的処理として読むが、完全除去と畳み込みの省略条件を残す。

\Needspace{4\baselineskip}\paragraph{選択肢の検討（査読待ち初稿）}\mbox{}\par

\noindent\textbf{1：適切} 空間フィルタを二次元画像のエッジ・平滑化・差分に用いる。二次元データに対し空間方向の重み付けや差分を行う。一次元・三次元にも拡張でき、二次元しか扱えないという意味ではない。\par

\noindent\textbf{2：条件付きで適切} 注目点付近の移動平均でランダム成分を除去する。独立で平均ゼロの雑音をM点平均すれば分散は1/Mに減る。有限Mで雑音が厳密にゼロになる保証はない。中央窓以外に因果的な後方窓もある。\par

\noindent\textbf{3：条件付きで適切} 線形システムのインパルス応答を重ねて任意入力への応答を求める。線形性により各時刻のインパルスへの応答を重ねることはできる。単一の応答を時間移動する畳み込み表現には時不変性も必要であり、時変系には二時刻の応答を用いる。\par

\noindent\textbf{4：不適切} 低域通過は短周期を残し長周期を減衰させる。周期Tと周波数fは逆数の関係である。低周波数は長周期なので、記述の通過側と減衰側は逆である。\par

\noindent\textbf{5：条件付きで適切} 一周期二標本未満では元データを再構成できない。一般の帯域制限信号を一意に再構成するにはスペクトルの重なりを避ける標本化が必要である。特別な事前情報を持つ信号まで一律に不可能とする定理ではなく、最高周波数と標本化間隔を明示して読む。\par

\paragraph{再利用する考え方} フィルタの名称を周波数で定義してから周期へ変換する。

\paragraph{共通理論} \hyperref[theory:exploration:filters]{第\ref*{theory:exploration:filters}節}, \hyperref[theory:exploration:lti]{第\ref*{theory:exploration:lti}節}, \hyperref[theory:exploration:sampling]{第\ref*{theory:exploration:sampling}節}

\par 参考：\cite{mit_signals2010,mit_moving2012,scipy_notch,scipy_fft,usgs_processing2000}。

\paragraph{関連問題} \hyperref[q:h28-reg-applied-20]{H28-REG-APPLIED-20}

\par 難易度 3（校準初稿）、優先度未評価。

\subsection{平成24年度 問19}
\label{q:h24-reg-applied-19}

\noindent\textbf{原典：} sources/H24/17\_applied\_science.pdf、PDF物理7頁。

\textbf{公式正答：} 5（single）。

\par\textbf{検討未確定：} 移動平均でランダム成分を除去するとは統計的な低減の意味である。信号まで平均化で弱める場合と、相関雑音が残る場合を区別する。

\paragraph{問題の要旨} 相互相関・重合・移動平均・FFT・ノッチの目的を対応させる。

\Needspace{4\baselineskip}\paragraph{確認した条件・記述（独自要約）}\mbox{}\par

\noindent\textbf{原典照合：} 否定の対象、指数、周波数と周期、五つの候補の列順を原典画像で照合。\par

\noindent\textbf{前提：} 二つの時系列の相関、複数記録の平均化、DFT、狭帯域阻止を区別する。\par

\paragraph{問われる核心} 相互相関・重合・移動平均・FFT・ノッチの目的を対応させる。

\paragraph{必要知識} 二つの時系列の相関、複数記録の平均化、DFT、狭帯域阻止を区別する。

\paragraph{考え方・正答を選ぶ理由（査読待ち初稿）} 5が不適切である。ノッチは特定周波数の周囲を阻止するものであり、その周波数だけを通過させる説明は逆である。ランダム雑音の抑制は統計的効果であり有限平均での完全消失ではない。

\Needspace{4\baselineskip}\paragraph{選択肢の検討（査読待ち初稿）}\mbox{}\par

\noindent\textbf{1：適切} 相互相関から二波形の類似と時間遅れを検出する。一方の波形をずらして積の和又は平均を求め、相関の極大から遅れを推定する。極大が常に唯一とは限らず、周期性・雑音に注意する。\par

\noindent\textbf{2：条件付きで適切} 複数記録の重合でランダム雑音を抑える。信号の位相をそろえ、雑音が独立又は無相関なら振幅S/Nは回数の平方根だけ改善する。相関雑音や未補正の走時差がある場合は同じ効果を保証しない。\par

\noindent\textbf{3：条件付きで適切} 注目点付近の移動平均でランダム成分を除去する。独立で平均ゼロの雑音をM点平均すれば分散は1/Mに減る。有限Mで雑音が厳密にゼロになる保証はない。中央窓以外に因果的な後方窓もある。\par

\noindent\textbf{4：適切} FFTで離散時系列のフーリエ係数を高速に計算する。FFTはDFTと異なる変換ではなく、DFTを効率よく計算する算法である。標本化によるaliasingを自動的に取り除く算法ではない。\par

\noindent\textbf{5：不適切} ノッチで指定周波数成分だけを通過させる。ノッチは狭帯域阻止である。通過と遮断を逆にしており、帯域通過の説明と取り違えている。\par

\paragraph{再利用する考え方} 名称ではなく保持する成分と抑える成分を書く。

\paragraph{共通理論} \hyperref[theory:exploration:filters]{第\ref*{theory:exploration:filters}節}, \hyperref[theory:exploration:correlation]{第\ref*{theory:exploration:correlation}節}, \hyperref[theory:exploration:stack]{第\ref*{theory:exploration:stack}節}, \hyperref[theory:exploration:fft]{第\ref*{theory:exploration:fft}節}

\par 参考：\cite{mit_signals2010,mit_moving2012,scipy_notch,scipy_fft,usgs_processing2000}。

\paragraph{関連問題} \hyperref[q:h29-reg-applied-35]{H29-REG-APPLIED-35}

\par 難易度 3（校準初稿）、優先度未評価。

\subsection{平成29年度 問35}
\label{q:h29-reg-applied-35}

\noindent\textbf{原典：} sources/H29/17\_applied\_science.pdf、PDF物理12頁。

\textbf{公式正答：} 5（single）。

\par\textbf{検討未確定：} 移動平均でランダム成分を除去するとは統計的な低減の意味である。信号まで平均化で弱める場合と、相関雑音が残る場合を区別する。

\paragraph{問題の要旨} 相互相関・重合・移動平均・FFT・ノッチの目的を対応させる。

\Needspace{4\baselineskip}\paragraph{確認した条件・記述（独自要約）}\mbox{}\par

\noindent\textbf{原典照合：} 否定の対象、指数、周波数と周期、五つの候補の列順を原典画像で照合。\par

\noindent\textbf{前提：} 二つの時系列の相関、複数記録の平均化、DFT、狭帯域阻止を区別する。\par

\paragraph{問われる核心} 相互相関・重合・移動平均・FFT・ノッチの目的を対応させる。

\paragraph{必要知識} 二つの時系列の相関、複数記録の平均化、DFT、狭帯域阻止を区別する。

\paragraph{考え方・正答を選ぶ理由（査読待ち初稿）} 5が不適切である。ノッチは特定周波数の周囲を阻止するものであり、その周波数だけを通過させる説明は逆である。ランダム雑音の抑制は統計的効果であり有限平均での完全消失ではない。

\Needspace{4\baselineskip}\paragraph{選択肢の検討（査読待ち初稿）}\mbox{}\par

\noindent\textbf{1：適切} 相互相関から二波形の類似と時間遅れを検出する。一方の波形をずらして積の和又は平均を求め、相関の極大から遅れを推定する。極大が常に唯一とは限らず、周期性・雑音に注意する。\par

\noindent\textbf{2：条件付きで適切} 複数記録の重合でランダム雑音を抑える。信号の位相をそろえ、雑音が独立又は無相関なら振幅S/Nは回数の平方根だけ改善する。相関雑音や未補正の走時差がある場合は同じ効果を保証しない。\par

\noindent\textbf{3：条件付きで適切} 注目点付近の移動平均でランダム成分を除去する。独立で平均ゼロの雑音をM点平均すれば分散は1/Mに減る。有限Mで雑音が厳密にゼロになる保証はない。中央窓以外に因果的な後方窓もある。\par

\noindent\textbf{4：適切} FFTで離散時系列のフーリエ係数を高速に計算する。FFTはDFTと異なる変換ではなく、DFTを効率よく計算する算法である。標本化によるaliasingを自動的に取り除く算法ではない。\par

\noindent\textbf{5：不適切} ノッチで指定周波数成分だけを通過させる。ノッチは狭帯域阻止である。通過と遮断を逆にしており、帯域通過の説明と取り違えている。\par

\paragraph{再利用する考え方} 名称ではなく保持する成分と抑える成分を書く。

\paragraph{共通理論} \hyperref[theory:exploration:filters]{第\ref*{theory:exploration:filters}節}, \hyperref[theory:exploration:correlation]{第\ref*{theory:exploration:correlation}節}, \hyperref[theory:exploration:stack]{第\ref*{theory:exploration:stack}節}, \hyperref[theory:exploration:fft]{第\ref*{theory:exploration:fft}節}

\par 参考：\cite{mit_signals2010,mit_moving2012,scipy_notch,scipy_fft,usgs_processing2000}。

\paragraph{関連問題} \hyperref[q:h24-reg-applied-19]{H24-REG-APPLIED-19}

\par 難易度 3（校準初稿）、優先度未評価。

\section{探査の順問題・逆問題と対象雑音}

\subsection{令和7年度 問31}
\label{q:r07-reg-applied-31}

\noindent\textbf{原典：} sources/R07/17\_applied\_science.pdf、PDF物理11頁。

\textbf{公式正答：} 1（single）。

\par\textbf{検討未確定：} 標本化の等号端点と帯域制限、理想再構成の条件が省略される。S/Nの平方根は改善倍率であり、元のS/Nが常に1ではない。

\paragraph{問題の要旨} 順問題と逆問題、対象依存の雑音、標本化と重合を区別する。

\Needspace{4\baselineskip}\paragraph{確認した条件・記述（独自要約）}\mbox{}\par

\noindent\textbf{原典照合：} 否定の対象、指数、周波数と周期、五つの候補の列順を原典画像で照合。\par

\noindent\textbf{前提：} 地下モデルからデータを予測するのが順問題、観測データから地下を推定するのが逆問題である。\par

\paragraph{問われる核心} 順問題と逆問題、対象依存の雑音、標本化と重合を区別する。

\paragraph{必要知識} 地下モデルからデータを予測するのが順問題、観測データから地下を推定するのが逆問題である。

\paragraph{考え方・正答を選ぶ理由（査読待ち初稿）} 1が不適切である。地下構造を求める解析を順問題、計測を逆問題とする対応が逆である。計測そのものはデータ取得であり、順問題とはモデルからデータを計算する写像を指す。

\Needspace{4\baselineskip}\paragraph{選択肢の検討（査読待ち初稿）}\mbox{}\par

\noindent\textbf{1：不適切} 地下推定の解析が順問題、データを得る計測が逆問題である。F(m)でモデルからデータを計算するのが順問題である。与えたデータからmを推定するのが逆問題であり、原典の役割は逆転している。\par

\noindent\textbf{2：適切} 入力信号によって発生した表面波でも解析対象でなければ雑音である。ノイズは対象に対する不要成分の呼び方である。表面波探査では信号となる同じ波が、反射法では相関した妨害波となり得る。\par

\noindent\textbf{3：条件付きで適切} 最高周波数の二倍以上で標本化すれば完全復元できる。帯域制限・理想的標本と再構成を想定した説明である。最高周波数成分を含む場合の等号端点、量子化、有限長、雑音を無視できない。安全な条件はfs>2fmaxであり、境界では位相によって標本が全てゼロとなる例がある。\par

\noindent\textbf{4：適切} ノッチで特定の周波数成分を遮断する。狭い阻止帯域で商用電源などの周期的な混入を抑える。指定帯域だけを通過させるフィルタではない。実装には有限の帯域幅と位相応答がある。\par

\noindent\textbf{5：条件付きで適切} 重合によるランダム雑音に対するS/Nが回数の平方根になる。同一信号の位相と振幅がそろい、雑音が無相関・平均ゼロ・同分散なら成り立つ。一般には元のS/Nに平方根Nを掛ける改善倍率を意味する。\par

\paragraph{再利用する考え方} ノイズをランダムに限らず目的から定義し、写像の向きで順逆を区別する。

\paragraph{共通理論} \hyperref[theory:exploration:inverse]{第\ref*{theory:exploration:inverse}節}, \hyperref[theory:exploration:stack]{第\ref*{theory:exploration:stack}節}, \hyperref[theory:exploration:sampling]{第\ref*{theory:exploration:sampling}節}

\par 参考：\cite{mit_signals2010,mit_moving2012,scipy_notch,scipy_fft,usgs_processing2000,pygimli_inversion}。

\par 難易度 3（校準初稿）、優先度未評価。

\section{四電極と接地抵抗の分離}

\subsection{平成25年度 問34}
\label{q:h25-reg-applied-34}

\noindent\textbf{原典：} sources/H25/17\_applied\_science.pdf、PDF物理12頁。

\textbf{公式正答：} 3（single）。

\paragraph{問題の要旨} 四電極法で接触抵抗の電圧降下を分離する理由を選ぶ。

\Needspace{4\baselineskip}\paragraph{確認した条件・記述（独自要約）}\mbox{}\par

\noindent\textbf{原典照合：} 29問21原典画像で否定、全五候補、数式・上付き単位・観測対象を照合。\par

\noindent\textbf{前提：} 電流注入電極と高入力抵抗の電位測定電極を分ける。電位電極電流は十分小さい。\par

\paragraph{問われる核心} 四電極法で接触抵抗の電圧降下を分離する理由を選ぶ。

\paragraph{必要知識} 電流注入電極と高入力抵抗の電位測定電極を分ける。電位電極電流は十分小さい。

\paragraph{考え方・正答を選ぶ理由（査読待ち初稿）} 3である。二端子では地盤に加えて電流電極の接触抵抗の電圧降下も測る。別の電位電極へ高入力抵抗の電圧計を接続すれば、その電極を流れる電流が小さく、接触抵抗による誤差を軽減する。送信側の接地抵抗は依然として電流確保に影響する。

\Needspace{4\baselineskip}\paragraph{選択肢の検討（査読待ち初稿）}\mbox{}\par

\noindent\textbf{1：不適切} 電磁カップリングを小さくする。電磁誘導の結合はケーブル配置・周波数等にも依存し、四電極化の主要な理由ではない。\par

\noindent\textbf{2：不適切} 測定電圧を大きくする。Vは電極幾何と地盤応答で決まり、電極を分けることだけで増大しない。\par

\noindent\textbf{3：適切} 接地抵抗の影響を軽減する。電流端子の接触電圧降下を電位差観測から分け、電位端子の小電流によりその接触降下も抑える。\par

\noindent\textbf{4：不適切} 地形の影響を無視する。地形は電流分布を変えるため、別電位電極でも順モデルの地形補正が必要である。\par

\noindent\textbf{5：不適切} 探査深度を大きくする。深度感度は配列と間隔等で決まり、電極端子を分ける理由そのものではない。\par

\paragraph{再利用する考え方} 送信電流を確保する問題と、電位測定の接触誤差を減らす問題を分ける。

\paragraph{共通理論} \hyperref[theory:exploration:fourpoint]{第\ref*{theory:exploration:fourpoint}節}

\par 参考：\cite{epa_resistivity}。

\par 難易度 2（校準初稿）、優先度未評価。

\section{GPRの周波数・分解能・探査深度}

\subsection{平成24年度 問21}
\label{q:h24-reg-applied-21}

\noindent\textbf{原典：} sources/H24/17\_applied\_science.pdf、PDF物理7頁。

\textbf{公式正答：} 3（single）。

\par\textbf{検討未確定：} 固定波形の理想鉛直分解能は深度そのものだけで単調低下するわけではない。実地の交換条件と検出限界を区別。

\paragraph{問題の要旨} GPRの深度・周波数・分解能と調査配置を比較する。

\Needspace{4\baselineskip}\paragraph{確認した条件・記述（独自要約）}\mbox{}\par

\noindent\textbf{原典照合：} 29問21原典画像で否定、全五候補、数式・上付き単位・観測対象を照合。\par

\noindent\textbf{前提：} 低損失非磁性媒質でv≈c/√εr、波長λ=v/f、パルス帯域Bに対する距離分解はおよそv/(2B)である。導電率等が到達深度を制限する。\par

\paragraph{問われる核心} GPRの深度・周波数・分解能と調査配置を比較する。

\paragraph{必要知識} 低損失非磁性媒質でv≈c/√εr、波長λ=v/f、パルス帯域Bに対する距離分解はおよそv/(2B)である。導電率等が到達深度を制限する。

\paragraph{考え方・正答を選ぶ理由（査読待ち初稿）} 3が不適切である。深部まで測るための低周波化、波の減衰・拡がりやノイズにより一般に深部の分解・検出が難しくなる。深度が増しても分解能が低下しないという一般保証はできない。

\Needspace{4\baselineskip}\paragraph{選択肢の検討（査読待ち初稿）}\mbox{}\par

\noindent\textbf{1：適切} 地表型と孔内型のレーダがある。地表からの反射と孔内アンテナによる反射・透過などの構成がある。\par

\noindent\textbf{2：適切} 周波数は探査深度と分解能へ関係。波長・帯域と吸収が分解と到達深度を制限する。\par

\noindent\textbf{3：不適切} 一般に深度が増しても分解能は低下しない。深部の低周波化や減衰・ビーム拡がりが分解と検出を難しくする。非低下の一般保証は不適切である。\par

\noindent\textbf{4：適切} 直線埋設物へできる限り直角の測線。横切る測線で位置の変化と点状反射の走時を捉えやすくする。未知方位では複数方位を測る。\par

\noindent\textbf{5：適切} 高分解能なので路面下空洞に利用。浅部の空洞境界の電気特性差と十分な波長・信号品質を使う。\par

\paragraph{再利用する考え方} 反射走時を深度と混同せず、帯域・速度・損失と探査目的を対応させる。

\paragraph{共通理論} \hyperref[theory:exploration:gpr]{第\ref*{theory:exploration:gpr}節}, \hyperref[theory:exploration:gprresolution]{第\ref*{theory:exploration:gprresolution}節}

\par 参考：\cite{epa_gpr}。

\par 難易度 3（校準初稿）、優先度未評価。

\subsection{令和2年度 問33}
\label{q:r02-reg-applied-33}

\noindent\textbf{原典：} sources/R02/17\_applied\_science.pdf、PDF物理12頁。

\textbf{公式正答：} 3（single）。

\par\textbf{検討未確定：} 2の分解能は能力と最小距離で比例方向が変わる。深度の厳密反比例と5の間接相関全面否定は留保。

\paragraph{問題の要旨} GPRのMHz帯と測線・速度・強度の推定限界を照合する。

\Needspace{4\baselineskip}\paragraph{確認した条件・記述（独自要約）}\mbox{}\par

\noindent\textbf{原典照合：} 29問21原典画像で否定、全五候補、数式・上付き単位・観測対象を照合。\par

\noindent\textbf{前提：} 低損失非磁性媒質でv≈c/√εr、波長λ=v/f、パルス帯域Bに対する距離分解はおよそv/(2B)である。導電率等が到達深度を制限する。\par

\paragraph{問われる核心} GPRのMHz帯と測線・速度・強度の推定限界を照合する。

\paragraph{必要知識} 低損失非磁性媒質でv≈c/√εr、波長λ=v/f、パルス帯域Bに対する距離分解はおよそv/(2B)である。導電率等が到達深度を制限する。

\paragraph{考え方・正答を選ぶ理由（査読待ち初稿）} 3の10kHzが不適切である。通常の地中レーダはMHz帯以上の高周波パルスを使い、10kHzを分類境界とする説明は三桁違う。2の比例は定量則ではなく分解能力と深度の傾向として限定する。

\Needspace{4\baselineskip}\paragraph{選択肢の検討（査読待ち初稿）}\mbox{}\par

\noindent\textbf{1：適切} profileにより構造、wide angleで電磁波速度を推定。走査波形を位置へ並べるprofileと、送受信距離を変えて走時から速度を求める構成を使い分ける。\par

\noindent\textbf{2：条件付きで適切} 周波数と分解能は比例、深度は逆比例。周波数が高いほど分解能力が良く、深度を制限する通常の傾向である。最小識別距離は波長に比例し、実深度が厳密に1/fとなる保証はない。\par

\noindent\textbf{3：不適切} 10kHz程度以上をGPR、以下を電磁探査とする。通常の目安は10MHz程度であり、kHzの桁を取り違える。物理的境界としての厳密性もない。\par

\noindent\textbf{4：適切} 直線埋設物へできる限り直角の測線。横切る測線で位置の変化と点状反射の走時を捉えやすくする。未知方位では複数方位を測る。\par

\noindent\textbf{5：条件付きで適切} 電気境界を測るため力学特性との相関はほとんどない。強度や硬さを直接測らないことは正しい。ただし含水・空隙・劣化など共通要因を通じた間接相関まで全面否定はできない。\par

\paragraph{再利用する考え方} 反射走時を深度と混同せず、帯域・速度・損失と探査目的を対応させる。

\paragraph{共通理論} \hyperref[theory:exploration:gpr]{第\ref*{theory:exploration:gpr}節}, \hyperref[theory:exploration:gprresolution]{第\ref*{theory:exploration:gprresolution}節}

\par 参考：\cite{epa_gpr}。

\paragraph{関連問題} \hyperref[q:r05-reg-applied-33]{R05-REG-APPLIED-33}

\par 難易度 3（校準初稿）、優先度未評価。

\subsection{令和3年度 問35}
\label{q:r03-reg-applied-35}

\noindent\textbf{原典：} sources/R03/17\_applied\_science.pdf、PDF物理13頁。

\textbf{公式正答：} 4（single）。

\par\textbf{検討未確定：} 原典のF・mは画像照合済み。比誘電率の空気と真空、10MHz分類目安及び各用途の媒質条件を区別。

\paragraph{問題の要旨} GPRの比誘電率の基準と誘電率の単位を照合する。

\Needspace{4\baselineskip}\paragraph{確認した条件・記述（独自要約）}\mbox{}\par

\noindent\textbf{原典照合：} 29問21原典画像で否定、全五候補、数式・上付き単位・観測対象を照合。\par

\noindent\textbf{前提：} 低損失非磁性媒質でv≈c/√εr、波長λ=v/f、パルス帯域Bに対する距離分解はおよそv/(2B)である。導電率等が到達深度を制限する。\par

\paragraph{問われる核心} GPRの比誘電率の基準と誘電率の単位を照合する。

\paragraph{必要知識} 低損失非磁性媒質でv≈c/√εr、波長λ=v/f、パルス帯域Bに対する距離分解はおよそv/(2B)である。導電率等が到達深度を制限する。

\paragraph{考え方・正答を選ぶ理由（査読待ち初稿）} 4が不適切である。D=εEからεの単位は(C/m\textsuperscript{2})/(V/m)=\allowbreak{}F/mで、F・mではない。比誘電率は真空ε0に対する比ε/ε0である。空気は通常近似として真空に近いが厳密な定義ではない。

\Needspace{4\baselineskip}\paragraph{選択肢の検討（査読待ち初稿）}\mbox{}\par

\noindent\textbf{1：適切} 誘電率・導電率の異なる境界からの反射を測る。電磁インピーダンスと減衰の違いを波形に捉える。地質名や機械強度を直接測る方式ではない。\par

\noindent\textbf{2：適切} アンテナ移動の波形を順に並べ記録断面を作る。横軸は測線位置、縦軸は通常往復走時であり、速度モデルから深度へ変換する。\par

\noindent\textbf{3：条件付きで適切} 中心周波数10MHz程度以上をGPRの目安とする。通常のGPRの周波数帯を示す分類目安である。電磁探査との厳密な物理境界ではない。\par

\noindent\textbf{4：不適切} 誘電率単位F・m、比誘電率は空気に対する比。誘電率はF/mであり、比誘電率の基準は真空である。空気の基準は通常の近似に過ぎない。\par

\noindent\textbf{5：条件付きで適切} 埋設物・空洞・地層・鉄筋・遺跡・水等を調査。各対象が電磁特性の差を持ち、十分な到達深度と分解能を得られる場合の用途である。すべての現場で適用可能とはしない。\par

\paragraph{再利用する考え方} 反射走時を深度と混同せず、帯域・速度・損失と探査目的を対応させる。

\paragraph{共通理論} \hyperref[theory:exploration:gpr]{第\ref*{theory:exploration:gpr}節}, \hyperref[theory:exploration:gprresolution]{第\ref*{theory:exploration:gprresolution}節}

\par 参考：\cite{epa_gpr}。

\par 難易度 3（校準初稿）、優先度未評価。

\subsection{令和5年度 問33}
\label{q:r05-reg-applied-33}

\noindent\textbf{原典：} sources/R05/17\_applied\_science.pdf、PDF物理12頁。

\textbf{公式正答：} 3（single）。

\par\textbf{検討未確定：} 分解能の表現は最小距離と識別能力を区別する。実深度のf依存を厳密な比例則とせず、力学特性との共通要因による間接相関を留保。

\paragraph{問題の要旨} GPRの周波数と分解能力・到達深度の向きを判断する。

\Needspace{4\baselineskip}\paragraph{確認した条件・記述（独自要約）}\mbox{}\par

\noindent\textbf{原典照合：} 29問21原典画像で否定、全五候補、数式・上付き単位・観測対象を照合。\par

\noindent\textbf{前提：} 低損失非磁性媒質でv≈c/√εr、波長λ=v/f、パルス帯域Bに対する距離分解はおよそv/(2B)である。導電率等が到達深度を制限する。\par

\paragraph{問われる核心} GPRの周波数と分解能力・到達深度の向きを判断する。

\paragraph{必要知識} 低損失非磁性媒質でv≈c/√εr、波長λ=v/f、パルス帯域Bに対する距離分解はおよそv/(2B)である。導電率等が到達深度を制限する。

\paragraph{考え方・正答を選ぶ理由（査読待ち初稿）} 3が不適切である。高周波では通常波長が短く分解能力が高まる一方、実地では探査深度が制限される。原典の分解能逆比例・深度比例はこの通常の向きを逆にしている。最小識別距離と能力の語を混同しない。

\Needspace{4\baselineskip}\paragraph{選択肢の検討（査読待ち初稿）}\mbox{}\par

\noindent\textbf{1：条件付きで適切} 電気境界を測るため力学特性との相関はほとんどない。強度や硬さを直接測らないことは正しい。ただし含水・空隙・劣化など共通要因を通じた間接相関まで全面否定はできない。\par

\noindent\textbf{2：適切} 直線埋設物へできる限り直角の測線。横切る測線で位置の変化と点状反射の走時を捉えやすくする。未知方位では複数方位を測る。\par

\noindent\textbf{3：不適切} 中心周波数と分解能は逆比例、探査深度は比例。分解できる最小距離は通常高周波ほど小さく、能力としての分解能は良くなる。深度も高周波ほど大きいという説明は通常の交換条件に反する。\par

\noindent\textbf{4：適切} profileにより構造、wide angleで電磁波速度を推定。走査波形を位置へ並べるprofileと、送受信距離を変えて走時から速度を求める構成を使い分ける。\par

\noindent\textbf{5：条件付きで適切} 中心周波数10MHz程度以上をGPRの目安とする。通常のGPRの周波数帯を示す分類目安である。電磁探査との厳密な物理境界ではない。\par

\paragraph{再利用する考え方} 反射走時を深度と混同せず、帯域・速度・損失と探査目的を対応させる。

\paragraph{共通理論} \hyperref[theory:exploration:gpr]{第\ref*{theory:exploration:gpr}節}, \hyperref[theory:exploration:gprresolution]{第\ref*{theory:exploration:gprresolution}節}

\par 参考：\cite{epa_gpr}。

\paragraph{関連問題} \hyperref[q:r02-reg-applied-33]{R02-REG-APPLIED-33}, \hyperref[q:r06-reg-applied-35]{R06-REG-APPLIED-35}

\par 難易度 3（校準初稿）、優先度未評価。

\subsection{令和6年度 問35}
\label{q:r06-reg-applied-35}

\noindent\textbf{原典：} sources/R06/17\_applied\_science.pdf、PDF物理12頁。

\textbf{公式正答：} 2（single）。

\par\textbf{検討未確定：} 分解能の最小距離/能力と探査深度の交換条件を区別し、10MHzを全方式の厳密境界にしない。

\paragraph{問題の要旨} GPRの電気境界と周波数・断面・用途を比較する。

\Needspace{4\baselineskip}\paragraph{確認した条件・記述（独自要約）}\mbox{}\par

\noindent\textbf{原典照合：} 29問21原典画像で否定、全五候補、数式・上付き単位・観測対象を照合。\par

\noindent\textbf{前提：} 低損失非磁性媒質でv≈c/√εr、波長λ=v/f、パルス帯域Bに対する距離分解はおよそv/(2B)である。導電率等が到達深度を制限する。\par

\paragraph{問われる核心} GPRの電気境界と周波数・断面・用途を比較する。

\paragraph{必要知識} 低損失非磁性媒質でv≈c/√εr、波長λ=v/f、パルス帯域Bに対する距離分解はおよそv/(2B)である。導電率等が到達深度を制限する。

\paragraph{考え方・正答を選ぶ理由（査読待ち初稿）} 2が不適切である。高周波は通常分解能力を良くする一方、深度を大きくする保証はない。帯域と媒質導電性・ノイズを含む到達条件から判断する。

\Needspace{4\baselineskip}\paragraph{選択肢の検討（査読待ち初稿）}\mbox{}\par

\noindent\textbf{1：適切} 誘電率・導電率の異なる境界からの反射を測る。電磁インピーダンスと減衰の違いを波形に捉える。地質名や機械強度を直接測る方式ではない。\par

\noindent\textbf{2：不適切} 中心周波数と分解能は逆比例、探査深度は比例。分解できる最小距離は通常高周波ほど小さく、能力としての分解能は良くなる。深度も高周波ほど大きいという説明は通常の交換条件に反する。\par

\noindent\textbf{3：条件付きで適切} 埋設物・空洞・地層・鉄筋・遺跡・水等を調査。各対象が電磁特性の差を持ち、十分な到達深度と分解能を得られる場合の用途である。すべての現場で適用可能とはしない。\par

\noindent\textbf{4：適切} アンテナ移動の波形を順に並べ記録断面を作る。横軸は測線位置、縦軸は通常往復走時であり、速度モデルから深度へ変換する。\par

\noindent\textbf{5：条件付きで適切} 中心周波数10MHz程度以上をGPRの目安とする。通常のGPRの周波数帯を示す分類目安である。電磁探査との厳密な物理境界ではない。\par

\paragraph{再利用する考え方} 反射走時を深度と混同せず、帯域・速度・損失と探査目的を対応させる。

\paragraph{共通理論} \hyperref[theory:exploration:gpr]{第\ref*{theory:exploration:gpr}節}, \hyperref[theory:exploration:gprresolution]{第\ref*{theory:exploration:gprresolution}節}

\par 参考：\cite{epa_gpr}。

\paragraph{関連問題} \hyperref[q:r05-reg-applied-33]{R05-REG-APPLIED-33}

\par 難易度 3（校準初稿）、優先度未評価。

\section{逆問題の正則化・最適化と離散化}

\subsection{平成30年度 問35}
\label{q:h30-reg-applied-35}

\noindent\textbf{原典：} sources/H30/17\_applied\_science.pdf、PDF物理12頁。

\textbf{公式正答：} 5（single）。

\par\textbf{検討未確定：} 設問のインバージョンに関係するという広い語には有限要素順モデルも含められる。公式5を算法の直接的役割の比較として保持し、FEMを逆解析から排除しない。

\paragraph{問題の要旨} 逆問題の最適化・正則化と有限要素離散化を区別する。

\Needspace{4\baselineskip}\paragraph{確認した条件・記述（独自要約）}\mbox{}\par

\noindent\textbf{原典照合：} 否定の対象、指数、周波数と周期、五つの候補の列順を原典画像で照合。\par

\noindent\textbf{前提：} モデル更新算法と順モデル計算の離散化を区別する。ただし順モデルは逆問題を構成する内部計算でもある。\par

\paragraph{問われる核心} 逆問題の最適化・正則化と有限要素離散化を区別する。

\paragraph{必要知識} モデル更新算法と順モデル計算の離散化を区別する。ただし順モデルは逆問題を構成する内部計算でもある。

\paragraph{考え方・正答を選ぶ理由（査読待ち初稿）} 公式5は、他四項目が最適化又は逆解の安定化を直接担うのに対し、有限要素法が主に支配方程式の離散化を担うという比較による。ただしFEMで順問題を計算する逆解析は広くある。インバージョンにFEMが関係しないと断定しない。

\Needspace{4\baselineskip}\paragraph{選択肢の検討（査読待ち初稿）}\mbox{}\par

\noindent\textbf{1：適切} マルカート法。非線形最小二乗の更新を減衰項で安定化するLevenberg–Marquardt型の算法である。\par

\noindent\textbf{2：適切} 遺伝的アルゴリズム。モデル候補の集団を評価・選択・変化させて目的関数を探索する最適化法である。大域的最良解を常に保証するわけではない。\par

\noindent\textbf{3：適切} 正則化法。データだけで決まりにくいモデルへ滑らかさや大きさ等の制約を加え、逆問題の不安定性を抑える。\par

\noindent\textbf{4：適切} 共役勾配法。線形化された最小二乗等の連立系又は非線形目的関数の探索に用いる。対称正定値系のCGと非線形CGを区別する。\par

\noindent\textbf{5：条件付きで適切} 有限要素法。直接の最適化・正則化の名称ではなく支配方程式の離散化であるため公式はこれを選ぶ。しかし順モデルや感度計算として逆解析へ用いることができ、無関係という絶対判断は誤りである。\par

\paragraph{再利用する考え方} 目的関数、モデル更新、物理場の数値計算の三つの役割を別々に問う。

\paragraph{共通理論} \hyperref[theory:exploration:inverse]{第\ref*{theory:exploration:inverse}節}, \hyperref[theory:fem:background]{第\ref*{theory:fem:background}節}

\par 参考：\cite{pygimli_inversion,mit_roy_fea1999}。

\paragraph{関連問題} \hyperref[q:r01-reg-applied-35]{R01-REG-APPLIED-35}

\par 難易度 3（校準初稿）、優先度未評価。

\subsection{令和元年度 問35}
\label{q:r01-reg-applied-35}

\noindent\textbf{原典：} sources/R01/17\_applied\_science.pdf、PDF物理12頁。

\textbf{公式正答：} 2（single）。

\paragraph{問題の要旨} フィルタ・差分・共役勾配・平滑化とラテラル検層を区別する。

\Needspace{4\baselineskip}\paragraph{確認した条件・記述（独自要約）}\mbox{}\par

\noindent\textbf{原典照合：} 否定の対象、指数、周波数と周期、五つの候補の列順を原典画像で照合。\par

\noindent\textbf{前提：} 数値処理又はモデル拘束と、孔内で物理量を測る配置を区別する。\par

\paragraph{問われる核心} フィルタ・差分・共役勾配・平滑化とラテラル検層を区別する。

\paragraph{必要知識} 数値処理又はモデル拘束と、孔内で物理量を測る配置を区別する。

\paragraph{考え方・正答を選ぶ理由（査読待ち初稿）} 2が不適切である。ラテラル法は電極配置を用いる電気検層の測定方式である。他四項目は信号処理、方程式の離散化、最適化、解への拘束を担う。

\Needspace{4\baselineskip}\paragraph{選択肢の検討（査読待ち初稿）}\mbox{}\par

\noindent\textbf{1：適切} リニアフィルタ。入力の和と定数倍に対して同じ関係を保つ信号処理であり、畳み込み等の数値計算で実装できる。\par

\noindent\textbf{2：不適切} ラテラル法。USGSの電気検層資料では電極間隔と配置で定義されるlateral logとして確認できる。数値解法の算法を指す語ではない。\par

\noindent\textbf{3：適切} 差分法。微分演算を格子値の差へ置き換えて支配方程式の近似解を求める数値法である。\par

\noindent\textbf{4：適切} 共役勾配法。大型の連立系や最小化を反復して解く数値算法である。\par

\noindent\textbf{5：適切} 平滑化拘束。隣接するモデル値の差などを罰することで、逆解の不要な振動を抑える数値的な条件である。\par

\paragraph{再利用する考え方} 測定の幾何配置とデータ取得後の計算を分ける。

\paragraph{共通理論} \hyperref[theory:exploration:inverse]{第\ref*{theory:exploration:inverse}節}, \hyperref[theory:exploration:seismicterms]{第\ref*{theory:exploration:seismicterms}節}

\par 参考：\cite{usgs_lateral_logging,pygimli_inversion,mit_signals2010}。

\paragraph{関連問題} \hyperref[q:h30-reg-applied-35]{H30-REG-APPLIED-35}

\par 難易度 2（校準初稿）、優先度未評価。

\section{人工分極と自然電位の信号機構}

\subsection{平成28年度 問22}
\label{q:h28-reg-applied-22}

\noindent\textbf{原典：} sources/H28/17\_applied\_science.pdf、PDF物理9頁。

\textbf{公式正答：} 1（single）。

\par\textbf{検討未確定：} SPの負異常は典型的鉱化機構と基準点による。AMTの深部は相対表現であり一定深度を保証しない。

\paragraph{問題の要旨} SP鉱体異常の符号と電気・弾性波探査を比較する。

\Needspace{4\baselineskip}\paragraph{確認した条件・記述（独自要約）}\mbox{}\par

\noindent\textbf{原典照合：} 29問21原典画像で否定、全五候補、数式・上付き単位・観測対象を照合。\par

\noindent\textbf{前提：} SPは基準点に対する電位差であり、典型的な硫化鉱体上の負異常を考える。\par

\paragraph{問われる核心} SP鉱体異常の符号と電気・弾性波探査を比較する。

\paragraph{必要知識} SPは基準点に対する電位差であり、典型的な硫化鉱体上の負異常を考える。

\paragraph{考え方・正答を選ぶ理由（査読待ち初稿）} 1が不適切である。典型的な鉱化電池の鉱体上部では基準点に対し負の電位異常を利用する。全鉱体・全自然電位源の符号を固定する規則ではない。

\Needspace{4\baselineskip}\paragraph{選択肢の検討（査読待ち初稿）}\mbox{}\par

\noindent\textbf{1：不適切} 鉱体があると地表電位が周囲より高い。典型的な硫化鉱体探査は負の電位異常を利用し、周囲より高いという説明は逆である。酸化還元や流動の機構を特定する。\par

\noindent\textbf{2：適切} 比抵抗垂直探査は電極間隔を広げ深さ分布を推定。配置間隔を変えて深度感度を変え、水平多層等のモデルを推定する。\par

\noindent\textbf{3：条件付きで適切} AMTは自然磁場と地電流で深部比抵抗を推定。電場磁場の周波数応答で表層レーダより深い領域へ感度を持つ。実深度は周波数・抵抗率・信号品質による。\par

\noindent\textbf{4：適切} 弾性波tomographyは周囲の起振受振で伝播特性を画像化。多数経路走時などの逆解析で速度分布を求め、覆いの不足は分解能を制限する。\par

\noindent\textbf{5：適切} ダウンホールで地表起振・孔内受振の初動走時を測る。孔壁に結合した受振器の位置と走時から区間速度を推定する。経路と屈折の補正が必要である。\par

\paragraph{再利用する考え方} 自然電位の符号を単一の地下物体の存在へ直結させない。

\paragraph{共通理論} \hyperref[theory:exploration:ipsp]{第\ref*{theory:exploration:ipsp}節}, \hyperref[theory:exploration:dcarrays]{第\ref*{theory:exploration:dcarrays}節}, \hyperref[theory:exploration:borehole]{第\ref*{theory:exploration:borehole}節}

\par 参考：\cite{epa_sp,epa_resistivity,usgs_mt_metadata,geebec_ps}。

\par 難易度 3（校準初稿）、優先度未評価。

\subsection{令和3年度 問32}
\label{q:r03-reg-applied-32}

\noindent\textbf{原典：} sources/R03/17\_applied\_science.pdf、PDF物理12頁。

\textbf{公式正答：} 4（single）。

\par\textbf{検討未確定：} 4にはモデル依存の深度推定が可能という反例がある。公式4を保持するが「SPで深度解析はできない」という誤解を教えない。3の特定深さも有限感度体積の近似。

\paragraph{問題の要旨} 比抵抗の深度感度と自然電位法の逆解析限界を比較する。

\Needspace{4\baselineskip}\paragraph{確認した条件・記述（独自要約）}\mbox{}\par

\noindent\textbf{原典照合：} 29問21原典画像で否定、全五候補、数式・上付き単位・観測対象を照合。\par

\noindent\textbf{前提：} SPは自然電位差の分布を直接測り、源の深度には機構・形状・抵抗率のモデルが必要である。\par

\paragraph{問われる核心} 比抵抗の深度感度と自然電位法の逆解析限界を比較する。

\paragraph{必要知識} SPは自然電位差の分布を直接測り、源の深度には機構・形状・抵抗率のモデルが必要である。

\paragraph{考え方・正答を選ぶ理由（査読待ち初稿）} 公式の4を保持する。SP観測から地下水や鉱体の深さを簡単に直接求めるという読みなら不適切である。一方、半幅や源モデル、他観測との逆解析により深度推定の例はあり、深度を探査・解析するという広い原典表現自体が全面的に不可能とは言えない。疑義を査読対象にする。

\Needspace{4\baselineskip}\paragraph{選択肢の検討（査読待ち初稿）}\mbox{}\par

\noindent\textbf{1：適切} IPは印加電流の分極応答、硫化鉱物探査に用いる。金属・電解質界面等の分極で電位の時間減衰や周波数位相が生じ、鉱床や環境調査に使われる。\par

\noindent\textbf{2：適切} 垂直探査は水平多層の鉛直比抵抗を推定。電極間隔を変えて深度感度を変化させ、WennerやSchlumberger等の配置で層モデルを推定する。\par

\noindent\textbf{3：条件付きで適切} 水平探査は特定深さの水平比抵抗変化を調べる。固定間隔等で感度深度をおおむねそろえ横方向を調べる。ただし有限体積の重み付き応答で一点の厳密深度だけを測らない。\par

\noindent\textbf{4：条件付きで適切} SPから地下水や鉱床の深さを探査解析し簡単安全に測る。自然電位分布の簡単な受動測定はできるが、深度は非一意で源モデルや他データが必要である。深度解析が全く不可能とも言えず、この候補の否定範囲は未確定である。\par

\noindent\textbf{5：適切} 多電極の直流電位から比抵抗tomography。多数の注入・受信対の応答を順モデルへ合わせて比抵抗分布を推定する。\par

\paragraph{再利用する考え方} 公式の選択と、科学的に否定できる命題の範囲を分けて残す。

\paragraph{共通理論} \hyperref[theory:exploration:ipsp]{第\ref*{theory:exploration:ipsp}節}, \hyperref[theory:exploration:dcarrays]{第\ref*{theory:exploration:dcarrays}節}

\par 参考：\cite{epa_sp,epa_ip,epa_resistivity}。

\par 難易度 3（校準初稿）、優先度未評価。

\section{探査の観測量と物理手法の対応}

\subsection{令和元年度再試験 問32}
\label{q:r01-ret-applied-32}

\noindent\textbf{原典：} sources/R01\_retest/17\_applied\_science.pdf、PDF物理11頁。

\textbf{公式正答：} 4（single）。

\paragraph{問題の要旨} 地震探査を重力・磁気・電気・放射能探査から観測量で選ぶ。

\Needspace{4\baselineskip}\paragraph{確認した条件・記述（独自要約）}\mbox{}\par

\noindent\textbf{原典照合：} 五候補の順序、最適切と不適切の方向、波種、深度目安、曲線名、最長と最短を原典画像で照合。\par

\noindent\textbf{前提：} 最も適切な説明を選ぶ。弾性波の人工震源と地下からの応答を対象とする。\par

\paragraph{問われる核心} 地震探査を重力・磁気・電気・放射能探査から観測量で選ぶ。

\paragraph{必要知識} 最も適切な説明を選ぶ。弾性波の人工震源と地下からの応答を対象とする。

\paragraph{考え方・正答を選ぶ理由（査読待ち初稿）} 4が適切である。爆薬・バイブレータ・エアガンで弾性波を送り、反射や屈折から構造を推定する。地質を調べるという目的が同じでも、他の四つは測定する物理量が異なる。

\Needspace{4\baselineskip}\paragraph{選択肢の検討（査読待ち初稿）}\mbox{}\par

\noindent\textbf{1：不適切} ブーゲー異常で基盤・褶曲・断層を考察する。密度差による重力探査であり弾性波の走時・波形を直接測る地震探査ではない。\par

\noindent\textbf{2：不適切} 磁性鉱床・堆積盆地・地殻熱構造を調べる。磁気異常や磁化の性質を使う探査の説明である。\par

\noindent\textbf{3：不適切} 電気的構造と深度別の比抵抗を知る。電気・電磁応答から導電率又は比抵抗を推定する説明である。\par

\noindent\textbf{4：適切} 人工起振で弾性波を地下へ送り資源・土木調査に用いる。地震探査の震源と観測量に対応する。海上のエアガンなど場所ごとに装置は異なる。\par

\noindent\textbf{5：不適切} U・Th・K等の放射性元素を探査する。放射能探査の観測対象であり弾性波とは異なる。天然核種のガンマ線等を利用する。\par

\paragraph{再利用する考え方} 手法の目的が同じ場合も入力・出力の物理量で区別する。

\paragraph{共通理論} \hyperref[theory:exploration:observables]{第\ref*{theory:exploration:observables}節}

\par 参考：\cite{epa_seismic_reflection,segj_integrated2017}。

\par 難易度 2（校準初稿）、優先度未評価。

\section{移動平均の雑音低減と周波数零点}

\subsection{平成28年度 問20}
\label{q:h28-reg-applied-20}

\noindent\textbf{原典：} sources/H28/17\_applied\_science.pdf、PDF物理8頁。

\textbf{公式正答：} 2（single）。

\par\textbf{検討未確定：} 2には周期と窓長を合わせれば除去できる具体的反例がある。公式2を保持するが、数学的に周期雑音除去が不可能とは説明しない。3の時不変性と5の帯域・端点条件も省略される。

\paragraph{問題の要旨} 移動平均と周期雑音除去の関係を検討する。

\Needspace{4\baselineskip}\paragraph{確認した条件・記述（独自要約）}\mbox{}\par

\noindent\textbf{原典照合：} 否定の対象、指数、周波数と周期、五つの候補の列順を原典画像で照合。\par

\noindent\textbf{前提：} 通常の移動平均は平滑化を行うが、窓長に応じた周波数応答の零点を持つ。\par

\paragraph{問われる核心} 移動平均と周期雑音除去の関係を検討する。

\paragraph{必要知識} 通常の移動平均は平滑化を行うが、窓長に応じた周波数応答の零点を持つ。

\paragraph{考え方・正答を選ぶ理由（査読待ち初稿）} 公式は2であり、一般的な移動平均を周期雑音専用の選択除去と説明する意図を不適切と読んでいる。しかし周期と窓長が一致する正弦波は移動平均で厳密に消せる。したがって周期雑音を除去できないという一般則へ書き換えず、原典の曖昧さを残す。

\Needspace{4\baselineskip}\paragraph{選択肢の検討（査読待ち初稿）}\mbox{}\par

\noindent\textbf{1：適切} 空間フィルタを二次元画像のエッジ・平滑化・差分に用いる。二次元データに対し空間方向の重み付けや差分を行う。一次元・三次元にも拡張でき、二次元しか扱えないという意味ではない。\par

\noindent\textbf{2：条件付きで適切} 移動平均は時系列の周期的雑音を除去するために用いる。M点等重み平均の周波数応答はsin(Mω/2)/(M sin(ω/2))に位相因子を掛けたものである。ω=2πk/Mの非DC成分は消えるため、窓長を周期へ合わせる反例がある。一般に任意の周期雑音が消えるとは限らない。\par

\noindent\textbf{3：条件付きで適切} 線形システムのインパルス応答を重ねて任意入力への応答を求める。線形性により各時刻のインパルスへの応答を重ねることはできる。単一の応答を時間移動する畳み込み表現には時不変性も必要であり、時変系には二時刻の応答を用いる。\par

\noindent\textbf{4：条件付きで適切} 低域通過で高周波を減衰させ低周波を通す。低域通過の基本的な役割である。無減衰という文言は理想通過帯域のモデルであり、実装の位相・リップルは別に確認する。\par

\noindent\textbf{5：条件付きで適切} 一周期二標本未満では元データを再構成できない。一般の帯域制限信号を一意に再構成するにはスペクトルの重なりを避ける標本化が必要である。特別な事前情報を持つ信号まで一律に不可能とする定理ではなく、最高周波数と標本化間隔を明示して読む。\par

\paragraph{再利用する考え方} 移動平均の効果を目的名で決めず、入力周波数と窓長の応答を計算する。

\paragraph{共通理論} \hyperref[theory:exploration:moving]{第\ref*{theory:exploration:moving}節}, \hyperref[theory:exploration:lti]{第\ref*{theory:exploration:lti}節}, \hyperref[theory:exploration:sampling]{第\ref*{theory:exploration:sampling}節}

\par 参考：\cite{mit_signals2010,mit_moving2012,scipy_notch,scipy_fft,usgs_processing2000}。

\paragraph{関連問題} \hyperref[q:h23-reg-applied-28]{H23-REG-APPLIED-28}

\par 難易度 3（校準初稿）、優先度未評価。

\section{物理探査の複数手法と幾何・時間}

\subsection{平成25年度 問33}
\label{q:h25-reg-applied-33}

\noindent\textbf{原典：} sources/H25/17\_applied\_science.pdf、PDF物理12頁。

\textbf{公式正答：} 5（single）。

\par\textbf{検討未確定：} 地形起伏があれば反射法は絶対不可とはしない。速度比だけから層厚精度を保証せず、隠れ層と観測条件を保持。

\paragraph{問題の要旨} 重力・速度検層・屈折法と起伏地の反射法の制約を比較する。

\Needspace{4\baselineskip}\paragraph{確認した条件・記述（独自要約）}\mbox{}\par

\noindent\textbf{原典照合：} 29問21原典画像で否定、全五候補、数式・上付き単位・観測対象を照合。\par

\noindent\textbf{前提：} 地形起伏は走時・観測配置へ影響し、浅層反射法では静補正や起振受振条件を整える。\par

\paragraph{問われる核心} 重力・速度検層・屈折法と起伏地の反射法の制約を比較する。

\paragraph{必要知識} 地形起伏は走時・観測配置へ影響し、浅層反射法では静補正や起振受振条件を整える。

\paragraph{考え方・正答を選ぶ理由（査読待ち初稿）} 5が不適切である。急峻で凹凸の多い地形を浅層反射法の一般的な適用対象とする説明は、走時補正・起振結合・配置の難しさを落としている。ただし適切な補正と計測により適用可能な例を否定しない。

\Needspace{4\baselineskip}\paragraph{選択肢の検討（査読待ち初稿）}\mbox{}\par

\noindent\textbf{1：適切} 重力分布から地下密度構造を推定。重力異常に対して密度と形状のモデルを作り逆解析する。\par

\noindent\textbf{2：適切} 単孔で波速の深さ分布を求める速度検層。地表又は孔内起振と孔内受振の走時を用いる。\par

\noindent\textbf{3：条件付きで適切} 屈折波から地下速度構造を解析。通常は速い下層の境界を進み戻る臨界屈折波の初動等を使う。任意の低速度層まで検出できるわけではない。\par

\noindent\textbf{4：条件付きで適切} 下方増速で速度比大きいほど分割と厚さ精度が良い。層走時勾配の差が明瞭になり識別しやすい傾向である。隠れ層、薄層、配置、ノイズによる精度制限もある。\par

\noindent\textbf{5：不適切} 浅層反射法は急峻で凹凸の多い場所へ適用される。通常の適用上は地形起伏が観測と補正の障害となり、適性を積極的に示す一般記述は不適切である。実施そのものが絶対不可能とはしない。\par

\paragraph{再利用する考え方} 適用の一般的有利条件と、補正して実施できる可能性を分ける。

\paragraph{共通理論} \hyperref[theory:exploration:refraction]{第\ref*{theory:exploration:refraction}節}, \hyperref[theory:exploration:reflectiongeometry]{第\ref*{theory:exploration:reflectiongeometry}節}

\par 参考：\cite{epa_seismic_reflection,epa_seismic_refraction,usgs_static2004,geebec_ps}。

\par 難易度 3（校準初稿）、優先度未評価。

\subsection{平成26年度 問19}
\label{q:h26-reg-applied-19}

\noindent\textbf{原典：} sources/H26/17\_applied\_science.pdf、PDF物理8頁。

\textbf{公式正答：} 5（single）。

\par\textbf{検討未確定：} 含水比・飽和度・間隙流体を区別し、P波速度の全材料での単調変化を主張しない。P波初動の利点は通常の方式に限定する。

\paragraph{問題の要旨} 各物理探査の観測量とP波速度の含水依存を比較する。

\Needspace{4\baselineskip}\paragraph{確認した条件・記述（独自要約）}\mbox{}\par

\noindent\textbf{原典照合：} 29問21原典画像で否定、全五候補、数式・上付き単位・観測対象を照合。\par

\noindent\textbf{前提：} P波速度は弾性率と密度で決まり、飽和や孔隙流体は両者と接触状態を変え得る。\par

\paragraph{問われる核心} 各物理探査の観測量とP波速度の含水依存を比較する。

\paragraph{必要知識} P波速度は弾性率と密度で決まり、飽和や孔隙流体は両者と接触状態を変え得る。

\paragraph{考え方・正答を選ぶ理由（査読待ち初稿）} 5が不適切である。P波速度は飽和・含水状態で変わり得る。Vp=√((K+4µ/3)/ρ)の剛性と密度が含水で変化するため、一般に不変とすることはできない。単純な含水比だけから全地盤の単調変化は予測しない。

\Needspace{4\baselineskip}\paragraph{選択肢の検討（査読待ち初稿）}\mbox{}\par

\noindent\textbf{1：適切} 重力が密度・深度・形状を反映する。密度差の体積積分が重力へ寄与し、同じ密度でも深度と形状で観測異常が変わる。\par

\noindent\textbf{2：条件付きで適切} 通常の屈折法は地表でP波初動を測る。地表から広い対象を調査しP初動を使う構成が一般的である。S波方式や隠れ層の限界もある。\par

\noindent\textbf{3：適切} ダウンホールは孔内速度測定方式。地表起振・孔内受振の走時を使う。\par

\noindent\textbf{4：適切} ガンマ線で深度方向の密度検層。線源と受信器を降下させ孔壁周囲の校正かさ密度を測る。\par

\noindent\textbf{5：不適切} P波速度は一般に含水比によらず不変。孔隙流体による有効弾性率と密度・接触の変化があるため、含水依存を一般に除外できない。\par

\paragraph{再利用する考え方} 一つの観測物性を水分と切り離さず、波速と密度の感度を考える。

\paragraph{共通理論} \hyperref[theory:exploration:borehole]{第\ref*{theory:exploration:borehole}節}, \hyperref[theory:exploration:waveproperties]{第\ref*{theory:exploration:waveproperties}節}

\par 参考：\cite{epa_gamma_density,geebec_ps,epa_seismic_refraction}。

\par 難易度 3（校準初稿）、優先度未評価。

\subsection{平成29年度 問33}
\label{q:h29-reg-applied-33}

\noindent\textbf{原典：} sources/H29/17\_applied\_science.pdf、PDF物理11頁。

\textbf{公式正答：} 3（single）。

\par\textbf{検討未確定：} 磁気単独の非一意性を形状推定の絶対不可能性にしない。GPRの10MHzは典型的分類目安である。

\paragraph{問題の要旨} ばね式重力計・MT・GPR・海上音波の観測を比較する。

\Needspace{4\baselineskip}\paragraph{確認した条件・記述（独自要約）}\mbox{}\par

\noindent\textbf{原典照合：} 29問21原典画像で否定、全五候補、数式・上付き単位・観測対象を照合。\par

\noindent\textbf{前提：} ばね伸びの比較による重力差と、絶対加速度の時間長さ測定を分ける。\par

\paragraph{問われる核心} ばね式重力計・MT・GPR・海上音波の観測を比較する。

\paragraph{必要知識} ばね伸びの比較による重力差と、絶対加速度の時間長さ測定を分ける。

\paragraph{考え方・正答を選ぶ理由（査読待ち初稿）} 3が不適切である。ばねの伸びを使う多くの重力探査用装置は相対重力計である。絶対計は自由落下等の運動を長さと時間で測る。

\Needspace{4\baselineskip}\paragraph{選択肢の検討（査読待ち初稿）}\mbox{}\par

\noindent\textbf{1：条件付きで適切} 磁気探査で鉄の概略位置は測れるが細形状種類は判別できない。磁化と形状・深度の非一意性があるため磁気だけで製品種を一意に識別しない。追加観測で形状推定ができないという絶対禁止ではない。\par

\noindent\textbf{2：適切} MTは自然磁場と地電流から深部比抵抗を探査。電磁応答のインピーダンスと周波数から比抵抗モデルを求める。\par

\noindent\textbf{3：不適切} ばね伸びの装置は絶対重力計。標準点や基準との重力差を測る相対方式であり絶対方式と取り違える。\par

\noindent\textbf{4：条件付きで適切} GPRは透過反射屈折を使い中心周波数10MHz程度以上。通常のGPRの高周波帯の説明である。10MHzは厳密な全装置・全媒質の境界ではない。\par

\noindent\textbf{5：適切} 海上ストリーマ受振器はハイドロフォン。水中の音圧変動を受信する素子を曳航ケーブルへ多数配置する。\par

\paragraph{再利用する考え方} センサが測る量と、逆解析で推定する地下構造を分ける。

\paragraph{共通理論} \hyperref[theory:exploration:observables]{第\ref*{theory:exploration:observables}節}, \hyperref[theory:exploration:gpr]{第\ref*{theory:exploration:gpr}節}

\par 参考：\cite{gsi_gravity_definitions,usgs_mt_metadata,epa_gpr,usgs_marine_seismic}。

\par 難易度 3（校準初稿）、優先度未評価。

\subsection{令和元年度 問32}
\label{q:r01-reg-applied-32}

\noindent\textbf{原典：} sources/R01/17\_applied\_science.pdf、PDF物理11頁。

\textbf{公式正答：} 5（single）。

\par\textbf{検討未確定：} JISの規格本文と出題時/現行版差は未照合。重力ののみは処理目的と読み、残差誤差がないとはしない。

\paragraph{問題の要旨} 探査の時間差・環境振動と表面波のS波速度を比較する。

\Needspace{4\baselineskip}\paragraph{確認した条件・記述（独自要約）}\mbox{}\par

\noindent\textbf{原典照合：} 29問21原典画像で否定、全五候補、数式・上付き単位・観測対象を照合。\par

\noindent\textbf{前提：} Rayleigh波の分散は主にVs分布へ関係し、時間差調査は同じ測定条件を比較する。\par

\paragraph{問われる核心} 探査の時間差・環境振動と表面波のS波速度を比較する。

\paragraph{必要知識} Rayleigh波の分散は主にVs分布へ関係し、時間差調査は同じ測定条件を比較する。

\paragraph{考え方・正答を選ぶ理由（査読待ち初稿）} 5が不適切である。表面波探査は分散曲線から主にS波速度構造を推定する。Rayleigh波にP成分もあるが、主推定対象をP波速度とする説明は違う。

\Needspace{4\baselineskip}\paragraph{選択肢の検討（査読待ち初稿）}\mbox{}\par

\noindent\textbf{1：条件付きで適切} 重力補正で密度構造による異常のみ抽出。潮汐・標高・地形等の寄与を補正し地下密度異常を解釈する。補正誤差や地域背景が完全に消えるという保証ではない。\par

\noindent\textbf{2：適切} GPRは数m浅部の埋設物空洞遺跡に多用。典型的な浅部の利用であり、媒質・装置によってより深い例もある。\par

\noindent\textbf{3：適切} time lapseは繰返し探査差から地下変化を推定。同じ条件の反復観測で変化を見る。装置ドリフトや位置・気象差も区別する。\par

\noindent\textbf{4：条件付きで適切} 環境振動は感覚補正の振動レベルとJIS計器を使う。人の周波数感度を補正した加速度レベルで評価する。出題時と現在の規格版、計器性能の全文は追加照合対象である。\par

\noindent\textbf{5：不適切} 表面波分散から主にP波速度構造を推定。通常の表面波法の主推定物性はS波速度Vsである。P波速度は感度を持つ他パラメータとして扱う。\par

\paragraph{再利用する考え方} 観測波の粒子運動と、その逆解析が主に拘束する物性を区別する。

\paragraph{共通理論} \hyperref[theory:exploration:surfacewaves]{第\ref*{theory:exploration:surfacewaves}節}, \hyperref[theory:exploration:multimethod]{第\ref*{theory:exploration:multimethod}節}

\par 参考：\cite{epa_surfacewaves,epa_gpr,rion_vibration2016,gsi_gravity_definitions}。

\paragraph{関連問題} \hyperref[q:r05-reg-applied-35]{R05-REG-APPLIED-35}

\par 難易度 3（校準初稿）、優先度未評価。

\subsection{令和5年度 問35}
\label{q:r05-reg-applied-35}

\noindent\textbf{原典：} sources/R05/17\_applied\_science.pdf、PDF物理13頁。

\textbf{公式正答：} 2（single）。

\par\textbf{検討未確定：} 磁気数m・表面波20mは設定依存、熱探査は温度差と季節・地況による。1mを水脈の上限深度としない。

\paragraph{問題の要旨} 表面波のVsと磁気・tomography・地温の感度を比較する。

\Needspace{4\baselineskip}\paragraph{確認した条件・記述（独自要約）}\mbox{}\par

\noindent\textbf{原典照合：} 29問21原典画像で否定、全五候補、数式・上付き単位・観測対象を照合。\par

\noindent\textbf{前提：} 観測の空間覆い、密度差、温度差と、Rayleigh分散の主Vs感度を考える。\par

\paragraph{問われる核心} 表面波のVsと磁気・tomography・地温の感度を比較する。

\paragraph{必要知識} 観測の空間覆い、密度差、温度差と、Rayleigh分散の主Vs感度を考える。

\paragraph{考え方・正答を選ぶ理由（査読待ち初稿）} 2が不適切である。Rayleigh波分散から通常主にS波速度分布を推定する。20mは典型装置・設定の深度であり、全表面波法の深度限界ではない。

\Needspace{4\baselineskip}\paragraph{選択肢の検討（査読待ち初稿）}\mbox{}\par

\noindent\textbf{1：条件付きで適切} 地上磁気探査で水平分解能数mの浅部調査。細かい測点と近い対象で詳細調査が可能である。実分解能は深度、磁化、観測配置、ノイズによる。\par

\noindent\textbf{2：不適切} Rayleigh波から20m程度までP波速度分布。主推定対象はS波速度Vsである。深度の数値も測定波長・配置に依存する。\par

\noindent\textbf{3：適切} 地下tomographyの構造像は測定幾何に依存。全方向から測れず、覆いのない方向・領域の分解は制限される。\par

\noindent\textbf{4：条件付きで適切} 密度差が大きい対象は重力で検出しやすい。同じ形状・深度・ノイズ条件なら密度差が重力異常を増す。深い小対象との無条件比較ではない。\par

\noindent\textbf{5：条件付きで適切} 1m深地温で地熱温泉や地すべり地下水を調査。流動水と平常地温の温度差を利用する。測温1mと地下水源の深度は別で、適期や地表環境の補正が必要である。\par

\paragraph{再利用する考え方} 典型的数値を固定限界とせず、主観測量・感度・配置をそろえて比較する。

\paragraph{共通理論} \hyperref[theory:exploration:surfacewaves]{第\ref*{theory:exploration:surfacewaves}節}, \hyperref[theory:exploration:multimethod]{第\ref*{theory:exploration:multimethod}節}

\par 参考：\cite{epa_surfacewaves,gsi_gravity_definitions,naro_thermal2008,landslide_thermal2007}。

\paragraph{関連問題} \hyperref[q:r01-reg-applied-32]{R01-REG-APPLIED-32}

\par 難易度 3（校準初稿）、優先度未評価。

\section{CSAMTの表皮深度・近接場・平面波}

\subsection{平成25年度 問35}
\label{q:h25-reg-applied-35}

\noindent\textbf{原典：} sources/H25/17\_applied\_science.pdf、PDF物理13頁。

\textbf{公式正答：} 2（single）。

\par\textbf{検討未確定：} 3の距離によらずという全称にも疑義がある。CSAMT遠方条件の明確な欠落を2の根拠とするが、TDEMの無条件深度増大を承認しない。VLF局と周波数は方式の典型である。

\paragraph{問題の要旨} MT・CSAMT・TDEM・Slingram・VLFの深度条件を比較する。

\Needspace{4\baselineskip}\paragraph{確認した条件・記述（独自要約）}\mbox{}\par

\noindent\textbf{原典照合：} 29問21原典画像で否定、全五候補、数式・上付き単位・観測対象を照合。\par

\noindent\textbf{前提：} CSAMTの平面波近似には送受信距離が表皮深度より十分長い条件が必要である。\par

\paragraph{問われる核心} MT・CSAMT・TDEM・Slingram・VLFの深度条件を比較する。

\paragraph{必要知識} CSAMTの平面波近似には送受信距離が表皮深度より十分長い条件が必要である。

\paragraph{考え方・正答を選ぶ理由（査読待ち初稿）} 2が不適切である。δ=√(2ρ/(ωµ))は低周波ほど増す。固定送受信距離で周波数を下げ続けると遠方条件が破れ、近接場となるため距離によらず深く測れるとは言えない。3も装置規模・時間・ノイズから独立ではなく、その表現の疑義を保持する。

\Needspace{4\baselineskip}\paragraph{選択肢の検討（査読待ち初稿）}\mbox{}\par

\noindent\textbf{1：条件付きで適切} MTは自然電磁場を使い低周波で深部を探る。人工源を要せず低周波に深部感度があるが、信号強度と電気構造が推定を制限する。\par

\noindent\textbf{2：不適切} CSAMTは距離によらず低周波化で深度増大。有限送信源の近接場と平面波条件を無視する。固定距離では低周波ほどδが大きくなる。\par

\noindent\textbf{3：条件付きで適切} TDEMは遮断後応答を測り距離によらず深度増大。遮断後の時間応答を使いcentral loopでも深部に感度を持つ。ただし寸法・送信強度・測定時間・ノイズを無視した深度保証はできない。\par

\noindent\textbf{4：適切} Slingramは小型送受信コイルで浅部を迅速測定。周波数領域の二次磁場の振幅・位相を使い、浅部導電性分布を機動的に調べる。\par

\noindent\textbf{5：適切} VLFは20kHz付近、方向に注意して浅部探査。遠方送信局の電磁波と導電性構造の応答を使い、送信局に対する構造方位が感度へ影響する。\par

\paragraph{再利用する考え方} 深さを周波数だけから決めず、有限源と電磁拡散の条件を確認する。

\paragraph{共通理論} \hyperref[theory:exploration:nearfield]{第\ref*{theory:exploration:nearfield}節}, \hyperref[theory:exploration:fdtd]{第\ref*{theory:exploration:fdtd}節}

\par 参考：\cite{usgs_csamt2012,epa_tdem_archive,epa_fdem,usgs_mt_metadata}。

\par 難易度 3（校準初稿）、優先度未評価。

\subsection{平成26年度 問22}
\label{q:h26-reg-applied-22}

\noindent\textbf{原典：} sources/H26/17\_applied\_science.pdf、PDF物理9頁。

\textbf{公式正答：} 3（single）。

\par\textbf{検討未確定：} Archieの骨格・表面伝導条件、TEM一次場遮断の過渡とstatic shiftの比較を留保。生じにくいという原典限定を保持。

\paragraph{問題の要旨} Archie・電磁結合・CSAMT近接場・VLF・TEMを比較する。

\Needspace{4\baselineskip}\paragraph{確認した条件・記述（独自要約）}\mbox{}\par

\noindent\textbf{原典照合：} 29問21原典画像で否定、全五候補、数式・上付き単位・観測対象を照合。\par

\noindent\textbf{前提：} 遠方条件の指標は送受信距離/δであり、同一周波数では高比抵抗ほどδが増す。\par

\paragraph{問われる核心} Archie・電磁結合・CSAMT近接場・VLF・TEMを比較する。

\paragraph{必要知識} 遠方条件の指標は送受信距離/δであり、同一周波数では高比抵抗ほどδが増す。

\paragraph{考え方・正答を選ぶ理由（査読待ち初稿）} 3が不適切である。間隔が小さいほど近接場になり、同じ周波数・間隔では高比抵抗ほどδが大きく平面波条件が破れやすい。原典の低比抵抗を高比抵抗へ書き換えず、誤記述として保存する。

\Needspace{4\baselineskip}\paragraph{選択肢の検討（査読待ち初稿）}\mbox{}\par

\noindent\textbf{1：条件付きで適切} Archieは孔隙水だけの電流を仮定する。絶縁性骨格・表面伝導が無視できる孔隙ネットワークを想定する。低抵抗流体はその条件を満たしやすいが、有効孔隙率が大きければ十分とは限らない。\par

\noindent\textbf{2：適切} 電磁カップリングがIP誤差になる。送受信ケーブル間の誘導応答は分極による位相・減衰応答と重なり誤差へ寄与する。\par

\noindent\textbf{3：不適切} CSAMT近接場は低比抵抗で生じやすい。δは√ρに比例するため同一条件では高比抵抗の方が有限距離による近接場になりやすい。\par

\noindent\textbf{4：適切} VLFは3〜30kHz、VLF-EMとVLF-MT。磁場変化を測る方式と電場磁場比を用いる方式を区別する周波数帯・分類である。\par

\noindent\textbf{5：条件付きで適切} TEMは時間応答、一次磁場なしで測りstatic shiftが生じにくい。送信遮断後の二次応答を観測し、MT電場の局所歪みとは異なる。turn-off過程・送信器残留応答などは別途扱い、生じにくいを完全ゼロへ強めない。\par

\paragraph{再利用する考え方} 近接場は高/低抵抗を記憶せずδと送受信距離から判断する。

\paragraph{共通理論} \hyperref[theory:exploration:nearfield]{第\ref*{theory:exploration:nearfield}節}, \hyperref[theory:exploration:archie]{第\ref*{theory:exploration:archie}節}, \hyperref[theory:exploration:fdtd]{第\ref*{theory:exploration:fdtd}節}

\par 参考：\cite{usgs_csamt2012,epa_ip,epa_tdem_archive,epa_borehole_resistivity}。

\paragraph{関連問題} \hyperref[q:h29-reg-applied-32]{H29-REG-APPLIED-32}

\par 難易度 3（校準初稿）、優先度未評価。

\subsection{平成28年度 問23}
\label{q:h28-reg-applied-23}

\noindent\textbf{原典：} sources/H28/17\_applied\_science.pdf、PDF物理9頁。

\textbf{公式正答：} 3（single）。

\par\textbf{検討未確定：} 原典は真空の誘電率をµと定義していることを画像確認。粘土の表面伝導と既存VLF局の条件を保持。

\paragraph{問題の要旨} Wenner式・ArchieとMT式の透磁率・誘電率を照合する。

\Needspace{4\baselineskip}\paragraph{確認した条件・記述（独自要約）}\mbox{}\par

\noindent\textbf{原典照合：} 29問21原典画像で否定、全五候補、数式・上付き単位・観測対象を照合。\par

\noindent\textbf{前提：} 見掛比抵抗ρa=\allowbreak{}2πaV/Iとρa=|Ex/Hy|\textsuperscript{2}/(ωµ)を使い、µは透磁率とする。\par

\paragraph{問われる核心} Wenner式・ArchieとMT式の透磁率・誘電率を照合する。

\paragraph{必要知識} 見掛比抵抗ρa=\allowbreak{}2πaV/Iとρa=|Ex/Hy|\textsuperscript{2}/(ωµ)を使い、µは透磁率とする。

\paragraph{考え方・正答を選ぶ理由（査読待ち初稿）} 3が不適切である。式のµは真空の透磁率µ0であって誘電率ε0ではない。式の形が正しくても原典の変数定義が誤る。Wenner式とMT式の次元を独立に照合する。

\Needspace{4\baselineskip}\paragraph{選択肢の検討（査読待ち初稿）}\mbox{}\par

\noindent\textbf{1：適切} Wennerのρa=\allowbreak{}2πaV/I。等間隔aの四電極を均質半空間の点電極として配置した幾何係数2πaである。\par

\noindent\textbf{2：条件付きで適切} 粘土の含水岩はArchie予測より低比抵抗。粘土表面の追加導電がある場合、孔隙水のみのモデルより伝導が増える。全試料の固定差ではない。\par

\noindent\textbf{3：不適切} MT式のµを真空の誘電率と定義する。E/HからΩを持つインピーダンスを作りωµで割る。µはH/mの透磁率でありF/mの誘電率ではない。\par

\noindent\textbf{4：適切} TEMはループ又は接地両端電流源で深部探査。磁気ループ方式だけでなく接地電流源の時間領域応答を測る構成もある。\par

\noindent\textbf{5：条件付きで適切} VLFは3〜30kHzの既存局で送信条件を制御できない。既存通信送信局を利用する受信探査では局の位置・周波数・運用を調査者が任意には設定できない。専用人工源の別方式を含めない。\par

\paragraph{再利用する考え方} 式の記号の名前・単位と励起方式を照合し、形だけで正しいと判定しない。

\paragraph{共通理論} \hyperref[theory:exploration:dcarrays]{第\ref*{theory:exploration:dcarrays}節}, \hyperref[theory:exploration:archie]{第\ref*{theory:exploration:archie}節}, \hyperref[theory:exploration:nearfield]{第\ref*{theory:exploration:nearfield}節}

\par 参考：\cite{epa_resistivity,usgs_csamt2012,epa_borehole_resistivity,epa_tdem_archive}。

\par 難易度 3（校準初稿）、優先度未評価。

\section{反射走時の双曲線・局所近似・配置}

\subsection{平成28年度 問21}
\label{q:h28-reg-applied-21}

\noindent\textbf{原典：} sources/H28/17\_applied\_science.pdf、PDF物理8頁。

\textbf{公式正答：} 1（single）。

\par\textbf{検討未確定：} 放物線は小オフセットの正当な局所二次近似となる。原文の近似に範囲がなく、公式1を通常の双曲線近似の知識として保持し条件不足を査読対象にする。20mも装置条件付きの目安。

\paragraph{問題の要旨} 反射走時の双曲線と放物線近似、ミラージ、微動、表面波を比較する。

\Needspace{4\baselineskip}\paragraph{確認した条件・記述（独自要約）}\mbox{}\par

\noindent\textbf{原典照合：} 五候補の順序、最適切と不適切の方向、波種、深度目安、曲線名、最長と最短を原典画像で照合。\par

\noindent\textbf{前提：} 単一水平反射境界の走時t\textsuperscript{2}=t0\textsuperscript{2}+x\textsuperscript{2}/V\textsuperscript{2}を基本モデルとする。\par

\paragraph{問われる核心} 反射走時の双曲線と放物線近似、ミラージ、微動、表面波を比較する。

\paragraph{必要知識} 単一水平反射境界の走時t\textsuperscript{2}=t0\textsuperscript{2}+x\textsuperscript{2}/V\textsuperscript{2}を基本モデルとする。

\paragraph{考え方・正答を選ぶ理由（査読待ち初稿）} 公式1は通常の反射走時を双曲線とすることを問う。しかしt=\allowbreak{}t0 sqrt(1+x\textsuperscript{2}/(V\textsuperscript{2}t0\textsuperscript{2}))を小xで展開すればt≈t0+x\textsuperscript{2}/(2V\textsuperscript{2}t0)という放物線近似が成立する。近似されるという原文を絶対的に数学の誤りと断定しない。

\Needspace{4\baselineskip}\paragraph{選択肢の検討（査読待ち初稿）}\mbox{}\par

\noindent\textbf{1：条件付きで適切} 水平二層の反射走時は放物線で近似される。正確な単一反射の走時は双曲線である。一方、小オフセットの二次近似は放物線であり、適用範囲を指定すれば近似として成立する。\par

\noindent\textbf{2：適切} 連続増速のミラージ構造では初動走時が曲線になりトモグラフィが適する。離散境界の折れ線モデルでなく深度ごとに屈折しながら曲がる経路となる。格子速度を更新する走時トモグラフィは連続的な速度変化も扱える。解像度は波線分布に依存する。\par

\noindent\textbf{3：適切} 屈折法は通常初動を読み解析する。直接波・臨界屈折波のうち最も早く到達する時刻を距離ごとに読む。反射波を同じ初動としない。\par

\noindent\textbf{4：適切} 微動アレイで自然・人為的な微振動からS波構造を推定する。複数点の位相・空間相関から表面波の分散を求め、主にS波速度構造を推定する。\par

\noindent\textbf{5：条件付きで適切} 表面波から20m程度までのS波分布を求める。通常の浅部能動測定の典型的対象である。最大波長と配列・震源が変われば探査深度も変わり、20mを絶対上限とはしない。\par

\paragraph{再利用する考え方} 曲線名を暗記するだけでなく走時式を展開し、近似範囲を確認する。

\paragraph{共通理論} \hyperref[theory:exploration:reflectiongeometry]{第\ref*{theory:exploration:reflectiongeometry}節}, \hyperref[theory:exploration:surfacewaves]{第\ref*{theory:exploration:surfacewaves}節}

\par 参考：\cite{epa_seismic_reflection,jica_refraction2015,moe_surface2004,epa_surfacewaves}。

\paragraph{関連問題} \hyperref[q:r05-reg-applied-32]{R05-REG-APPLIED-32}

\par 難易度 3（校準初稿）、優先度未評価。

\subsection{令和5年度 問32}
\label{q:r05-reg-applied-32}

\noindent\textbf{原典：} sources/R05/17\_applied\_science.pdf、PDF物理11頁。

\textbf{公式正答：} 4（single）。

\par\textbf{検討未確定：} 放物線は小オフセットの局所近似として成立する。公式4を一般走時式の比較として保持し範囲不足を残す。CMPによる相関雑音抑制もmoveoutの差に依存する。

\paragraph{問題の要旨} ミラージ、測線方向、S波屈折、反射曲線、相関雑音を比較する。

\Needspace{4\baselineskip}\paragraph{確認した条件・記述（独自要約）}\mbox{}\par

\noindent\textbf{原典照合：} 五候補の順序、最適切と不適切の方向、波種、深度目安、曲線名、最長と最短を原典画像で照合。\par

\noindent\textbf{前提：} CMPの通常の双曲線近似と小オフセットの局所二次近似を分ける。\par

\paragraph{問われる核心} ミラージ、測線方向、S波屈折、反射曲線、相関雑音を比較する。

\paragraph{必要知識} CMPの通常の双曲線近似と小オフセットの局所二次近似を分ける。

\paragraph{考え方・正答を選ぶ理由（査読待ち初稿）} 公式4である。水平成層の通常近似は双曲線であり、放物線を一般的走時式とする説明は不十分である。一方、小オフセットのTaylor二次近似は放物線となるため、原文の近似的という語に対する範囲不足を疑義として保持する。

\Needspace{4\baselineskip}\paragraph{選択肢の検討（査読待ち初稿）}\mbox{}\par

\noindent\textbf{1：適切} 連続的な深度増速構造をミラージュ構造と呼ぶ。速度層境界を持つ折れ線モデルでなく、曲がる波線の走時曲線になる連続勾配モデルである。\par

\noindent\textbf{2：条件付きで適切} 屈折測線を等高線とほぼ垂直で一定勾配に配置するのが望ましい。地質構造と地形を評価しやすい代表的な測線設計である。全地形・全目的でその配置が唯一又は可能とはいえない。\par

\noindent\textbf{3：適切} 土質のS波屈折を比較的平坦な地形で多く行う。起振と受振及び地形補正の制約を考えた通常の調査条件である。\par

\noindent\textbf{4：条件付きで適切} 多層CMPの反射走時は近似的に放物線となる。通常は有効速度を用いる双曲線で近似する。小オフセットなら双曲線の二次展開は放物線となるため、その条件を除外せず公式判断と分ける。\par

\noindent\textbf{5：条件付きで適切} CMP重合がランダムとコヒーレント雑音の抑制に貢献する。適切なNMOで反射波をそろえ、異なるmoveoutの多重反射等は重合で弱められる。反射と同じmoveoutの相関雑音まで必ず消えるわけではない。\par

\paragraph{再利用する考え方} 近似の次数と範囲、雑音のmoveoutを示して判断する。

\paragraph{共通理論} \hyperref[theory:exploration:reflectiongeometry]{第\ref*{theory:exploration:reflectiongeometry}節}, \hyperref[theory:exploration:refraction]{第\ref*{theory:exploration:refraction}節}, \hyperref[theory:exploration:stack]{第\ref*{theory:exploration:stack}節}

\par 参考：\cite{epa_seismic_reflection,jica_refraction2015,usgs_processing2000}。

\paragraph{関連問題} \hyperref[q:h28-reg-applied-21]{H28-REG-APPLIED-21}

\par 難易度 3（校準初稿）、優先度未評価。

\section{反射分解能・薄層分離と存在検出}

\subsection{平成26年度 問20}
\label{q:h26-reg-applied-20}

\noindent\textbf{原典：} sources/H26/17\_applied\_science.pdf、PDF物理8頁。

\textbf{公式正答：} 1（single）。

\par\textbf{検討未確定：} 原典1の反射法は浅層の目安を想定した公式判断と読めるが、調査条件を一律固定できない。2は検出という語で分解能を説明する省略があり、薄層の振幅検出を否定しない。

\paragraph{問題の要旨} 反射法の配置・薄層分解能・NMOと屈折法のブラインドを比較する。

\Needspace{4\baselineskip}\paragraph{確認した条件・記述（独自要約）}\mbox{}\par

\noindent\textbf{原典照合：} 五候補の順序、最適切と不適切の方向、波種、深度目安、曲線名、最長と最短を原典画像で照合。\par

\noindent\textbf{前提：} 最大オフセットの目安と必要条件、薄層の検出と二境界の分離を区別する。\par

\paragraph{問われる核心} 反射法の配置・薄層分解能・NMOと屈折法のブラインドを比較する。

\paragraph{必要知識} 最大オフセットの目安と必要条件、薄層の検出と二境界の分離を区別する。

\paragraph{考え方・正答を選ぶ理由（査読待ち初稿）} 1が不適切である。浅層反射の展開長の通常目安を探査深度程度とする資料に対し、二倍を必須としている。必要距離は目的と速度構造次第であり、二倍が絶対禁止という意味でもない。2の四分の一波長は二境界の分離の目安であり、薄層の存在検出はそれ以下でも可能となる。

\Needspace{4\baselineskip}\paragraph{選択肢の検討（査読待ち初稿）}\mbox{}\par

\noindent\textbf{1：不適切} 反射法の最大オフセットを探査深度の二倍程度にする必要がある。浅層反射の通常目安の倍数を取り違える。一律に二倍を必要とする原理はなく、波群分離・速度解析・調査目的から設計する。\par

\noindent\textbf{2：条件付きで適切} 反射法で波長の四分の一未満の地層の検出が難しい。通常の分解能目安はλ/4程度で上面と下面を別々の反射として分離する話である。振幅変化で存在を検出できることと区別し、薄層は全て不可視とはしない。\par

\noindent\textbf{3：適切} CMP記録からNMO用の重合速度を速度解析する。異なるオフセットでの走時差をそろえる速度を推定する。重合速度を無条件に真の区間速度と同一視しない。\par

\noindent\textbf{4：適切} 深部へ増速する層構造でも薄い中間層が初動に現れない。臨界屈折波が生じても、直接波や深い層を通る波より早くなる距離範囲がなければ初動に現れない。低速度逆転だけが見えない原因ではない。\par

\noindent\textbf{5：条件付きで適切} 屈折法で全走時が一致することを確認する。相反的な同じ測線で起振点を往復した端点間走時の整合などを確認する。異なる受振点の全時刻が同じになるという意味ではなく、走時読み取り・経路の条件を要する。\par

\paragraph{再利用する考え方} 配置の数値、分離能力、検出の有無、速度の種類を別々に判断する。

\paragraph{共通理論} \hyperref[theory:exploration:reflectiongeometry]{第\ref*{theory:exploration:reflectiongeometry}節}, \hyperref[theory:exploration:resolution]{第\ref*{theory:exploration:resolution}節}, \hyperref[theory:exploration:refraction]{第\ref*{theory:exploration:refraction}節}

\par 参考：\cite{jica_refraction2015,usgs_refraction2013,epa_surfacewaves,xia_rayleigh1999,coastal_reflection2016,seg_vertical_resolution,usgs_processing2000}。

\paragraph{関連問題} \hyperref[q:r03-reg-applied-33]{R03-REG-APPLIED-33}

\par 難易度 3（校準初稿）、優先度未評価。

\subsection{令和3年度 問33}
\label{q:r03-reg-applied-33}

\noindent\textbf{原典：} sources/R03/17\_applied\_science.pdf、PDF物理12頁。

\textbf{公式正答：} 1（single）。

\paragraph{問題の要旨} 反射の分解能を走時切片と区別する。

\Needspace{4\baselineskip}\paragraph{確認した条件・記述（独自要約）}\mbox{}\par

\noindent\textbf{原典照合：} 五候補の順序、最適切と不適切の方向、波種、深度目安、曲線名、最長と最短を原典画像で照合。\par

\noindent\textbf{前提：} 薄層の上下面の分離、横方向の反射寄与範囲、入力waveletを区別する。\par

\paragraph{問われる核心} 反射の分解能を走時切片と区別する。

\paragraph{必要知識} 薄層の上下面の分離、横方向の反射寄与範囲、入力waveletを区別する。

\paragraph{考え方・正答を選ぶ理由（査読待ち初稿）} 1が不適切である。intercept timeは屈折走時の直線の切片であり、反射の二イベントの分離や横方向の寄与範囲を表す基準ではない。

\Needspace{4\baselineskip}\paragraph{選択肢の検討（査読待ち初稿）}\mbox{}\par

\noindent\textbf{1：不適切} インターセプトタイム。屈折波の走時直線を距離ゼロへ外挿した切片である。反射分解能の直接的な基準と異なる。\par

\noindent\textbf{2：適切} レイリー基準。近接する二つの応答を分離できるかという分解能の基準に関係する。表面波としてのRayleighと用語の役割を分ける。\par

\noindent\textbf{3：適切} ウェーブレット。反射の応答形状と有効帯域が近いイベントを分離できる程度に影響する。\par

\noindent\textbf{4：適切} リッカー基準。反射waveletの重なりを用いる分解能の基準に関係する。具体的閾値を基準名から一律数値へ補わない。\par

\noindent\textbf{5：適切} フレネルゾーン。同程度の位相で反射へ寄与する面上の範囲であり、横方向の分解能に関係する。migration前後と開口の条件を区別する。\par

\paragraph{再利用する考え方} 同じ人名の用語でも波の種類、wavelet、分離基準を混同しない。

\paragraph{共通理論} \hyperref[theory:exploration:resolution]{第\ref*{theory:exploration:resolution}節}

\par 参考：\cite{seg_vertical_resolution,seg_lateral_resolution,usgs_refraction2013}。

\paragraph{関連問題} \hyperref[q:h26-reg-applied-20]{H26-REG-APPLIED-20}

\par 難易度 2（校準初稿）、優先度未評価。

\section{屈折法の地山評価と主用途}

\subsection{令和2年度 問32}
\label{q:r02-reg-applied-32}

\noindent\textbf{原典：} sources/R02/17\_applied\_science.pdf、PDF物理11頁。

\textbf{公式正答：} 2（single）。

\par\textbf{検討未確定：} 屈折法にはS波もあり、強震動に寄与する地下情報を得る可能性はある。通常のP波地山調査の主用途という比較を明示し、用途の絶対排除をしない。

\paragraph{問題の要旨} 屈折法の主用途と強震動のS波モデルを区別する。

\Needspace{4\baselineskip}\paragraph{確認した条件・記述（独自要約）}\mbox{}\par

\noindent\textbf{原典照合：} 五候補の順序、最適切と不適切の方向、波種、深度目安、曲線名、最長と最短を原典画像で照合。\par

\noindent\textbf{前提：} 通常のP波屈折による地山の硬軟・割れ目・風化等の評価を主用途として比較する。\par

\paragraph{問われる核心} 屈折法の主用途と強震動のS波モデルを区別する。

\paragraph{必要知識} 通常のP波屈折による地山の硬軟・割れ目・風化等の評価を主用途として比較する。

\paragraph{考え方・正答を選ぶ理由（査読待ち初稿）} 公式2である。強震動・耐震評価では主にS波構造や減衰・震源条件が必要であり、通常の地山P波屈折をその主たる方法としない。他の四項目は岩盤・土構造物の地質条件の評価に結び付く。ただしS波屈折や複数手法の統合が耐震評価に役立つ可能性まで否定しない。

\Needspace{4\baselineskip}\paragraph{選択肢の検討（査読待ち初稿）}\mbox{}\par

\noindent\textbf{1：適切} ダムの基礎岩盤の評価。風化・破砕・硬軟の構造を把握する通常の地山評価に関係する。\par

\noindent\textbf{2：条件付きで適切} 強震動予測や耐震設計。P波屈折の代表的な主用途の比較では公式の選択対象である。S波屈折等の応用と地下モデルの補助情報としての役割はあり、全く利用できないという意味ではない。\par

\noindent\textbf{3：適切} のり面勾配等の土構造物の基礎調査。風化帯や基盤等の地質評価を補助する。速度だけで設計勾配を機械的に決定しない。\par

\noindent\textbf{4：適切} 不安定土塊の工学的判断。土塊・基盤の速度構造が地質状態の判断を助ける。力学強度や地下水等の他情報も要する。\par

\noindent\textbf{5：適切} トンネル沿いの地山評価。岩盤の風化・破砕帯等の分布を速度情報から把握する代表的用途である。\par

\paragraph{再利用する考え方} 主としてという比較を、利用可能性が一切ないという全称へ変えない。

\paragraph{共通理論} \hyperref[theory:exploration:waveproperties]{第\ref*{theory:exploration:waveproperties}節}, \hyperref[theory:exploration:surfacewaves]{第\ref*{theory:exploration:surfacewaves}節}

\par 参考：\cite{segj_integrated2017,jica_refraction2015,japansea_spac2017}。

\par 難易度 3（校準初稿）、優先度未評価。

\section{屈折初動のはぎとり・相反走時・連続勾配}

\subsection{平成23年度 問32}
\label{q:h23-reg-applied-32}

\noindent\textbf{原典：} sources/H23/17\_applied\_science.pdf、PDF物理11頁。

\textbf{公式正答：} 3（single）。

\par\textbf{検討未確定：} オフセットと深度の対応は目安であり一律の必要条件ではない。全走時の一致は相反的往復観測の整合であり、全測点の同時到達を意味しない。

\paragraph{問題の要旨} 表面波・浅層反射・はぎとり・全走時・ミラージ構造を比較する。

\Needspace{4\baselineskip}\paragraph{確認した条件・記述（独自要約）}\mbox{}\par

\noindent\textbf{原典照合：} 五候補の順序、最適切と不適切の方向、波種、深度目安、曲線名、最長と最短を原典画像で照合。\par

\noindent\textbf{前提：} 通常の屈折法の初動走時、浅層反射の配置目安、波長依存の表面波分散を考える。\par

\paragraph{問われる核心} 表面波・浅層反射・はぎとり・全走時・ミラージ構造を比較する。

\paragraph{必要知識} 通常の屈折法の初動走時、浅層反射の配置目安、波長依存の表面波分散を考える。

\paragraph{考え方・正答を選ぶ理由（査読待ち初稿）} 3が不適切である。はぎとり法は屈折波と直接波の初動走時の解析であり、反射波から同じ手法で境界を求める説明は対応が誤っている。全走時一致は往復の整合を意味し、各受振点の時刻が同じという命題ではない。

\Needspace{4\baselineskip}\paragraph{選択肢の検討（査読待ち初稿）}\mbox{}\par

\noindent\textbf{1：適切} レイリー波の波長別位相速度からS波速度分布を推定する。地盤の深度方向の不均質が分散曲線へ反映され、主にS波速度モデルを推定できる。P波速度・密度・層厚にも依存し、S波だけの関数ではない。\par

\noindent\textbf{2：条件付きで適切} 浅層反射で最大オフセットを探査深度程度に設定する。浅層反射の一般的な展開目安として読む。速度解析・近接雑音・反射角・地形・重合数で調整し、全ての計測に必要な厳密等式ではない。\par

\noindent\textbf{3：不適切} 反射波の走時をはぎとり法で解析して境界深度を求める。萩原のはぎとり法は直接波と屈折波の初動走時から層構造を求める手法である。反射波の走時を同じ手法の入力とする説明は誤りである。\par

\noindent\textbf{4：条件付きで適切} 屈折法で全走時が一致することを確認する。相反的な同じ測線で起振点を往復した端点間走時の整合などを確認する。異なる受振点の全時刻が同じになるという意味ではなく、走時読み取り・経路の条件を要する。\par

\noindent\textbf{5：適切} 連続増速のミラージ構造では初動走時が曲線になりトモグラフィが適する。離散境界の折れ線モデルでなく深度ごとに屈折しながら曲がる経路となる。格子速度を更新する走時トモグラフィは連続的な速度変化も扱える。解像度は波線分布に依存する。\par

\paragraph{再利用する考え方} 反射と屈折を、入力する波と走時解析の前提から分ける。

\paragraph{共通理論} \hyperref[theory:exploration:refraction]{第\ref*{theory:exploration:refraction}節}, \hyperref[theory:exploration:reflectiongeometry]{第\ref*{theory:exploration:reflectiongeometry}節}, \hyperref[theory:exploration:surfacewaves]{第\ref*{theory:exploration:surfacewaves}節}

\par 参考：\cite{jica_refraction2015,usgs_refraction2013,epa_surfacewaves,xia_rayleigh1999,coastal_reflection2016}。

\paragraph{関連問題} \hyperref[q:h29-reg-applied-34]{H29-REG-APPLIED-34}

\par 難易度 3（校準初稿）、優先度未評価。

\subsection{平成29年度 問34}
\label{q:h29-reg-applied-34}

\noindent\textbf{原典：} sources/H29/17\_applied\_science.pdf、PDF物理12頁。

\textbf{公式正答：} 3（single）。

\par\textbf{検討未確定：} オフセットと深度の対応は目安であり一律の必要条件ではない。全走時の一致は相反的往復観測の整合であり、全測点の同時到達を意味しない。

\paragraph{問題の要旨} 表面波・浅層反射・はぎとり・全走時・ミラージ構造を比較する。

\Needspace{4\baselineskip}\paragraph{確認した条件・記述（独自要約）}\mbox{}\par

\noindent\textbf{原典照合：} 五候補の順序、最適切と不適切の方向、波種、深度目安、曲線名、最長と最短を原典画像で照合。\par

\noindent\textbf{前提：} 通常の屈折法の初動走時、浅層反射の配置目安、波長依存の表面波分散を考える。\par

\paragraph{問われる核心} 表面波・浅層反射・はぎとり・全走時・ミラージ構造を比較する。

\paragraph{必要知識} 通常の屈折法の初動走時、浅層反射の配置目安、波長依存の表面波分散を考える。

\paragraph{考え方・正答を選ぶ理由（査読待ち初稿）} 3が不適切である。はぎとり法は屈折波と直接波の初動走時の解析であり、反射波から同じ手法で境界を求める説明は対応が誤っている。全走時一致は往復の整合を意味し、各受振点の時刻が同じという命題ではない。

\Needspace{4\baselineskip}\paragraph{選択肢の検討（査読待ち初稿）}\mbox{}\par

\noindent\textbf{1：適切} レイリー波の波長別位相速度からS波速度分布を推定する。地盤の深度方向の不均質が分散曲線へ反映され、主にS波速度モデルを推定できる。P波速度・密度・層厚にも依存し、S波だけの関数ではない。\par

\noindent\textbf{2：条件付きで適切} 浅層反射で最大オフセットを探査深度程度に設定する。浅層反射の一般的な展開目安として読む。速度解析・近接雑音・反射角・地形・重合数で調整し、全ての計測に必要な厳密等式ではない。\par

\noindent\textbf{3：不適切} 反射波の走時をはぎとり法で解析して境界深度を求める。萩原のはぎとり法は直接波と屈折波の初動走時から層構造を求める手法である。反射波の走時を同じ手法の入力とする説明は誤りである。\par

\noindent\textbf{4：条件付きで適切} 屈折法で全走時が一致することを確認する。相反的な同じ測線で起振点を往復した端点間走時の整合などを確認する。異なる受振点の全時刻が同じになるという意味ではなく、走時読み取り・経路の条件を要する。\par

\noindent\textbf{5：適切} 連続増速のミラージ構造では初動走時が曲線になりトモグラフィが適する。離散境界の折れ線モデルでなく深度ごとに屈折しながら曲がる経路となる。格子速度を更新する走時トモグラフィは連続的な速度変化も扱える。解像度は波線分布に依存する。\par

\paragraph{再利用する考え方} 反射と屈折を、入力する波と走時解析の前提から分ける。

\paragraph{共通理論} \hyperref[theory:exploration:refraction]{第\ref*{theory:exploration:refraction}節}, \hyperref[theory:exploration:reflectiongeometry]{第\ref*{theory:exploration:reflectiongeometry}節}, \hyperref[theory:exploration:surfacewaves]{第\ref*{theory:exploration:surfacewaves}節}

\par 参考：\cite{jica_refraction2015,usgs_refraction2013,epa_surfacewaves,xia_rayleigh1999,coastal_reflection2016}。

\paragraph{関連問題} \hyperref[q:h23-reg-applied-32]{H23-REG-APPLIED-32}

\par 難易度 3（校準初稿）、優先度未評価。

\section{地震速度境界・地質と地盤状態}

\subsection{平成25年度 問32}
\label{q:h25-reg-applied-32}

\noindent\textbf{原典：} sources/H25/17\_applied\_science.pdf、PDF物理12頁。

\textbf{公式正答：} 3（single）。

\par\textbf{検討未確定：} P>Sは安定な等方弾性体の条件、増速と初動のみは通常の屈折法の枠組みである。異方性や他の解析を無条件に排除しない。

\paragraph{問題の要旨} 屈折法の増速前提と地下水・風化・波種・初動を区別する。

\Needspace{4\baselineskip}\paragraph{確認した条件・記述（独自要約）}\mbox{}\par

\noindent\textbf{原典照合：} 五候補の順序、最適切と不適切の方向、波種、深度目安、曲線名、最長と最短を原典画像で照合。\par

\noindent\textbf{前提：} 通常の地盤の等方弾性を想定し、臨界屈折の層状解釈とP波の飽和依存を区別する。\par

\paragraph{問われる核心} 屈折法の増速前提と地下水・風化・波種・初動を区別する。

\paragraph{必要知識} 通常の地盤の等方弾性を想定し、臨界屈折の層状解釈とP波の飽和依存を区別する。

\paragraph{考え方・正答を選ぶ理由（査読待ち初稿）} 3が不適切である。空隙を空気又は水が満たす状態は圧縮性等へ影響し、P波速度は地下水の有無・飽和状態に左右され得る。速度だけで岩種を決めず、同じ岩種でも割れ目と風化の程度を考える。

\Needspace{4\baselineskip}\paragraph{選択肢の検討（査読待ち初稿）}\mbox{}\par

\noindent\textbf{1：条件付きで適切} 通常の屈折法で深部ほど速度が高いと仮定する。臨界屈折の層状モデルの前提である。あらゆる地下構造が必ず増速するという自然法則ではない。\par

\noindent\textbf{2：適切} 同じ岩種でも浅部速度の違いが風化・変質・割れ目の違いを反映する。弾性係数、接触状態、空隙等の変化が速度へ影響する。同じ岩種なら一定速度とは限らない。\par

\noindent\textbf{3：不適切} 土質のP波速度は地下水の有無に左右されない。空隙流体と飽和状態が圧縮性に影響するため無関係という断定は誤りである。\par

\noindent\textbf{4：条件付きで適切} P波速度は常にS波速度より大きい。安定な等方弾性体ではVP\textsuperscript{2}=(K+4G/3)/ρ、VS\textsuperscript{2}=G/ρよりVP>VSである。異方性のモードや流体のS波不存在まで常に二つの速度を比較できるわけではない。\par

\noindent\textbf{5：条件付きで適切} 解析に初動走時のみを用いる。通常の屈折初動解析の説明である。反射・全波形等を含める別の解析まで禁止する意味ではない。\par

\paragraph{再利用する考え方} 波種、弾性係数、空隙流体を結び、速度を岩種の識別子にしない。

\paragraph{共通理論} \hyperref[theory:exploration:refraction]{第\ref*{theory:exploration:refraction}節}, \hyperref[theory:exploration:waveproperties]{第\ref*{theory:exploration:waveproperties}節}

\par 参考：\cite{usgs_refraction2013,segj_integrated2017}。

\paragraph{関連問題} \hyperref[q:r03-reg-applied-34]{R03-REG-APPLIED-34}

\par 難易度 3（校準初稿）、優先度未評価。

\subsection{令和3年度 問34}
\label{q:r03-reg-applied-34}

\noindent\textbf{原典：} sources/R03/17\_applied\_science.pdf、PDF物理13頁。

\textbf{公式正答：} 2（single）。

\par\textbf{検討未確定：} S波屈折は平坦な土質調査という代表的運用の比較であり、急峻地形で絶対実施できないと断定しない。速度からの地質推定は複数解を持つ。

\paragraph{問題の要旨} 速度境界と地質境界、S波屈折の地形条件と用途を比較する。

\Needspace{4\baselineskip}\paragraph{確認した条件・記述（独自要約）}\mbox{}\par

\noindent\textbf{原典照合：} 五候補の順序、最適切と不適切の方向、波種、深度目安、曲線名、最長と最短を原典画像で照合。\par

\noindent\textbf{前提：} 通常の土質地盤のS波屈折は比較的平坦な配置で測ることが多い。地形により走時解釈や起振・受振方向が難しくなる。\par

\paragraph{問われる核心} 速度境界と地質境界、S波屈折の地形条件と用途を比較する。

\paragraph{必要知識} 通常の土質地盤のS波屈折は比較的平坦な配置で測ることが多い。地形により走時解釈や起振・受振方向が難しくなる。

\paragraph{考え方・正答を選ぶ理由（査読待ち初稿）} 2が不適切である。土質のS波屈折を急峻地形で多く行うという説明を、典型的な比較的平坦な条件と取り違えている。地形を克服する特別な測定が不可能という意味ではない。

\Needspace{4\baselineskip}\paragraph{選択肢の検討（査読待ち初稿）}\mbox{}\par

\noindent\textbf{1：適切} 屈折法の速度境界は地質境界と必ずしも一致しない。速度は岩種だけでなく風化・割れ目・飽和等に依存する。同じ岩種内でも速度が変わり、異なる地質が近い速度を持つ場合もある。\par

\noindent\textbf{2：不適切} 土質のS波屈折を比較的急峻な地形で多く行う。典型的なS波の土質調査の地形条件は比較的平坦である。急峻が通常条件という説明は逆であり、起振・受振・経路の制約が増す。\par

\noindent\textbf{3：適切} トンネル・ダム・のり面等の基礎調査へ屈折法を用いる。通常のP波初動を使って地山の速度構造を求め、工学的地質評価の一情報とする。\par

\noindent\textbf{4：条件付きで適切} 速度構造から硬軟・割れ目・風化・破砕帯等を推定する。物性と地質状態の相関を利用する推定である。速度だけから一意に種類や強度を決めず、露頭・ボーリング等で拘束する。\par

\noindent\textbf{5：適切} 速度層は岩種・膠結・風化・変質・亀裂の複合で決まる。速度への寄与が複数あるため、速度層を岩種ごとのラベルと同一視しない。\par

\paragraph{再利用する考え方} 地形の典型条件と技術的な実施可能性を分け、速度の多要因性を保つ。

\paragraph{共通理論} \hyperref[theory:exploration:waveproperties]{第\ref*{theory:exploration:waveproperties}節}, \hyperref[theory:exploration:refraction]{第\ref*{theory:exploration:refraction}節}

\par 参考：\cite{jica_refraction2015,segj_integrated2017,usgs_refraction2013}。

\paragraph{関連問題} \hyperref[q:h25-reg-applied-32]{H25-REG-APPLIED-32}, \hyperref[q:r06-reg-applied-32]{R06-REG-APPLIED-32}

\par 難易度 3（校準初稿）、優先度未評価。

\subsection{令和6年度 問32}
\label{q:r06-reg-applied-32}

\noindent\textbf{原典：} sources/R06/17\_applied\_science.pdf、PDF物理11頁。

\textbf{公式正答：} 1（single）。

\par\textbf{検討未確定：} オフセット深度の対応と双曲線は通常条件の目安・近似である。受振器群の効果は方向・間隔・波長に依存する。

\paragraph{問題の要旨} 速度境界、連続勾配、浅層展開、反射双曲線、受振器群を比較する。

\Needspace{4\baselineskip}\paragraph{確認した条件・記述（独自要約）}\mbox{}\par

\noindent\textbf{原典照合：} 五候補の順序、最適切と不適切の方向、波種、深度目安、曲線名、最長と最短を原典画像で照合。\par

\noindent\textbf{前提：} 地震波速度は岩種だけでなく、割れ目・風化・飽和などの状態にも依存する。\par

\paragraph{問われる核心} 速度境界、連続勾配、浅層展開、反射双曲線、受振器群を比較する。

\paragraph{必要知識} 地震波速度は岩種だけでなく、割れ目・風化・飽和などの状態にも依存する。

\paragraph{考え方・正答を選ぶ理由（査読待ち初稿）} 1が不適切である。屈折法の速度境界と地質学の境界がよく一致するという一般化はできない。同じ岩種内の風化境界が速度を変え、異なる岩種間で速度差が小さい場合もある。

\Needspace{4\baselineskip}\paragraph{選択肢の検討（査読待ち初稿）}\mbox{}\par

\noindent\textbf{1：不適切} 屈折法の層境界は地質境界とよく一致する。速度差で定まる境界を地質分類の境界と同一視している。岩種以外の要因が速度へ影響するため一般的な一致を保証しない。\par

\noindent\textbf{2：適切} 連続増速のミラージ構造では初動走時が曲線になりトモグラフィが適する。離散境界の折れ線モデルでなく深度ごとに屈折しながら曲がる経路となる。格子速度を更新する走時トモグラフィは連続的な速度変化も扱える。解像度は波線分布に依存する。\par

\noindent\textbf{3：条件付きで適切} 浅層反射で最大オフセットを探査深度程度に設定する。浅層反射の一般的な展開目安として読む。速度解析・近接雑音・反射角・地形・重合数で調整し、全ての計測に必要な厳密等式ではない。\par

\noindent\textbf{4：条件付きで適切} 水平多層CMPの走時は近似的に双曲線となる。通常のNMO近似で有効速度とゼロオフセット走時を用いる。大オフセットや強い不均質・異方性では非双曲線成分が重要になる。\par

\noindent\textbf{5：条件付きで適切} 受振器群を結合し表面波等を抑え反射を強調する。配列の空間応答により波群を選別し独立雑音も平均化する。方向・間隔が対象に適していなければ必要信号を弱める場合がある。\par

\paragraph{再利用する考え方} 速度構造と地質構造、近似モデルと実地の制約を別々に評価する。

\paragraph{共通理論} \hyperref[theory:exploration:waveproperties]{第\ref*{theory:exploration:waveproperties}節}, \hyperref[theory:exploration:reflectiongeometry]{第\ref*{theory:exploration:reflectiongeometry}節}, \hyperref[theory:exploration:stack]{第\ref*{theory:exploration:stack}節}

\par 参考：\cite{segj_integrated2017,jica_refraction2015,usgs_processing2000,epa_seismic_reflection}。

\paragraph{関連問題} \hyperref[q:r03-reg-applied-34]{R03-REG-APPLIED-34}

\par 難易度 3（校準初稿）、優先度未評価。

\section{地震探査の走時・画像化用語}

\subsection{平成24年度 問23}
\label{q:h24-reg-applied-23}

\noindent\textbf{原典：} sources/H24/17\_applied\_science.pdf、PDF物理7頁。

\textbf{公式正答：} 4（single）。

\paragraph{問題の要旨} 地震探査の走時・反射用語と重力補正を区別する。

\Needspace{4\baselineskip}\paragraph{確認した条件・記述（独自要約）}\mbox{}\par

\noindent\textbf{原典照合：} 否定の対象、指数、周波数と周期、五つの候補の列順を原典画像で照合。\par

\noindent\textbf{前提：} 地震探査は波の走時・波形、重力探査は加速度の異常を測る。\par

\paragraph{問われる核心} 地震探査の走時・反射用語と重力補正を区別する。

\paragraph{必要知識} 地震探査は波の走時・波形、重力探査は加速度の異常を測る。

\paragraph{考え方・正答を選ぶ理由（査読待ち初稿）} 4のブーゲー補正が不適切である。地震探査に関係する四つは、屈折波の走時切片、初動に現れにくい層、反射を支配するインピーダンス、地形や表層に由来する静的走時ずれに対応する。

\Needspace{4\baselineskip}\paragraph{選択肢の検討（査読待ち初稿）}\mbox{}\par

\noindent\textbf{1：適切} 原点走時。屈折波の走時直線を距離ゼロへ外挿した切片である。直線の傾きと合わせて層厚等を推定する。実測した距離ゼロの初動時刻とは同一視しない。\par

\noindent\textbf{2：適切} ブラインド層。屈折法の初動走時に現れにくい層である。薄い中間層が初動を与えない場合や低速度層を区別し、用語の範囲は資料により異なる。\par

\noindent\textbf{3：適切} 音響インピーダンス。垂直入射の簡単なモデルではZ=ρVであり、境界両側の差が反射振幅を支配する。重力で用いる密度そのものだけとは異なる。\par

\noindent\textbf{4：不適切} ブーゲー補正。測点と基準面の間の質量が重力へ及ぼす効果を補正する重力探査の処理である。地震波走時に対する静補正ではない。\par

\noindent\textbf{5：適切} 静補正。地形や表層低速度層などによる走時ずれを一定の時間シフトとしてそろえる処理である。重力の地形・Bouguer補正と観測量が異なる。\par

\paragraph{再利用する考え方} 処理名だけでなく何を補正し、単位は何かを問う。

\paragraph{共通理論} \hyperref[theory:exploration:seismicterms]{第\ref*{theory:exploration:seismicterms}節}, \hyperref[theory:gravity:corrections]{第\ref*{theory:gravity:corrections}節}

\par 参考：\cite{usgs_refraction2013,epa_seismic_reflection,usgs_static2004,usgs_processing2000}。

\paragraph{関連問題} \hyperref[q:h30-reg-applied-33]{H30-REG-APPLIED-33}

\par 難易度 2（校準初稿）、優先度未評価。

\subsection{平成30年度 問33}
\label{q:h30-reg-applied-33}

\noindent\textbf{原典：} sources/H30/17\_applied\_science.pdf、PDF物理11頁。

\textbf{公式正答：} 5（single）。

\paragraph{問題の要旨} 地震探査の走時・反射用語と重力補正を区別する。

\Needspace{4\baselineskip}\paragraph{確認した条件・記述（独自要約）}\mbox{}\par

\noindent\textbf{原典照合：} 否定の対象、指数、周波数と周期、五つの候補の列順を原典画像で照合。\par

\noindent\textbf{前提：} 地震探査は波の走時・波形、重力探査は加速度の異常を測る。\par

\paragraph{問われる核心} 地震探査の走時・反射用語と重力補正を区別する。

\paragraph{必要知識} 地震探査は波の走時・波形、重力探査は加速度の異常を測る。

\paragraph{考え方・正答を選ぶ理由（査読待ち初稿）} 5のブーゲー補正が不適切である。地震探査に関係する四つは、屈折波の走時切片、初動に現れにくい層、反射を支配するインピーダンス、地形や表層に由来する静的走時ずれに対応する。

\Needspace{4\baselineskip}\paragraph{選択肢の検討（査読待ち初稿）}\mbox{}\par

\noindent\textbf{1：適切} 原点走時。屈折波の走時直線を距離ゼロへ外挿した切片である。直線の傾きと合わせて層厚等を推定する。実測した距離ゼロの初動時刻とは同一視しない。\par

\noindent\textbf{2：適切} ブラインド層。屈折法の初動走時に現れにくい層である。薄い中間層が初動を与えない場合や低速度層を区別し、用語の範囲は資料により異なる。\par

\noindent\textbf{3：適切} 音響インピーダンス。垂直入射の簡単なモデルではZ=ρVであり、境界両側の差が反射振幅を支配する。重力で用いる密度そのものだけとは異なる。\par

\noindent\textbf{4：適切} 静補正。地形や表層低速度層などによる走時ずれを一定の時間シフトとしてそろえる処理である。重力の地形・Bouguer補正と観測量が異なる。\par

\noindent\textbf{5：不適切} ブーゲー補正。測点と基準面の間の質量が重力へ及ぼす効果を補正する重力探査の処理である。地震波走時に対する静補正ではない。\par

\paragraph{再利用する考え方} 処理名だけでなく何を補正し、単位は何かを問う。

\paragraph{共通理論} \hyperref[theory:exploration:seismicterms]{第\ref*{theory:exploration:seismicterms}節}, \hyperref[theory:gravity:corrections]{第\ref*{theory:gravity:corrections}節}

\par 参考：\cite{usgs_refraction2013,epa_seismic_reflection,usgs_static2004,usgs_processing2000}。

\paragraph{関連問題} \hyperref[q:h24-reg-applied-23]{H24-REG-APPLIED-23}

\par 難易度 2（校準初稿）、優先度未評価。

\subsection{令和元年度 問34}
\label{q:r01-reg-applied-34}

\noindent\textbf{原典：} sources/R01/17\_applied\_science.pdf、PDF物理12頁。

\textbf{公式正答：} 4（single）。

\par\textbf{独立解答と公式正答に相違あり：} 独立解答 5。初回独立5対公式4を保持し資料確認後の訂正4を別記する。rigid-boom論文は検索収録の当該原理を確認し、PDF取得は403及びHTML challengeで未達である。日本語表記と装置体系の最終照合を査読対象とする。

\paragraph{問題の要旨} 静補正・深度走時・CMP・リジッドブーム・下方接続を区別する。

\Needspace{4\baselineskip}\paragraph{確認した条件・記述（独自要約）}\mbox{}\par

\noindent\textbf{原典照合：} 否定の対象、指数、周波数と周期、五つの候補の列順を原典画像で照合。\par

\noindent\textbf{前提：} 原典の4はリジットブームシステム、5は下方接続と表記される。下方接続を重力・磁気にしか使わない手法と限定しない。\par

\paragraph{問われる核心} 静補正・深度走時・CMP・リジッドブーム・下方接続を区別する。

\paragraph{必要知識} 原典の4はリジットブームシステム、5は下方接続と表記される。下方接続を重力・磁気にしか使わない手法と限定しない。

\paragraph{考え方・正答を選ぶ理由（査読待ち初稿）} 公式4である。初回5は下方接続をポテンシャル場だけに限定した誤認であった。Stanfordの波動下方接続とmigrationの説明を確認すると、地下へ波動場を外挿する処理も地震探査に関係する。リジッドブームEMIは送受信コイルを用いる電磁誘導計測の配置であり、地震データ処理の他四項目と異なる。

\Needspace{4\baselineskip}\paragraph{選択肢の検討（査読待ち初稿）}\mbox{}\par

\noindent\textbf{1：適切} 静補正。地形や表層低速度層などによる走時ずれを一定の時間シフトとしてそろえる処理である。重力の地形・Bouguer補正と観測量が異なる。\par

\noindent\textbf{2：適切} 深度走時。速度モデルと対応する深度での伝播時間を結ぶ量であり、地震探査の時間・深度対応に用いる。\par

\noindent\textbf{3：適切} CMP重合。共通中点の反射記録を集め、走時差を補正して重ねる反射法の処理である。強い傾斜や不均質では中点が同じ反射点を意味するとは限らない。\par

\noindent\textbf{4：不適切} リジットブームシステム。原典表記を保存する。rigid-boom EMIの一次論文では一次磁場と導体の二次磁場を送受信コイルで扱う電磁計測として記述される。地震波の走時処理ではない。\par

\noindent\textbf{5：適切} 下方接続。地震波動場を深部へ外挿してmigration画像を得る方法に関係する。重力・磁気のポテンシャル場にも下方接続があるため、名前の共通性だけで地震分野から除外しない。\par

\paragraph{再利用する考え方} 同じ数学的処理が複数の物理分野で使われることを許し、計測装置の物理原理を確認する。

\paragraph{共通理論} \hyperref[theory:exploration:seismicterms]{第\ref*{theory:exploration:seismicterms}節}, \hyperref[theory:exploration:continuation]{第\ref*{theory:exploration:continuation}節}, \hyperref[theory:exploration:electromagnetic]{第\ref*{theory:exploration:electromagnetic}節}

\par 参考：\cite{usgs_refraction2013,epa_seismic_reflection,usgs_static2004,usgs_processing2000,stanford_downward2001,rigid_boom2019}。

\paragraph{関連問題} \hyperref[q:r04-reg-applied-32]{R04-REG-APPLIED-32}

\par 難易度 3（校準初稿）、優先度未評価。

\subsection{令和4年度 問32}
\label{q:r04-reg-applied-32}

\noindent\textbf{原典：} sources/R04/17\_applied\_science.pdf、PDF物理11頁。

\textbf{公式正答：} 3（single）。

\paragraph{問題の要旨} 深度走時・migration・静補正・CMPとブーゲー補正を区別する。

\Needspace{4\baselineskip}\paragraph{確認した条件・記述（独自要約）}\mbox{}\par

\noindent\textbf{原典照合：} 否定の対象、指数、周波数と周期、五つの候補の列順を原典画像で照合。\par

\noindent\textbf{前提：} 走時・反射波形の処理と重力の質量補正を分ける。\par

\paragraph{問われる核心} 深度走時・migration・静補正・CMPとブーゲー補正を区別する。

\paragraph{必要知識} 走時・反射波形の処理と重力の質量補正を分ける。

\paragraph{考え方・正答を選ぶ理由（査読待ち初稿）} 3のブーゲー補正が不適切である。migrationは反射・回折記録を地下の反射位置へ戻す処理であり、波動場を下方接続する方法もある。

\Needspace{4\baselineskip}\paragraph{選択肢の検討（査読待ち初稿）}\mbox{}\par

\noindent\textbf{1：適切} 深度走時。速度モデルと深度から伝播時間を対応させる地震探査の量である。\par

\noindent\textbf{2：適切} マイグレーション。傾斜面や回折の見掛け位置を地下の適切な位置へ再配置する画像化処理である。観測時間を一律定数で換算するだけの処理ではない。\par

\noindent\textbf{3：不適切} ブーゲー補正。測点と基準面の間の質量が重力へ及ぼす効果を補正する重力探査の処理である。地震波走時に対する静補正ではない。\par

\noindent\textbf{4：適切} 静補正。地形や表層低速度層などによる走時ずれを一定の時間シフトとしてそろえる処理である。重力の地形・Bouguer補正と観測量が異なる。\par

\noindent\textbf{5：適切} CMP重合。共通中点の記録を集め、NMO等の走時差を整えて重ねる。単に未補正の任意波形を足す方法ではない。\par

\paragraph{再利用する考え方} 同じ補正という語でも観測量と誤差の原因を分ける。

\paragraph{共通理論} \hyperref[theory:exploration:seismicterms]{第\ref*{theory:exploration:seismicterms}節}, \hyperref[theory:exploration:continuation]{第\ref*{theory:exploration:continuation}節}, \hyperref[theory:gravity:corrections]{第\ref*{theory:gravity:corrections}節}

\par 参考：\cite{usgs_refraction2013,epa_seismic_reflection,usgs_static2004,usgs_processing2000,stanford_downward2001}。

\paragraph{関連問題} \hyperref[q:r01-reg-applied-34]{R01-REG-APPLIED-34}

\par 難易度 2（校準初稿）、優先度未評価。

\section{探査の重合とS/N平方根則}

\subsection{平成26年度 問21}
\label{q:h26-reg-applied-21}

\noindent\textbf{原典：} sources/H26/17\_applied\_science.pdf、PDF物理9頁。

\textbf{公式正答：} 3（single）。

\par\textbf{検討未確定：} 平方根則は振幅S/Nと無相関雑音・同位相信号の条件付きであり、群設置も方向・間隔に依存する。

\paragraph{問題の要旨} S/Nの平方根則・群設置・ADC・ノッチ・逆解析を区別する。

\Needspace{4\baselineskip}\paragraph{確認した条件・記述（独自要約）}\mbox{}\par

\noindent\textbf{原典照合：} 否定の対象、指数、周波数と周期、五つの候補の列順を原典画像で照合。\par

\noindent\textbf{前提：} 同じ信号に独立同分散の雑音が重なる場合の振幅S/Nを考える。ADCの標本化周波数は信号帯域に対して設定する。\par

\paragraph{問われる核心} S/Nの平方根則・群設置・ADC・ノッチ・逆解析を区別する。

\paragraph{必要知識} 同じ信号に独立同分散の雑音が重なる場合の振幅S/Nを考える。ADCの標本化周波数は信号帯域に対して設定する。

\paragraph{考え方・正答を選ぶ理由（査読待ち初稿）} 3が不適切である。標本化周波数を最高周波数と同じにすると、一般には異なる連続信号が同じ標本を与える。帯域制限を前提に二倍を超える周波数を用いる。他の候補は代表的な処理を説明する。

\Needspace{4\baselineskip}\paragraph{選択肢の検討（査読待ち初稿）}\mbox{}\par

\noindent\textbf{1：条件付きで適切} 重合のS/N改善が回数の平方根に比例する。同一位相信号はN倍、独立雑音の標準偏差は平方根N倍なので比は平方根N倍となる。\par

\noindent\textbf{2：条件付きで適切} 複数センサーの出力を群として結合し雑音を低減する。独立雑音の平均化に加え、センサーの間隔と入射方向に応じた空間フィルタ効果を利用する。対象の反射波を強めるには位相・配置が適切でなければならない。\par

\noindent\textbf{3：不適切} ADCの標本化周波数を信号の最高周波数と同じにする。一般の信号に必要なNyquist条件を満たさずaliasingを起こす。同じ周波数にするのは正しい一般的設定ではない。\par

\noindent\textbf{4：適切} ノッチで特定の周波数成分を遮断する。狭い阻止帯域で商用電源などの周期的な混入を抑える。指定帯域だけを通過させるフィルタではない。実装には有限の帯域幅と位相応答がある。\par

\noindent\textbf{5：適切} モデルの理論値と観測値の差を非線形最小二乗等で小さくする。順モデルF(m)の残差を目的関数にして地下物性を更新する代表的な逆解析である。線形問題や他のノルムもあり、非線形最小二乗だけが唯一の方法ではない。\par

\paragraph{再利用する考え方} 雑音を振幅で測るか電力で測るかを先に固定する。

\paragraph{共通理論} \hyperref[theory:exploration:stack]{第\ref*{theory:exploration:stack}節}, \hyperref[theory:exploration:sampling]{第\ref*{theory:exploration:sampling}節}, \hyperref[theory:exploration:inverse]{第\ref*{theory:exploration:inverse}節}

\par 参考：\cite{mit_signals2010,mit_moving2012,scipy_notch,scipy_fft,usgs_processing2000,pygimli_inversion}。

\paragraph{関連問題} \hyperref[q:h30-reg-applied-34]{H30-REG-APPLIED-34}

\par 難易度 3（校準初稿）、優先度未評価。

\subsection{平成27年度 問20}
\label{q:h27-reg-applied-20}

\noindent\textbf{原典：} sources/H27/17\_applied\_science.pdf、PDF物理7頁。

\textbf{公式正答：} 1（single）。

\par\textbf{検討未確定：} S/Nの振幅・電力定義が省略される。標本化の等号端点と完全復元、理想帯域通過の実装上の限界を区別する。

\paragraph{問題の要旨} 重合回数・コヒーレンス・標本化・帯域通過・FFTを比較する。

\Needspace{4\baselineskip}\paragraph{確認した条件・記述（独自要約）}\mbox{}\par

\noindent\textbf{原典照合：} 否定の対象、指数、周波数と周期、五つの候補の列順を原典画像で照合。\par

\noindent\textbf{前提：} 振幅S/N、周波数ごとの二信号の関係、理想フィルタを区別する。\par

\paragraph{問われる核心} 重合回数・コヒーレンス・標本化・帯域通過・FFTを比較する。

\paragraph{必要知識} 振幅S/N、周波数ごとの二信号の関係、理想フィルタを区別する。

\paragraph{考え方・正答を選ぶ理由（査読待ち初稿）} 1が不適切である。振幅S/Nは独立雑音のとき平方根N倍であり、N倍ではない。コヒーレンスは相互スペクトルを両信号の自己スペクトルで規格化して比較する。

\Needspace{4\baselineskip}\paragraph{選択肢の検討（査読待ち初稿）}\mbox{}\par

\noindent\textbf{1：不適切} スタックのS/Nが回数Nに比例する。信号はN倍、雑音の振幅尺度は平方根N倍なので振幅S/Nは平方根N倍である。電力比ならN倍となるため比の定義を明示する。\par

\noindent\textbf{2：適切} コヒーレンスは二つの時系列の周波数ごとの類似を示す相関係数である。周波数別の結び付きであり、代表的な二乗コヒーレンスは相互スペクトルの絶対値二乗を自己スペクトルの積で割る。時間領域の遅れ別相関そのものとは区別する。\par

\noindent\textbf{3：条件付きで適切} 最高周波数の二倍以上で標本化すれば完全復元できる。帯域制限・理想的標本と再構成を想定した説明である。最高周波数成分を含む場合の等号端点、量子化、有限長、雑音を無視できない。安全な条件はfs>2fmaxであり、境界では位相によって標本が全てゼロとなる例がある。\par

\noindent\textbf{4：条件付きで適切} 帯域通過は設定帯域を減衰させず、それ以外を遮断する。理想帯域通過の定義として読む。実際の設計では通過帯域のリップル、遷移帯域、有限減衰を持つため完全遮断を一般保証しない。\par

\noindent\textbf{5：適切} FFTは二のべき乗の標本数で効率的である。radix-2では段数がlog2Nとなる。二のべき乗でなければ計算不能という意味ではなく、混合基数やBluestein法もある。\par

\paragraph{再利用する考え方} 振幅比と電力比、理想定義と実装の差を明示する。

\paragraph{共通理論} \hyperref[theory:exploration:stack]{第\ref*{theory:exploration:stack}節}, \hyperref[theory:exploration:correlation]{第\ref*{theory:exploration:correlation}節}, \hyperref[theory:exploration:sampling]{第\ref*{theory:exploration:sampling}節}, \hyperref[theory:exploration:fft]{第\ref*{theory:exploration:fft}節}

\par 参考：\cite{mit_signals2010,mit_moving2012,scipy_notch,scipy_fft,usgs_processing2000,scipy_coherence}。

\par 難易度 3（校準初稿）、優先度未評価。

\subsection{平成30年度 問34}
\label{q:h30-reg-applied-34}

\noindent\textbf{原典：} sources/H30/17\_applied\_science.pdf、PDF物理11頁。

\textbf{公式正答：} 1（single）。

\par\textbf{検討未確定：} 群設置と平方根則は雑音相関・信号整列に依存し、標本化の等号と帯域制限を省略して完全復元を保証できない。

\paragraph{問題の要旨} 重合のS/N指数と電源雑音・群設置・標本化・逆解析を照合する。

\Needspace{4\baselineskip}\paragraph{確認した条件・記述（独自要約）}\mbox{}\par

\noindent\textbf{原典照合：} 否定の対象、指数、周波数と周期、五つの候補の列順を原典画像で照合。\par

\noindent\textbf{前提：} 振幅S/Nは同一位相信号と独立雑音を想定する。\par

\paragraph{問われる核心} 重合のS/N指数と電源雑音・群設置・標本化・逆解析を照合する。

\paragraph{必要知識} 振幅S/Nは同一位相信号と独立雑音を想定する。

\paragraph{考え方・正答を選ぶ理由（査読待ち初稿）} 1が不適切である。平方根Nの向上をNの二乗とする誤りである。電力S/NでもN倍でありNの二乗にはならない。周期雑音には狭帯域阻止、物性推定には順モデルの残差最小化を対応させる。

\Needspace{4\baselineskip}\paragraph{選択肢の検討（査読待ち初稿）}\mbox{}\par

\noindent\textbf{1：不適切} スタックのS/N改善が回数Nの二乗に比例する。信号の振幅Nに対して雑音の標準偏差は平方根Nである。振幅比も電力比もNの二乗とはならない。\par

\noindent\textbf{2：条件付きで適切} 複数センサーの出力を群として結合し雑音を低減する。独立雑音の平均化に加え、センサーの間隔と入射方向に応じた空間フィルタ効果を利用する。対象の反射波を強めるには位相・配置が適切でなければならない。\par

\noindent\textbf{3：適切} ノッチで特定の周波数成分を遮断する。狭い阻止帯域で商用電源などの周期的な混入を抑える。指定帯域だけを通過させるフィルタではない。実装には有限の帯域幅と位相応答がある。\par

\noindent\textbf{4：条件付きで適切} 最高周波数の二倍以上で標本化すれば完全復元できる。帯域制限・理想的標本と再構成を想定した説明である。最高周波数成分を含む場合の等号端点、量子化、有限長、雑音を無視できない。安全な条件はfs>2fmaxであり、境界では位相によって標本が全てゼロとなる例がある。\par

\noindent\textbf{5：適切} モデルの理論値と観測値の差を非線形最小二乗等で小さくする。順モデルF(m)の残差を目的関数にして地下物性を更新する代表的な逆解析である。線形問題や他のノルムもあり、非線形最小二乗だけが唯一の方法ではない。\par

\paragraph{再利用する考え方} 加算した信号と雑音を別々に計算して指数の誤りを検出する。

\paragraph{共通理論} \hyperref[theory:exploration:stack]{第\ref*{theory:exploration:stack}節}, \hyperref[theory:exploration:filters]{第\ref*{theory:exploration:filters}節}, \hyperref[theory:exploration:inverse]{第\ref*{theory:exploration:inverse}節}

\par 参考：\cite{mit_signals2010,mit_moving2012,scipy_notch,scipy_fft,usgs_processing2000,pygimli_inversion}。

\paragraph{関連問題} \hyperref[q:h26-reg-applied-21]{H26-REG-APPLIED-21}, \hyperref[q:r04-reg-applied-33]{R04-REG-APPLIED-33}

\par 難易度 3（校準初稿）、優先度未評価。

\subsection{令和4年度 問33}
\label{q:r04-reg-applied-33}

\noindent\textbf{原典：} sources/R04/17\_applied\_science.pdf、PDF物理11頁。

\textbf{公式正答：} 5（single）。

\par\textbf{検討未確定：} 完全復元の等号端点・帯域制限、群配置の方向と間隔を要する。表面波抑制を独立雑音の平方根則だけで説明しない。

\paragraph{問題の要旨} モデル更新・ノッチ・受振器群・標本化・垂直重合を区別する。

\Needspace{4\baselineskip}\paragraph{確認した条件・記述（独自要約）}\mbox{}\par

\noindent\textbf{原典照合：} 否定の対象、指数、周波数と周期、五つの候補の列順を原典画像で照合。\par

\noindent\textbf{前提：} 同一起振点の反復で対象反射の位相をそろえる。群配置は対象波と妨害波の空間応答を考える。\par

\paragraph{問われる核心} モデル更新・ノッチ・受振器群・標本化・垂直重合を区別する。

\paragraph{必要知識} 同一起振点の反復で対象反射の位相をそろえる。群配置は対象波と妨害波の空間応答を考える。

\paragraph{考え方・正答を選ぶ理由（査読待ち初稿）} 5が不適切である。垂直重合の振幅S/Nは理論上平方根N倍でありNの二乗ではない。表面波は反射波に対する妨害波となるが、空間的に相関した波なので独立ランダム雑音の平均化と同じ説明だけでは足りない。

\Needspace{4\baselineskip}\paragraph{選択肢の検討（査読待ち初稿）}\mbox{}\par

\noindent\textbf{1：適切} モデルの理論値と観測値の差を非線形最小二乗等で小さくする。順モデルF(m)の残差を目的関数にして地下物性を更新する代表的な逆解析である。線形問題や他のノルムもあり、非線形最小二乗だけが唯一の方法ではない。\par

\noindent\textbf{2：適切} ノッチで特定の周波数成分を遮断する。狭い阻止帯域で商用電源などの周期的な混入を抑える。指定帯域だけを通過させるフィルタではない。実装には有限の帯域幅と位相応答がある。\par

\noindent\textbf{3：条件付きで適切} 受振器群で表面波を抑え反射波を強めランダム雑音を減らす。複数受振器の結合には空間応答があり、反射と表面波の見掛け速度・波長に応じた選別を行える。配列が悪ければ必要信号も弱まるため無条件ではない。\par

\noindent\textbf{4：条件付きで適切} 最高周波数の二倍以上で標本化すれば完全復元できる。帯域制限・理想的標本と再構成を想定した説明である。最高周波数成分を含む場合の等号端点、量子化、有限長、雑音を無視できない。安全な条件はfs>2fmaxであり、境界では位相によって標本が全てゼロとなる例がある。\par

\noindent\textbf{5：不適切} 垂直スタックのS/Nが回数Nの二乗に比例する。独立同分散雑音の振幅S/N改善は平方根N倍である。Nの二乗は信号電力だけの増加を雑音の増加と比較せず取り出した誤りである。\par

\paragraph{再利用する考え方} 信号の種類、統計的相関、配列の空間応答を分ける。

\paragraph{共通理論} \hyperref[theory:exploration:stack]{第\ref*{theory:exploration:stack}節}, \hyperref[theory:exploration:filters]{第\ref*{theory:exploration:filters}節}, \hyperref[theory:exploration:inverse]{第\ref*{theory:exploration:inverse}節}

\par 参考：\cite{mit_signals2010,mit_moving2012,scipy_notch,scipy_fft,usgs_processing2000,pygimli_inversion}。

\paragraph{関連問題} \hyperref[q:h30-reg-applied-34]{H30-REG-APPLIED-34}

\par 難易度 3（校準初稿）、優先度未評価。

\section{表面波分散・S波感度と統合逆解析}

\subsection{令和7年度 問35}
\label{q:r07-reg-applied-35}

\noindent\textbf{原典：} sources/R07/17\_applied\_science.pdf、PDF物理13頁。

\textbf{公式正答：} 2（single）。

\par\textbf{検討未確定：} 均質半無限体は非分散の退化例であり1を全称にしない。P波・密度の感度はゼロではない。微動併用で長波長が得られることとモード同定の確認が必要。

\paragraph{問題の要旨} 表面波の分散、物性感度、反復逆解析、波長別深度、能動と微動の併用を比較する。

\Needspace{4\baselineskip}\paragraph{確認した条件・記述（独自要約）}\mbox{}\par

\noindent\textbf{原典照合：} 五候補の順序、最適切と不適切の方向、波種、深度目安、曲線名、最長と最短を原典画像で照合。\par

\noindent\textbf{前提：} 通常の不均質水平成層モデルとRayleigh分散の主なS波速度感度を考える。\par

\paragraph{問われる核心} 表面波の分散、物性感度、反復逆解析、波長別深度、能動と微動の併用を比較する。

\paragraph{必要知識} 通常の不均質水平成層モデルとRayleigh分散の主なS波速度感度を考える。

\paragraph{考え方・正答を選ぶ理由（査読待ち初稿）} 2が不適切である。Rayleigh分散は一般にS波速度へ強い感度を持ち、P波の影響がS波より大きいとする説明は逆である。波長を広い範囲で得ると浅部と深部を拘束しやすくなるが、観測に合うモデルが唯一とは限らない。

\Needspace{4\baselineskip}\paragraph{選択肢の検討（査読待ち初稿）}\mbox{}\par

\noindent\textbf{1：条件付きで適切} 水平成層のRayleigh位相速度は周波数による分散を持つ。深度で物性が変わる成層の通常モデルでは分散を持つ。均質等方半無限体や全層同物性の退化例は非分散なので、全てのモデルへの全称としない。\par

\noindent\textbf{2：不適切} Rayleigh位相速度へのP波の影響がS波より大きい。通常の地盤分散はS波速度が支配的に影響し、P波だけを主たる未知量とする説明は誤りである。\par

\noindent\textbf{3：適切} 理論と観測の位相速度差を減らすようモデル修正と計算を反復する。分散曲線を順計算し地下構造パラメータを更新する逆問題である。初期値・誤差と拘束にも依存する。\par

\noindent\textbf{4：適切} 探査深度は最大波長、浅部解像度は最小波長と強く関係する。長波長ほど深部に感度を持ち、短波長は浅部の小さな構造を拘束する傾向がある。一波長が一深度だけを測るわけではない。\par

\noindent\textbf{5：条件付きで適切} 能動表面波と微動を併用して探査深度を増やす。微動で長波長側を補える観測条件なら深部の拘束が増す。低周波が得られない場合やモードが不整合の場合は併用だけで深度増加を保証しない。\par

\paragraph{再利用する考え方} 分散の主感度、波長の感度範囲、データ併用の条件を示し唯一解を仮定しない。

\paragraph{共通理論} \hyperref[theory:exploration:surfacewaves]{第\ref*{theory:exploration:surfacewaves}節}, \hyperref[theory:exploration:arraylimits]{第\ref*{theory:exploration:arraylimits}節}, \hyperref[theory:exploration:inverse]{第\ref*{theory:exploration:inverse}節}

\par 参考：\cite{epa_surfacewaves,xia_rayleigh1999,japansea_spac2017,array_limits2002,pygimli_inversion}。

\paragraph{関連問題} \hyperref[q:r06-reg-applied-33]{R06-REG-APPLIED-33}

\par 難易度 3（校準初稿）、優先度未評価。

\section{探査波動の伝播原理}

\subsection{令和2年度 問35}
\label{q:r02-reg-applied-35}

\noindent\textbf{原典：} sources/R02/17\_applied\_science.pdf、PDF物理12頁。

\textbf{公式正答：} 2（single）。

\par\textbf{検討未確定：} 相反性と重ね合わせは媒体・線形性の条件付きである。等価層は異なる分野で使われる語でもあるため、ポテンシャル場の基本的用法と波動伝播の原理を区別する。

\paragraph{問題の要旨} 弾性波の伝播原理とポテンシャル場の等価層を区別する。

\Needspace{4\baselineskip}\paragraph{確認した条件・記述（独自要約）}\mbox{}\par

\noindent\textbf{原典照合：} 否定の対象、指数、周波数と周期、五つの候補の列順を原典画像で照合。\par

\noindent\textbf{前提：} 線形で相反的な媒体における波動の重ね合わせ・波面構成・伝播経路を対象とする。\par

\paragraph{問われる核心} 弾性波の伝播原理とポテンシャル場の等価層を区別する。

\paragraph{必要知識} 線形で相反的な媒体における波動の重ね合わせ・波面構成・伝播経路を対象とする。

\paragraph{考え方・正答を選ぶ理由（査読待ち初稿）} 2の等価層の原理が不適切である。USGSの等価層は重力の密度層や磁気の双極子層でポテンシャル場を表現するものとして確認できる。これは地震波伝播を記述する基本原理の四項目とは役割が異なる。

\Needspace{4\baselineskip}\paragraph{選択肢の検討（査読待ち初稿）}\mbox{}\par

\noindent\textbf{1：適切} フェルマの原理。波線経路に対する走時の停留条件を述べる。あらゆる経路が絶対的最短時間になるという一般化は避ける。\par

\noindent\textbf{2：不適切} 等価層の原理。重力・磁気の場を等価な密度又は双極子の層で表す考え方である。弾性波の波面や伝播経路を定める基本則とは異なる。\par

\noindent\textbf{3：適切} 相反原理。相反性が成立する線形媒体では震源と受振点を交換した応答が対応する。媒体や成分の条件を保つ必要がある。\par

\noindent\textbf{4：適切} ホイヘンスの原理。波面上の点からの二次波の構成で次の波面を考える。地震波の伝播・画像化の理解に関係する。\par

\noindent\textbf{5：条件付きで適切} 重ね合わせの原理。線形の波動方程式では解の線形結合も解となる。破壊を伴う大きな非線形変形へ無条件に適用しない。\par

\paragraph{再利用する考え方} 原理の名前を、場の表現・波面・経路・線形性・相反性という役割へ対応させる。

\paragraph{共通理論} \hyperref[theory:exploration:principles]{第\ref*{theory:exploration:principles}節}

\par 参考：\cite{usgs_equivalent2007,stanford_downward2001,mit_signals2010}。

\paragraph{関連問題} \hyperref[q:r07-reg-applied-32]{R07-REG-APPLIED-32}

\par 難易度 3（校準初稿）、優先度未評価。

\subsection{令和7年度 問32}
\label{q:r07-reg-applied-32}

\noindent\textbf{原典：} sources/R07/17\_applied\_science.pdf、PDF物理12頁。

\textbf{公式正答：} 3（single）。

\par\textbf{検討未確定：} 重ね合わせは線形媒体に限る。スネルの法則は波種と境界条件を指定して用いる。

\paragraph{問題の要旨} ホイヘンス・重ね合わせ・レンツ・スネル・フェルマを区別する。

\Needspace{4\baselineskip}\paragraph{確認した条件・記述（独自要約）}\mbox{}\par

\noindent\textbf{原典照合：} 否定の対象、指数、周波数と周期、五つの候補の列順を原典画像で照合。\par

\noindent\textbf{前提：} 地震探査で扱う弾性波の波面、線形性、屈折、波線の走時を対象とする。\par

\paragraph{問われる核心} ホイヘンス・重ね合わせ・レンツ・スネル・フェルマを区別する。

\paragraph{必要知識} 地震探査で扱う弾性波の波面、線形性、屈折、波線の走時を対象とする。

\paragraph{考え方・正答を選ぶ理由（査読待ち初稿）} 3のレンツの法則は電磁誘導が磁束変化を妨げる向きに生じることを表す。電磁探査には関係するが、他四項目のような弾性波伝播の基本則ではない。

\Needspace{4\baselineskip}\paragraph{選択肢の検討（査読待ち初稿）}\mbox{}\par

\noindent\textbf{1：適切} ホイヘンスの原理。二次波で波面を構成する考え方であり、波動伝播の理解に関係する。\par

\noindent\textbf{2：条件付きで適切} 重ね合わせの原理。線形波動では複数波の和を解として扱うことができる。大変形や非線形な破壊過程は同じ前提ではない。\par

\noindent\textbf{3：不適切} レンツの法則。電磁誘導の向きを磁束変化との関係から決める法則であり、地震波の反射・屈折の基本則とは異なる。\par

\noindent\textbf{4：適切} スネルの法則。境界の位相整合により角度と波速を結ぶ。地震波ではP波・S波へのモード変換を含む各波速で考える。\par

\noindent\textbf{5：適切} フェルマの原理。波線の走時が経路変分に対して停留する条件を用いる。最小という単語だけで複数経路の可能性を捨てない。\par

\paragraph{再利用する考え方} 電磁気の誘導と弾性波の伝播を、扱う物理量から分ける。

\paragraph{共通理論} \hyperref[theory:exploration:principles]{第\ref*{theory:exploration:principles}節}, \hyperref[theory:optics:refraction]{第\ref*{theory:optics:refraction}節}

\par 参考：\cite{mit_signals2010,stanford_downward2001,epa_seismic_reflection}。

\paragraph{関連問題} \hyperref[q:r02-reg-applied-35]{R02-REG-APPLIED-35}

\par 難易度 3（校準初稿）、優先度未評価。

